%% GEB LaTeX v2 — Part 04 / 7
%% 自动生成  split_v2.py
%% 独立编译（从 GEB_LaTeX_v2/ 目录）：
%%   xelatex -output-directory=split part04.tex

%% GEB 精排模板 ── XeLaTeX
\PassOptionsToPackage{unicode}{hyperref}
\PassOptionsToPackage{hyphens}{url}

%% ─────────────────────────────────────────────
%%  文档类：memoir（精排书籍首选）
%% ─────────────────────────────────────────────
\documentclass[
  12pt,
  twoside,
  openright,
  oldfontcommands,
]{memoir}

%% memoir 纸张尺寸（必须在字体等 usepackage 之前）
\setstocksize{260mm}{185mm}   % 大16开
\settrimmedsize{\stockheight}{\stockwidth}{*}
\settrims{0pt}{0pt}

%% ─────────────────────────────────────────────
%%  字体
%% ─────────────────────────────────────────────
\usepackage{fontspec}
\usepackage{xeCJK}

% ※（U+203B）Georgia 无此字符，强制用 CJK 字体渲染
\xeCJKDeclareCharClass{CJK}{"203B}

% geb.ttf：EPUB 自定义字体，含 Tai Viet 占位字形（U+AA9E–U+AAAF）
% 以及 U+31CF、U+39DF 等特殊字形
\newfontfamily\gebfont{geb.ttf}[Path=./fonts/]

% 中文正文：思源宋体（Noto Serif CJK SC）
\setCJKmainfont[
  BoldFont    = {Noto Serif CJK SC SemiBold},
  ItalicFont  = {Kaiti SC},          % 楷体作斜体
]{Noto Serif CJK SC}

% 中文无衬线（标题等）
\setCJKsansfont[
  BoldFont = {Heiti SC},
]{Heiti SC}

% 中文等宽
\setCJKmonofont{PingFang SC}

% 西文正文：Palatino（等高数字，基线与思源宋体对齐更佳）
\setmainfont[
  BoldFont       = {Palatino Bold},
  ItalicFont     = {Palatino Italic},
  BoldItalicFont = {Palatino Bold Italic},
]{Palatino}

% 西文无衬线
\setsansfont{Helvetica Neue}

% 数学字体
\usepackage{amsmath,amssymb}
\usepackage{unicode-math}

%% ─────────────────────────────────────────────
%%  版面：大16开（185×260mm）页边距（memoir 原生布局）
%% ─────────────────────────────────────────────
\setlrmarginsandblock{25mm}{20mm}{*}
\setulmarginsandblock{22mm}{25mm}{*}
\setheaderspaces{*}{7mm}{*}
\setheadfoot{\baselineskip}{12mm}
\checkandfixthelayout

%% ─────────────────────────────────────────────
%%  颜色
%% ─────────────────────────────────────────────
\usepackage{xcolor}
\definecolor{accent}{HTML}{1A3A5C}     % 深蓝（章号、页眉线）
\definecolor{accentlight}{HTML}{4A7DAA}% 中蓝（节标题）
\definecolor{rulegray}{HTML}{BBBBBB}   % 分隔线

%% ─────────────────────────────────────────────
%%  章节样式
%% ─────────────────────────────────────────────
\usepackage{titlesec}

% 章：大号宋体居中，上方装饰线，下方细实线
\titleformat{\chapter}[display]
  {\normalfont\centering}
  {\color{accent}\small\CJKfamily{}\sffamily 第\ \thechapter\ 章}
  {0pt}
  {\color{accent}\LARGE\bfseries\CJKfamily{}\rmfamily}
  [\vspace{4pt}\textcolor{rulegray}{\titlerule[0.4pt]}]

\titlespacing*{\chapter}{0pt}{30pt}{40pt}

% 节：宋体加粗左对齐，左侧蓝色竖线装饰
\titleformat{\section}
  {\normalfont\large\bfseries\color{accentlight}}
  {}
  {0pt}
  {\colorbox{white}{\raisebox{0pt}{\rule[-1pt]{3pt}{14pt}}\hspace{6pt}}}
  []

\titlespacing*{\section}{0pt}{18pt}{8pt}

% 小节
\titleformat{\subsection}
  {\normalfont\normalsize\bfseries}
  {}
  {0pt}
  {}

%% ─────────────────────────────────────────────
%%  页眉页脚（memoir 内建）
%% ─────────────────────────────────────────────
\makepagestyle{geb}
\makeevenhead{geb}{\small\thepage}{}{\small\itshape\leftmark}
\makeoddhead {geb}{\small\itshape\rightmark}{}{\small\thepage}
\makeheadrule {geb}{\textwidth}{0.4pt}
\pagestyle{geb}
\renewcommand{\chaptermark}[1]{\markboth{第\ \thechapter\ 章\enspace #1}{第\ \thechapter\ 章\enspace #1}}
\renewcommand{\sectionmark}[1]{\markright{\thesection\enspace #1}}

%% ─────────────────────────────────────────────
%%  中文标准名称
%% ─────────────────────────────────────────────
\renewcommand{\contentsname}{目录}
\renewcommand{\figurename}{图}
\renewcommand{\tablename}{表}
\renewcommand{\appendixname}{附录}
\renewcommand{\bibname}{参考文献}
\renewcommand{\indexname}{索引}

%% ─────────────────────────────────────────────
%%  段落间距与行距（memoir 内建命令）
%% ─────────────────────────────────────────────
\setSpacing{1.3}
\setlength{\parindent}{2em}
\setlength{\parskip}{3pt plus 1pt minus 1pt}
\raggedbottom

%% ─────────────────────────────────────────────
%%  图片与浮动体
%% ─────────────────────────────────────────────
\usepackage{graphicx}
\usepackage{float}
\usepackage{caption}
\captionsetup{
  font=small,
  labelfont={bf,color=accent},
  labelsep=quad,
  width=0.88\linewidth,
  justification=centering,
  singlelinecheck=false,
}

%% pandocbounded：memoir 兼容版，限制图片最大宽度和高度
\newcommand*\pandocbounded[1]{%
  \begingroup
  \setkeys{Gin}{width=\linewidth,height=0.8\textheight,keepaspectratio}%
  #1%
  \endgroup
}
\def\fps@figure{htbp}

%% ─────────────────────────────────────────────
%%  引用块样式
%% ─────────────────────────────────────────────
\usepackage{mdframed}
\mdfdefinestyle{gebquote}{
  leftline=true,
  rightline=false,
  topline=false,
  bottomline=false,
  linewidth=2pt,
  linecolor=accentlight,
  backgroundcolor=black!3,
  innerleftmargin=12pt,
  innerrightmargin=8pt,
  innertopmargin=6pt,
  innerbottommargin=6pt,
}
\renewenvironment{quote}
  {\begin{mdframed}[style=gebquote]\small}
  {\end{mdframed}}

%% ─────────────────────────────────────────────
%%  表格
%% ─────────────────────────────────────────────
\usepackage{longtable,booktabs,array,multirow,calc}
\usepackage{etoolbox}
\makeatletter
\patchcmd\longtable{\par}{\if@noskipsec\mbox{}\fi\par}{}{}
\makeatother
\IfFileExists{footnotehyper.sty}{\usepackage{footnotehyper}}{\usepackage{footnote}}
\makesavenoteenv{longtable}

% 美化 booktabs 规则
\setlength{\aboverulesep}{0pt}
\setlength{\belowrulesep}{0pt}
\setlength{\heavyrulewidth}{0.8pt}
\setlength{\lightrulewidth}{0.4pt}

%% ─────────────────────────────────────────────
%%  脚注
%% ─────────────────────────────────────────────
\usepackage[
  bottom,
  hang,
  flushmargin,
]{footmisc}
\setlength{\footnotemargin}{1.5em}
\renewcommand{\footnoterule}{%
  \vspace{4pt}\noindent\textcolor{rulegray}{\rule{0.3\linewidth}{0.4pt}}\vspace{4pt}%
}

%% ─────────────────────────────────────────────
%%  目录样式
%% ─────────────────────────────────────────────
\usepackage{titletoc}

%% memoir 的 chapternumberline/partnumberline 需映射到 \numberline
%% 才能被 titletoc 识别为「有编号」条目
\makeatletter
\renewcommand{\chapternumberline}[1]{\numberline{#1}}
\renewcommand{\partnumberline}[1]{\numberline{#1}}
\makeatother

\titlecontents{chapter}[0em]
  {\addvspace{6pt}\bfseries\color{accent}}
  {\makebox[3em][l]{\thecontentslabel}}
  {}
  {\hfill\contentspage}

\titlecontents{section}[3em]
  {\addvspace{2pt}\small}
  {}
  {}
  {\titlerule*[0.5em]{.}\contentspage}

%% ─────────────────────────────────────────────
%%  超链接
%% ─────────────────────────────────────────────
\usepackage{bookmark}
\IfFileExists{xurl.sty}{\usepackage{xurl}}{}
\urlstyle{same}
\usepackage{hyperref}
\hypersetup{
  pdftitle={哥德尔、艾舍尔、巴赫——集异璧之大成},
  pdfauthor={〔美〕侯世达},
  pdflang={zh-CN},
  colorlinks=true,
  linkcolor={accent},
  urlcolor={accentlight},
  citecolor={accent},
  pdfcreator={XeLaTeX + GEB Template},
}

%% ─────────────────────────────────────────────
%%  杂项
%% ─────────────────────────────────────────────
\usepackage{microtype}
\setlength{\emergencystretch}{3em}
\providecommand{\tightlist}{%
  \setlength{\itemsep}{0pt}\setlength{\parskip}{0pt}}

% 对话角色加粗
\newcommand{\speaker}[1]{\textbf{#1}}

%% ─────────────────────────────────────────────
%%  GEB v2 字体命令
%% ─────────────────────────────────────────────
\newCJKfontfamily\kaitifamily[BoldFont={Kaiti SC Bold}]{Kaiti SC}
\newcommand{\kaititext}[1]{{\kaitifamily #1}}
\newcommand{\kaitismall}[1]{{\kaitifamily\small #1}}
\newcommand{\kaitibold}[1]{{\kaitifamily\bfseries #1}}
\newcommand{\heititext}[1]{{\sffamily\bfseries #1}}
\newcommand{\heitismall}[1]{{\sffamily\bfseries\small #1}}
\newcommand{\dialogname}[1]{{\sffamily\bfseries #1}}
\newcommand{\westerntext}[1]{{\rmfamily #1}}
\newcommand{\westernital}[1]{{\rmfamily\itshape #1}}

%% ─────────────────────────────────────────────
%%  GEB v2 排版环境（memoir 内建 hangparas，不需要 hanging 宏包）
%% ─────────────────────────────────────────────

% 对话悬挂缩进行
\newenvironment{hangdialog}{%
  \begin{hangparas}{2em}{1}%
}{%
  \par\end{hangparas}%
}

% 楷体舞台指示（居中小楷）
\newenvironment{stagedir}{%
  \begin{center}\kaitifamily\small%
}{%
  \end{center}%
}

% 对话引导语（楷体斜体缩进）
\newenvironment{dialogguide}{%
  \begin{adjustwidth}{1em}{0em}\kaitifamily\itshape\small%
}{%
  \end{adjustwidth}%
}

% 仿宋/楷体引用块（带左竖线）
\newenvironment{fsquote}{%
  \begin{mdframed}[style=gebquote]\kaitifamily\small%
}{%
  \end{mdframed}%
}

% 楷体加粗引用框
\newenvironment{ktquote}{%
  \begin{mdframed}[style=gebquote]\kaitifamily\small%
}{%
  \end{mdframed}%
}

% 禅宗公案块
\newenvironment{zenkoan}{%
  \begin{mdframed}[style=gebquote]\kaitifamily\small%
}{%
  \end{mdframed}%
}

% 禅宗公案诗
\newenvironment{zenkoanpoem}{%
  \begin{center}\kaitifamily\small%
}{%
  \end{center}%
}

% 图片副说明（小楷居中，灰色）
\newenvironment{imgsub}{%
  \begin{center}\kaitifamily\small\color{rulegray}%
}{%
  \end{center}%
}

% 对话诗歌
\newenvironment{dialogpoem}{%
  \begin{adjustwidth}{3em}{0em}\kaitifamily%
}{%
  \end{adjustwidth}%
}

% 章节前置数字大标题（如"第一章"）
\newenvironment{chappreheader}{%
  \begin{center}\large\sffamily\bfseries\color{accent}%
}{%
  \end{center}\vspace{-10pt}%
}

% 部分前置大标题
\newenvironment{partpreheader}{%
  \begin{center}\Large\sffamily\bfseries\color{accent}%
}{%
  \end{center}\vspace{-10pt}%
}

%% ─────────────────────────────────────────────
%%  扉页（自定义，替换 maketitle）
%% ─────────────────────────────────────────────
\newcommand{\gebmaketitle}{%
  \begin{titlingpage}
    \vspace*{6\baselineskip}
    \begin{center}
      {\fontsize{28}{34}\selectfont\bfseries\color{accent}
        哥德尔、艾舍尔、巴赫}\\[12pt]
      {\Large\color{rulegray}——集异璧之大成}\\[40pt]
      \textcolor{rulegray}{\rule{0.4\linewidth}{0.4pt}}\\[20pt]
      {\large \textit{Gödel, Escher, Bach: An Eternal Golden Braid}}\\[6pt]
      {\normalsize \color{accentlight} Douglas R.\ Hofstadter}\\[40pt]
      {\normalsize 〔美〕侯世达　著\\[4pt]
        本书翻译组　译}\\[60pt]
      {\small\color{rulegray} 商务印书馆}
    \end{center}
    \vfill
  \end{titlingpage}
}

%% pandoc 在 body 中会注入 \babel@toc，提供空定义避免报错
\makeatletter
\providecommand{\babel@toc}[2]{}
\makeatother

%% memoir + XeLaTeX 下需手动告知 xdvipdfmx 纸张尺寸
\AtBeginDocument{%
  \special{papersize=\the\paperwidth,\the\paperheight}%
}

%% ═══════════════════════════════════════════
%%  正文开始
%% ═══════════════════════════════════════════

\begin{document}
\mainmatter
\setcounter{chapter}{24}
\markboth{}{}

\chapter{无门与哥德尔}

什么是禅宗？

我不敢肯定我知道禅宗是什么。一方面，我觉得我对禅宗非常明了；但另一方面，我又觉得我永远也没法明白它。自从我大学一年级的语文老师在课堂上给我们大声朗读赵州的``无''，我就开始与禅宗式的生活观进行搏斗，也许我永远也不会停止这么做。对我来说，禅宗是智力流沙------晦涩、无意义、紊乱、无法无天。它撩人而又令人恼火。不过它也很幽默，让人耳目一新，富于吸引力。禅宗是有其特殊的意义、光芒和明晰性的。我希望在这一章里能把我的这些感受传达一些给读者。这样（虽然这似乎挺奇怪），我们就被直接引向哥德尔的理论了。

佛教禅宗的基本教条之一是：没有任何办法能刻划禅宗是什么。无论你用什么样的词语空间努力去涵盖禅宗，都不会成功，它总要再冒出去。看起来，阐释禅宗的所有努力似乎完全都是浪费时间。但禅宗信徒们并不这么看。比如说，禅宗的公案，虽然是用词语表达的，乃是禅宗探究的中心部分。公案是当作``触发器''的，它们自己并不含有足够的信息以得到顿悟，但它们可能足以解开人们心智中导致顿悟的机制。不过一般说来，禅宗的观点认为词语与真理是不相容的，或至少是词语不能捕捉到真理。

无门禅师

也许就是为了以一种极端的方式指出这一点，无门禅师在十三世纪时编辑了四十八个公案，并于每个公案后面附了评注及一首小``颂''。这部作品被称为《无门之门》，或《无门关》。有趣的是，无门与斐波那契的生活年代几乎完全相同：无门从1183年至1260年生活在中国，斐波那契于1180年到1250年生活在意大利。对于那些希望通过《无门关》而看懂或``理解''公案的人，看了《无门关》之后大概会惊诧不已：那些评注及诗句本应是澄清公案的意思的，实际上却和公案本身一样晦涩。看这个例子：\footnote{}

公案：

清凉大法眼，因僧斋前上参，眼以手指帘，时有二僧同去卷帘。

眼曰：``一得，一失。''

无门曰：

且道：是谁得谁失？若向者里着得一只眼，便知清凉国师败阙处。然虽如此，却忌向得失里商量！

颂曰：

卷起明明彻太空，太空犹未合吾宗。

争似从空都放下？绵绵密密不通风！

现在我们再来看一个：\footnote{}

公案：

五祖曰：``譬如水牯牛过窗棂，头，角，四蹄都过了，因什么尾巴过不得？''

无门曰：

若向者里颠倒着一只眼，下得一转语，可以上报四恩，下资三有。其或未然，更须照顾尾巴始得！

颂曰：

过去堕坑堑，回来却被坏。

者些尾巴子，直是甚奇怪！

我想你得承认，无门并没有真的把每件事都澄清。可以说，元语言（无门用的语言）与对象语言（公案的语言）没有什么太大差别。在有些人看来，无门的那些评注纯属故意装傻，也许就是为了表明花时间谈论禅宗是毫无用处的。

\pandocbounded{\includegraphics[keepaspectratio]{./media/OEBPS/Images/img_046.jpg}}

图46．三界

埃舍尔作（版画，1955）

不过，我们可以从不止一个层次来理解无门的评注。比如，考虑下面这个例子：\footnote{}

公案：

南泉和尚问云：``还有不与人说的法么？''

泉云：``有。''

僧云：``如何是不与人说的法？''

泉云：``不是心，不是佛，不是物！''

无门曰：

南泉被者一问，直得揣尽家私，郎当不少！

颂曰：

叮咛损君德，无言真有功。

任从沧海变，终不为君通！

这首诗里，无门似乎说了些禅宗里非常核心的东西，而非一些傻话。然而奇怪的是，这首诗是自指的，因此它不仅评论了南泉所说的话，也阐明了自身的无效性。这种悖论是禅宗的一大特点。它是``破坏逻辑头脑''的一种尝试。从公案中同样也能看到这种悖论特质。看了无门的评注之后，你觉得南泉真的那么肯定他给的回答吗？或者他的回答的正确性根本就是无关紧要的？或者正确性这种东西在禅宗里就没有位置？正确性与真理之间的区别在哪里？或许这种区别是不存在的？如果南泉说``没有不与人说的法''，这又如何呢？结果会有什么不同吗？他的话还会记录在公案里而流传下来吗？

\pandocbounded{\includegraphics[keepaspectratio]{./media/OEBPS/Images/img_047.jpg}}

图47．露珠

艾舍尔作（镂刻凹版，1947）

这里是另一个旨在破坏逻辑头脑的公案：\footnote{}

道悟趋禅师问曰：``吾欲求真理，吾应修至何等心境方可求之？''

师曰：``本来无心，故无心境可修；本来无真理，故无由求之。''

``既无心可修，无真可求，聚这些和尚此处习禅修行何故？''

师曰：``此处并无寸地，这些和尚何以聚得？吾口中无舌，何以集而教之？''

道悟问：``师何以言谎？''

师曰：``吾既无舌语人，何能言谎？''

道悟惘然曰：``吾不了师言。''

师曰：``吾亦不自了。''

如果有什么公案使人困惑，这就是一个。而且很可能它的目的恰恰就在于引起困惑，因为人的心智处于困惑状态时就会在某种程度上不合逻辑地运转。只有跨出逻辑，摆脱理论，人才能跃入顿悟境地。可是逻辑到底怎么不好了？为什么它会阻止人们达到顿悟？

\pandocbounded{\includegraphics[keepaspectratio]{./media/OEBPS/Images/img_048.jpg}}

图48．另一个世界

艾舍尔作（木雕版，1947）

禅宗反对二元论的斗争

回答这些问题需要对顿悟是什么有一定的了解。所谓顿悟，最简明扼要地说或许就是：超越二元论。那么二元论又是什么呢？二元论就是把世界从概念上分划为种种范畴。这么一种非常自然的倾向能否被超越？我在``分划''前面加了一个修饰语``从概念上''，可能会让人觉得这是种智力上的或意识中的努力，因而也许会给人一种印象：克服二元论只需简单地抑制思维即可（好像抑制思维实际上很简单似的！）。但把世界裂成各种范畴也发生在远远低于思维这种较高层次的低层次上。事实上，二元论不仅是概念上对世界的划分，同样也是感知觉上对世界的划分。换句话说就是，人类的感知觉本质上是种二元现象------这使得追求顿悟至少是种逆水行舟般的奋斗了。

在禅宗看来，二元论的核心就是词语------普通的词。对词的使用必然导致二元化，因为每个词很明显地就是代表了一个概念范畴。所以，禅宗的一个主要部分就是为反对依靠词语而斗争。反对使用词语的最有力的工具之一就是公案，其中词被如此彻底地误用，以至那些认真看待公案的人会晕头转向，理不清自己的神智。因此，说顿悟的敌人是逻辑也许是不对的，它应该是二元化，借助词语的思维。事实上还要更基本：是知觉。一旦你感知到一个客体，你就把它与世界的其余部分划分开了；你人为地把世界分成部分，你于是就远离了``道''。

下面这个公案展示了同词语的斗争：\footnote{}

公案：

首山和尚拈竹篦示众云：``汝等汝人，若唤作竹篦则触；不唤作竹篦则背。汝诸人且道：唤作什么？''

无门曰：

唤作竹篦则触；不唤作竹篦则背。不得有语，不得无语。速道，速道！

颂曰：

拈起竹篦，行杀活令。

背触交驰，佛祖乞命！

（这里的``佛祖''指的是禅宗的六位倍受尊崇的创建者，其中菩提达摩是初祖，慧能是六祖。）

\pandocbounded{\includegraphics[keepaspectratio]{./media/OEBPS/Images/img_049.jpg}}

图49．白天与黑夜

艾舍尔作（木刻，1938）

为什么称之为竹篦就违反了它的实在性？可能是因为这种划类显得像是捕获了实在，而实际上这样的陈述却根本连表面都未能触及。可以比较一下``5是个素数''这个说法。那么多东西------无数多的事实依据------都省略掉了。另一方面，不称之为竹篦，那也的确是在无视一个特定的事实，不管这个事实是多么微不足道。因此，词语把我们引向某些真理------或许，同样也引向某些虚假------但肯定不能引向所有真理。你若是依赖词语走向真理，那就如同依赖一个不完全的形式系统而走向真理。一个形式系统的确会给你一些真理，但正如我们马上就会看到的，无论一个形式系统多么强有力，都不可能给出所有真理。数学家们的困窘就在于：除了形式系统，还有什么可以依靠？而禅宗信徒的困窘则是：除了词语，还有什么可以依靠？无门把这一困境阐述得很清楚：``不得有语，不得无语。''

下面又是南泉了：\footnote{}

公案：

南泉因赵州问：``如何是道？''

泉云：``平常心是道。''

州云：``还可趣向否？''

泉云：``拟向即乖！''

州云：``不拟怎知是道？''

泉云：``道不属知，不属不知；知是妄觉，不知是无记。若真达不疑之道，犹如太虚廓然洞豁，岂可说是非耶！''{{[}见图50{]}}

\pandocbounded{\includegraphics[keepaspectratio]{./media/OEBPS/Images/img_050.jpg}}

图50．果皮

艾舍尔作（木雕，1955）

这个奇特的陈述似乎充满了悖论。这多少让人想起那个治嗝偏方：``绕着大树跑三圈，脑子里始终不想`乌鸦'这个词。''禅宗哲学似乎体现了这样一个观念：通向终极真理之路，就像那个唯一的治嗝偏方一样，会是充斥悖论的。

主义、无方式以及云门

如果词语不好，思维也不好，那么什么是好的？当然，这么问已经是彻头彻尾的二元论了，但我们是在讨论禅宗，无需假装出笃信禅宗的样子------因此我们可以试着来认真地回答这个问题。对于禅宗所追求的东西，我这里有个名称：``主义''。主义是反哲学的，是一神摒弃思维的存在方式。主义的大师是石头，树，蛤蟆。高等动物若想达到主义，就得经过一番奋斗，而且是永远不可能完全达到的。不过，人们偶尔还是会有幸瞥一眼主义的。也许下面这个公案就赐予了这么一瞥：\footnote{}

公案：

沩山和尚始在百丈会中充典座，百丈将选大沩主人，乃请同首座对众下语，出格者可往。百丈遂拈净瓶置地上，设问云：``不得唤作净瓶，汝唤作什么？''

首座乃云：``不可唤作也。''

百丈却问于山，山乃倒净瓶而去。

百丈笑云：``第一座输却山子也！''因命沩为开山。

摒弃感知，摒弃逻辑、词语、二元化的思维------这就是禅宗的实质，主义的实质。这即是``无''方式------非智能，非机械，就是``无''。赵州处于无方式中，而这就是为什么他的``无''废问了那个问题。对于云门禅师，无方式是再自然不过的了：\footnote{}

{（云门禅师）}上堂，拈拄杖曰：``拄杖子化为龙，吞却乾坤了也。山河大地，甚处得来？''

禅宗采纳整体论，并且推向逻辑上的极端。如果整体论是断言事物必须作为一个整体被理解，而非其各个部分的总和，那么禅宗走得更远，认为整个世界根本就不能被划分为一个个事物。划分世界就会误入岐途，因而就不能达到顿悟了。

公案：\footnote{}

一师因一僧问曰：``如何是道？''

师曰：``正眼前是道。''

``如何我不自见？''

``汝自虑故。''

``师知之否？''

师曰：``汝但见二分：言`我不自见'，、`师见之'，汝目障矣。''

``无我无你，可得见否？''

``无我无你，谁欲见之？''

很明显，禅师想传达这样一种观念，即顿悟状态意味着自我和宇宙之间的分界消解了。这将是二元论的真正终结，因为正如他所说的，任何一个有感知愿望的系统都将不复存在。但除了死亡，那还能是种什么状态？一个活生生的人怎么能消解他自己与外部世界之间的分界线呢？

禅宗与堕界

禅宗和尚拔对给一个行将圆寂的门人写了一封信，信中说：``汝之无终之终，一似雪片溶于清气之中。''雪片，曾是宇宙中完全可见的一个子系统，现在溶解于它曾依托的那个更大的系统中了。虽然它已不再是以一个清晰可见的子系统而存在，其实质却依然是存在的，并将这么保持下去。它漂浮在堕界中，与没打出的嗝在一起，与没人读的故事中的人物在一起\ldots\ldots 这是我对那番话的理解。

禅宗认识到了自身的局限，正如数学家们逐渐也认识到了公理化方法作为获得真理的方法其局限所在。这并不意味着禅宗对自身之外有什么东西有个确切的答案。数学家们不清楚在形式化推理之外还有什么有效的推理形式，禅宗也强不了多少。关于禅宗的界限，下面这个奇怪的公案给了一个清楚的禅宗说法，看上去很能代表南泉的精神：\footnote{}

师{（洞山禅师------译注）}谓众曰：``知有佛向上人，方有语话分。''

僧问：``如何是佛向上人？''

师曰：``非佛。''

总是能走得更远。顿悟并非禅宗的终点。而且并没有一付超越禅宗的良方。唯一坚实可靠的是，佛非道也。禅宗是一个系统，不可能成为它自己的元系统。总是有东西处在禅宗之外，那是无法在禅宗之内完全了解或说清楚的。

艾舍尔与禅宗

在质疑感知、设置无答案的荒唐哑谜方面，禅宗有个伴儿：画家艾舍尔。比如《白天与黑夜》（图49）这幅``背触交驰''（借用无门的话）的杰作。人们会问，``那些真的是鸟吗？那些真的是田野吗？那真的是白天吗？真的是黑夜吗？''然而我们又都知道这么问毫无意义。这幅画同禅宗的公案一样，是在努力打破逻辑头脑。艾舍尔还特别喜欢作一些矛盾的画，像《另一个世界》（图48）------把实在与非实在摆弄来摆弄去，与禅宗摆弄实在与非实在的方式完全一样。是否应该认真对待艾舍尔？是否应该认真对待禅宗？

在《露珠》（图47）中，对于反射有种精致的、俳句般的探究。另外还有两幅月亮映在平和的水面上的宁静图景：《泥塘》（图51）和《水面涟漪》（图52）。水中的月亮是许多公案的主题。下面是一个例子：\footnote{}

千代能尼师在园觉佛光大师会下学禅，久久不能获得参悟的结果。一个月明之夜，千代以一个旧桶提水，偶因桶箍破裂而桶底脱落，豁然获得大自在，作一偈以记其事：

扶持旧桶，桶底忽脱。

桶里无水，水中无月。

《三界》是一幅艾舍尔的画（图46），也是一个禅宗公案的主题：\footnote{}

{（瑞岩）}问：``三界竞起时如何？''

师{（岩头禅师）}曰：``坐却著。''

曰：``未审师意如何？''

师曰：``移取庐山来，即向汝道。''

\pandocbounded{\includegraphics[keepaspectratio]{./media/OEBPS/Images/img_051.jpg}}

图51．泥塘

艾舍尔作（木刻，1952）

``三比二''与艾舍尔

在《辞》（图149）里，相左的事物在不同的层次上被纳入了一个统一体。顺着这幅画我们看到逐渐的变迁：从黑鸟到白鸟到黑鱼到白鱼到黑蛙到白蛙到黑鸟\ldots\ldots 六步之后，回到了我们开始的地方！这是调和黑与白这两分，还是鸟、鱼、蛙这三分？或者是2的偶数性与3的奇数性之间的对立导致的六位一体？在音乐里，六个等长音符会造成节奏岐义------是两个一组的三组，还是三个一组的两组？这种岐义有个名称：``三比二''。肖邦是``三比二''的大师：见他的华尔兹（作品42号）、练习曲（作品25号之二）。在巴赫的作品中，出现了这种现象的有第五键盘组曲中的小步舞曲，以及那首不可思议的G小调第一无伴奏小提琴奏鸣曲的终曲。

\pandocbounded{\includegraphics[keepaspectratio]{./media/OEBPS/Images/img_052.jpg}}

图52．水面涟漪

艾舍尔作（油毡浮雕，1950）

逐步进入到《辞》的中央时，差别便渐渐开始模糊。于是到最后剩下的不是三个，也不是两个，而是一个单质：``辞''（画中是拉丁文VERBUM，就是``辞''的意思）在熠熠闪光------也许这就是顿悟的象征。有趣的是，``辞''不仅是个词，而且它的意义也即``词''------并不完全与禅宗的观念相符。另一方面，``辞''是画中唯一的一个词。洞山禅师曾说过，``全部佛经三藏可以由一个字来表达。''我一直在想，该是什么样的解码机制才能从一个字里抽出三藏经文？也许是个带有两个半球的东西。

因陀罗之网

\pandocbounded{\includegraphics[keepaspectratio]{./media/OEBPS/Images/img_053.jpg}}

图53．三个球Ⅱ

艾舍尔作（版画，1946）

最后，我们来看看《三个球之二》（图53），其中世界的每个部分都包含了、同时也被包含于其它部分：写字台反映出在它上面的球，球与球之间彼此反映，同时也反映出写字台，并反映出这幅画本身以及正在作画的艺术家。所有的事物彼此之间都有着无尽的联结。这里还只不过是给个提示，然而提示已经足够了。佛教中的因陀罗之网就象征着遍布宇宙的一张无穷无尽的线网，水平的线穿过空间，垂直的线穿过时间。每个线与线的交汇处都是一个个体，每个个体都是颗明珠。``神''的巨大光芒照耀并穿透每颗明珠，而且，每颗明珠不但在反映网上其它明球的光泽------也反映着遍布于宇宙的每一反映的每一反映。

这使我脑子里产生了一幅重正化粒子的图像：在每个电子中，都有虚光子、正电子、中微子、π介子\ldots\ldots；每个光子中又有虚电子、质子、中子、π介子\ldots\ldots；而在每个π介子中，有\ldots\ldots{}

但随即又来了另一幅图像：关于人。每个人都反映在许多其他人的脑子里，那些人又反映在另外一些人的脑子里，如此下去。

这两幅图像都能由扩充迁移网简明漂亮地表示出来。在粒子的情形，将是每个粒子范畴都有一个网；在人的情形，是每个人有个网。每个网都含有对其它网的调用，这样就在每个ATN周围包上了一个ATN的虚云。调用某一个网导致对别的网的调用，这个过程在终了之前会像多级瀑布般地流泻到任意远。

无门论无

让我们回到无门，把这番闲扯收归于禅宗。下面是他对赵州的``无''的评论：\footnote{}

参禅须透祖师关，妙语要穷心路决。祖关不透，尽是依草附木精灵。且道：如何是祖师关？只者一个``无''字，乃宗门一关也。遂目之曰：``禅宗无门关。''透得过者，非但亲见赵州，便可与历代祖师把手共行，眉毛厮结，同一眼见，同一耳闻，岂不庆快？莫有要透关底么？将三百六十骨节、八万四千毫窍，通身起个疑团，参个``无''字，昼夜提撕。莫作虚无会，莫作有无会，如吞了个热铁丸相似，吐又吐不出，荡尽从前恶知恶觉；久久纯熟，自然内外打成一片，如哑子得梦只许自知；蓦然打发，惊天动地，如夺得关将军大刀入手，逢佛杀佛，逢祖杀祖，于生死岸得大自在，向六道四生中游戏三昧。且作么生提撕？尽平生气力举个``无''字。若不间断，好似法烛，一点便着。

从无门到WU谜题

我们现在从赵州的``无''那无比漂渺的高度上降下来，降到单调又平凡的侯世达的WU\ldots\ldots 我知道你已为这个WU花费了不少精力（读第一章的时候），所以我现在很愿意摆出答案了。那时提出的问题是：

``WU还有定理性也无？''

答案并非一个含糊其词的``无''，而是一声响亮的``不''。为了指明这一点，我将借用二元论的、逻辑的思维。

在第一章中有两点观察是至关重要的：

⑴．WU谜题之所以有难度，很大程度上是因为它涉及了加长规则与缩短规则两者的交织；

⑵．解开这个谜题的希望还是有的，只是要借助一个工具：数论。某种意义上讲，数论在处理这种复杂度的问题时具有合适的深度。

第一章里我们没有非常仔细地从这种角度来分析WU谜题，现在我们就要这么做了。我们会见到那第二点观察（当推广一下，超出意义不大的WJU系统的时候）是数学界最富成果的发现之一，它改变了数学家们对自己那门学问的看法。

为使读者参考起来方便，我重述WJU系统如下：

符号：W，J，U

公理：WJ

规则：

Ⅰ．若xJ是定理，则xJU也是定理。

Ⅱ．若Wx是定理，则Wxx也是定理。

Ⅲ．在任何定理中，JJJ可以被换成U。

Ⅳ．UU可以从任何定理中删去。

无门告诉我们如何解开WU谜题

于是，按照上述两点观察，WU谜题无非是一道披着印符外衣的自然数问题。假若我们能想出一个办法，把它转换到数的领域里，我们大概就能解它了。让我们好好想想无门的话：``若向者里着得一只眼，便知清凉国师败阙处。''但``着得一只眼''又怎样呢？

你要是数一数定理中J的数目，很快你就会发现好像永远不会是0。换句话说，似乎无论折腾多少回加长与缩短的操作，我们永远也得不到一个清除了所有J的符号串。让我们把串中J的数目称作该串的``眼''数，因为J的存在就好像是在我们想要的WU串中间钻了一些``洞眼''。注意公理WJ中有一只``眼''。我们不但可以说明定理的``眼''数不可能是0，还可以说明，``眼''数永远不会是3的倍数。

作为开始，先注意规则Ⅰ与规则Ⅳ对``眼''数丝毫不发生影响。所以我们只需要考虑规则Ⅱ与Ⅲ。规则Ⅲ只不过是使``眼''数恰好减3，使用这个规则之后，得到的输出或许也可能是3的倍数------但仅只在输入已经是3的倍数这一情况下。简而言之，规则Ⅲ决不会无中生有地产生3的倍数。它必须在``眼''数已经是3的倍数的情况下才能产生``眼''数为3的倍数的串。同样的事情也发生在规则Ⅱ身上，因为它使``眼''数翻倍。理由是，若2n被3整除，那么------由于3不整除2------n就得被3整除（数论里的一个简单事实）。无论是规则Ⅱ还是规则Ⅲ，都不能无中生有地产生3的倍数。

而这恰恰就是解开WU谜题的关键所在！下面是我们已经知道了的事情：

⑴．最一开始的``眼''数是1（不是3的倍数）；

⑵．有两条规则丝毫不影响``眼''数；

⑶．另两条规则的确要影响``眼''数的值，但除非初始给出的是3的倍数，否则它俩也决不会产生3的倍数。

结论是------也是典型的继承性论证方式------``眼''数永远无法变成3的倍数。特别地，0就决不可能是``眼''数的值了。因此，WU不是WJU系统的定理。

请注意，即使是作为一个关于``眼''数的谜题，这个问题也会因为加长规则与缩短规则的交织而使人迷惑。零成了追求目标，

``眼''数可以增大（规则Ⅱ），可以减小（规则Ⅲ）。我们大概会想，不断地换着使用那些规则，最终或许能达到0，于是就这么干了下去，直到我们站出来分析一番我们面临的到底是什么。现在，多亏一个简单的数论推理过程，我们知道了那是不可能的。

对WJU系统进行哥德尔配数

并非所有以WU谜题为代表的那类问题都像WU谜题那么容易解。但我们已经看到了至少有这么一个这种问题是可以嵌入数论并用数论来解决的。我们现在就要看到，有那么一种方法，可以把所有的关于任何形式系统的问题都嵌入数论。这之所以能做到是因为哥德尔所发现的一种特殊的同构。下面我用WJU系统来说明这种方法。

我们从WJU系统的记号开始。先把每个符号都映射到一个新的符号上去：

\begin{longtable}[]{@{}lll@{}}
\toprule\noalign{}
\endhead
\bottomrule\noalign{}
\endlastfoot
W & 　\gebfont{⇔}　 & 3 \\
J & 　\gebfont{⇔}　 & 1 \\
U & 　\gebfont{⇔}　 & 0 \\
\end{longtable}

这个对应是任意选定的。若要一定说出个理由，可以说每个符号都多少有点像它所映上的那个新符号。每个数都称作相应字母的哥德尔数。现在我敢肯定你已经猜到了字母串的哥德尔数将是什么了：

\begin{longtable}[]{@{}lll@{}}
\toprule\noalign{}
\endhead
\bottomrule\noalign{}
\endlastfoot
WU & 　\gebfont{⇔}　 & 30 \\
WJJU & 　\gebfont{⇔}　 & 3110 \\
\multicolumn{3}{@{}l@{}}{%
等等} \\
\end{longtable}

这很简单。很清楚，这个记号间的映射是个保持信息的变换，就像是在两个不同的乐器上演奏同一旋律。

我们现在来看看WJU系统中的一个典型的推导，我们同时用两种记号写：

\begin{longtable}[]{@{}llllll@{}}
\toprule\noalign{}
\endhead
\bottomrule\noalign{}
\endlastfoot
⑴ & WJ & ------ & 公理 & ------ & 31 \\
⑵ & WJJ & ------ & 规则2 & ------ & 311 \\
⑶ & WJJJJ & ------ & 规则2 & ------ & 31111 \\
⑷ & WUJ & ------ & 规则3 & ------ & 301 \\
⑸ & WUJU & ------ & 规则1 & ------ & 3010 \\
⑹ & WUJUUJU & ------ & 规则2 & ------ & 3010010 \\
⑺ & WUJJU & ------ & 规则4 & ------ & 30110 \\
\end{longtable}

左边一列是使用我们已熟悉的四条印符规则而得到的。右边一列同样也可以当作是使用了一组类似的印符规则而生成的。不过右边的一列有种二重性。下面我来解释。

从印符和算术两个角度看问题

我们可以说第五个串（``3010''）是从第四个来的，方法是在右边放上一个``0''；另一方面，我们同样也可以把这一转换看成是由于算术运算而导致的------确切地说就是乘以10。当自然数按十进制写出时，乘以10和在右边放上一个0是没有差别的。我们可以利用这一点写出一条算术规则，与印符规则I相对应：

算术规则Ⅰa：一个数的十进制展开的右侧如果以1结尾，则可乘以10。

我们可以不必指称十进制展开中的符号，而用算术方式来描述最右边的数字：

算术规则Ⅰb：一个数被10除时的余数如果是1，则可乘以10。

我们可以仍是固守纯粹的印符规则格式，比如下面这个：

印符规则Ⅰ：从任何一个最右侧符号为``1''的定理都可以得到一个新的定理，方法是在那个``1''右侧拼上一个``0''。

它们的效果是一样的。这就是为什么右边一列具有``二重性''：可以视为一系列的印符操作，改变着符号排列模式；也可以视为一系列算术运算，改变着数的量级。但有充足有力的理由使我们对算术的版本更感兴趣。跨出一个纯印符系统，再走入另一个同构的印符系统，这没有多大意思；而离开印符领域然后步入一个与之同构的数论的某个部分则会激发出未开发的潜力。这就像一个人前半生一直通晓乐谱，但只是用眼睛看而已------然后，忽然有一天，有人告诉了他声音与乐谱之间的映射。那将是怎样一个崭新而又丰富的世界！在这里，就好像一个人前半生一直都熟知符号串的样子，但只是那些串的形状而已，从不知意义何在------然后，忽然有一天，有人告诉了他事物与串之间的映射。这是怎样的一种启迪！哥德尔配数的发现被人们比作笛卡尔关于平面曲线与二元方程之间同构的那个发现：一旦你见到了，真是令人难以置信地简单------而一个崭新宽广的世界就这么展开了。

不过，在我们跳到结论之前，你或许会愿意看看这个同构的较高层次的全貌是什么样。这也是个很好的练习。其思想在于使给出的各个算术规则与WJU系统的每个印符规则在效果上没有差别。

下面是实现的方案之一。规则中的m和k是任意的自然数，n是小于10\textsuperscript{m}的任何自然数。

规则1：若有了10m+1，则还可以有10×(10m+1)。

例子：从第4行到第5行。这时，m=30。

规则2：若有了3×10\textsuperscript{m}+n，则还可以有10\textsuperscript{m}×(3×10\textsuperscript{m}+n)+n。

例子：从第1行到第2行，这时m与n都是1。

规则3：若有了k×10\textsuperscript{m+3}+111×10\textsuperscript{m}+n，则还可以有k×10\textsuperscript{m+1}+n。

例子：从第3行到第4行，这时m与n都是1，而k是3。

规则4：若有了k×10(m+2)+n，则还可以有k×10(m)+n。

例子：从第6行到第7行。这时，m=2，n=10，k=301。

可别忘了我们的公理！没有它我们哪里都去不了。所以，我们说：

我们有31。

现在，右边一列可被视为一个资格充分的算术过程了，它来自一个新的算术系统，不妨称之为310系统：

\begin{longtable}[]{@{}lll@{}}
\toprule\noalign{}
\endhead
\bottomrule\noalign{}
\endlastfoot
⑴ & 31　 & 给定的 \\
⑵ & 311　 & 规则2(m=1, n=1) \\
⑶ & 31111　 & 规则3(m=2, n=11) \\
⑷ & 301　 & 规则3(m=1, n=1, k=3) \\
⑸ & 3010　 & 规则1(m=30) \\
⑹ & 3010010　 & 规则2(m=3, n=10) \\
⑺ & 30110　 & 规则4(m=2, n=10, k=301) \\
\end{longtable}

请看，加长与缩短规则还在追着我们，这个``310系统''也被缠上了。这个东西被传送到了数的领域，使得哥德尔数时而大时而小。如果细细观察，你会发现这些规则的设计是基于一个并不深奥的思想：整数的十进制表示中那些数码的左右移动是与乘以或除以10的方幂相关联的。这一简单的结论可以推广：

中心命题：若有一条印符规则，它在任一十进制表示的数中移动、改变、删除或插入数码，那么这条规则可同样由一条算术规则来替代，后者包括对10的方幂的算术运算以及加、减等运算。

简而言之：

{用于数字的印符规则实际上等同于用于数的算术规则。}

这一简单的结论是哥德尔方法的核心，其后果是震撼人心的。就是说，一旦对一个形式系统进行了哥德尔配数，立刻就有了一组规则，使我们得到完整的哥德尔同构。结果就是，我们可以把对任一形式系统的研究------事实上即对所有形式系统的研究------转成数论工作。

WJU可产生的数

就像一组印符规则会产生一集定理，一集相应的自然数可通过重复使用算术规则而生成。这些可产生的数在数论中扮演的角色与形式系统中定理扮演的角色是一样的。当然，什么样的数是可产生的，这取决于采用的是什么规则。``可产生的''数的可产生性是相对于一个算术规则的系统而言的。例如，诸如31，3010010，31111等等这样的数可以称为``WJU可产生的''数------一个笨拙的词，或许该简称为``WJU数''，意指这些数是把WJU系统通过哥德尔配数转入数论的产物。假若我们是对pq系统进行了哥德尔配数，然后把其规则``算术化''，那么我们可以称那些可产生的数为``pq数''------如此等等。

注意，可产生的数（对任一给定了的系统）是由递归方法定义的：给定一些已知的可产生的数，我们的规则就告诉我们如何得到更多可产生的数。于是，可产生的数组成的类就是在不断地扩大，很像斐波那契数列或Q数的增长方式。任一形式系统的可产生的数构成的集合都是一个递归可枚举集合。其补集------不可产生的数构成的集合怎么样呢？也总是递归可枚举的吗？那些不可产生的数有没有什么共同的算术特征？

当你把形式系统研究换成数论研究时，就会产生这类问题。对每个算术化了的系统都可以问，``我们能否用一种递归可枚举的办法刻划不可产生的数？''这些都是数论中很困难的问题。对现在那些已经算术化了系统，这种问题都会显出是些极难解决的问题。不过，假如存在解决办法，那也必将是依靠通常那种用于自然数的一步一步的推理。而这当然也就是在我们前些章讲过的那种典范形式中进行的。TNT整个看来似乎抓住了所有有效的数学思维过程，并都压入一个紧凑的系统中了。

咨询TNT以解答有关可产生的数的问题

那么，是不是说有可能解答任何有关形式系统问题的方法都已经在这么一个形式系统------TNT------里面了呢？似乎这是可以想象的。比如看下面这个问题：

WU是WJU系统中的定理吗？

找到答案等同于确定30是否是一个WJU数。由于这是一个数论陈述，我们该可以设想，经过一番艰苦工作就能想出办法把句子``30是一个WJU数''翻译成TNT记号，其过程与想出如何把其它数论句子翻译成TNT记号是差不多的。我该立刻提醒读者，这样的翻译确实存在，但极其复杂。你也许还记得，在第八章里我指出过：即使像``b是10的方幂''这么简单的算术谓词，要转成TNT记号的形式都是非常棘手的------而谓词``b是个WJU数''则还要复杂得多！可这还是能做到的。这个问题中的b可被数字SSSSSSSSSSSSSSSSSSSSSSSSSSSSSS0替换。得到的结果将是一个硕大无朋的TNT符号串，一个谈论WU谜题的TNT符号串。就让我们称之为``无朋''吧。通过无朋及与之类似的串，TNT现在可以``用码''来谈论WJU系统了。

无朋的两重性

为了从这个对原问题的奇特的转换中得到更多的好处，我们得寻找下面这个新问题的答案：

无朋是TNT的定理吗？

我们做的只不过是把一个相对短小的串（WU）换成了另一个（硕大的无朋），把一个简单的形式系统（WJU系统）换成了一个复杂的（TNT）。虽然把问题这么修饰了一番，解决并不见得容易些。实际上，TNT中的加长和缩短规则更多更全了，所以把问题这么改装一下很可能比原问题要难解决得多。甚至你可能会说通过无朋来看WU纯属故弄玄虚。不过，我们可以在不止一个层次上看待无朋。

事实上，这一点是很令人感兴趣的：无朋有两种不同的被动意义。首先，是我们已经给出了的：

30是个WJU数。

其次，我们知道这个陈述与下述陈述有关联（通过同构）：

WU是WJU系统中的定理。

所以我们可以合法地引用后者作为无朋的第二个被动意义。这也许会显得有些奇怪，因为毕竟无朋只不过是由一些加号、括号等等TNT中的符号组成的东西。它怎么可能表达出没有算术内容的陈述呢？

事实是，它能。正如一串单音在一支曲子里既可以构成旋律也可以构成和声，正如``BACH''可以解释成名字也可以解释成旋律，正如单独一个句子可以同时是精确地从结构上描述艾舍尔的一幅画、DNA的一个节断、巴赫的一支曲子、嵌有这个句子的一个对话，无朋也可以从（至少）两个完全不同的角度来理解。发生这种事情是由于下述两个事实：

事实1：像``WU是个定理''这样的陈述可以通过哥德尔同构编码成一个数论问题。

事实2：数论陈述可以翻译到TNT系统中去。

可以说，有了事实1，无朋就是一则编了码的消息，有了事实2，用于该编码的符号就是TNT的符号。

编码与隐含意义

有一种观点认为，一则编了码的消息与未编码的消息的不同之处在于：前者仅有其自身还不能表示什么------还需要有关编码的知识。现在我们可以来反驳这种观点了。事实上，在现实中不存在什么未编码的消息。只有用较熟悉的编码编的消息和用不太熟悉的编码编的消息。若要显露一则消息的意义，就必须用某种机制或同构从编成的编码中把它抽出来。发现解码的方法可能很困难，可一旦发现了，消息就会变得水一样清澈。当编码足够熟悉的时候，它就不显得像编码了，人们于是也就忘了有一个解码机制存在。这样，那则消息就与其意义等同了。

我们这里碰到的情况就是消息与其意义几乎完全被等同，以至于我们很难接受驻存在符号中的另选意义了。就是说，我们对TNT符号的偏见如此之深，只看到TNT符号串中的数论意义（并且仅仅是数论意义），以至于很难接受某些TNT符号串是关于WJU系统的陈述。但哥德尔的同构迫使我们认识到某些TNT串中的这第二层意义。

以较为熟悉的方式解码时，无朋载有的消息是：

30是个WJU数。

这是个数论陈述，来自对每个记号用约定的方法解释。

但由于发现了哥德尔配数，以及基于其上的整个哥德尔同构，我们从某种意义上是破译了一段编码，其中关于WJU系统的消息被编成了TNT串。哥德尔同构是一种新的信息揭示者，正如古代文本的释读是信息揭示者一样。用这种新的、不那么熟悉的机制解码时，无朋载有如下消息：

WU是WJU系统的定理。

这个故事的寓意是我们以前听到过的：意义是我们在辨认出同构时自动出现的副产品。所以无朋至少有两个被动意义------也许还多！

自食恶果：对TNT进行哥德尔配数

当然事情不是到此就结束了。我们才刚开始体会到哥德尔同构的潜力。自然而然的路子就是把TNT反映其它形式系统的能力转而用在它自己身上，就像乌龟使螃蟹的唱机转而打击其自身，以及使他的高脚杯转而破坏其自身一样。为了做到这一点，我们得对TNT像对WJU系统一样进行哥德尔配数，然后把其推理规则``算术化''。哥德尔配数是很容易的。比如，我们可以建立如下对应：

\begin{longtable}[]{@{}llllll@{}}
\toprule\noalign{}
\endhead
\bottomrule\noalign{}
\endlastfoot
符号 & & 密码子 & \multicolumn{3}{l@{}}{%
关于密码子的一些有助记忆的说法} \\
0 & \ldots{} & 666 & \multicolumn{3}{l@{}}{%
洒过农药后的昆虫数目} \\
S & \ldots{} & 123 & \multicolumn{3}{l@{}}{%
后继关系：1，2，3，\ldots\ldots{}} \\
= & \ldots{} & 111 & \multicolumn{3}{l@{}}{%
视觉上的相似：横过来看} \\
+ & \ldots{} & 112 & \multicolumn{3}{l@{}}{%
1+1=2} \\
· & \ldots{} & 236 & \multicolumn{3}{l@{}}{%
2×3=6} \\
( & \ldots{} & 362 & 以2结尾 & \multirow{6}{=}{\pandocbounded{\includegraphics[keepaspectratio]{./media/OEBPS/Images/Formula-right_bracket.png}}} & \multirow{2}{=}{~} \\
) & \ldots{} & 323 & 以3结尾 \\
＜ & \ldots{} & 212 & 以2结尾 & & \multirow{2}{=}{这三对形成一个模式} \\
＞ & \ldots{} & 213 & 以3结尾 \\
{[} & \ldots{} & 312 & 以2结尾 & & \multirow{2}{=}{~} \\
{]} & \ldots{} & 313 & 以3结尾 \\
a & \ldots{} & 262 & \multicolumn{3}{l@{}}{%
与\gebfont{∀}（626）相反} \\
\textquotesingle{} & \ldots{} & 163 & \multicolumn{3}{l@{}}{%
163是质数，而``质数''的``质''首笔是撇} \\
∧ & \ldots{} & 161 & \multicolumn{3}{l@{}}{%
``∧''是序列1-6-1的``形象''} \\
∨ & \ldots{} & 616 & \multicolumn{3}{l@{}}{%
``∨''是序列6-1-6的``形象''} \\
→ & \ldots{} & 633 & \multicolumn{3}{l@{}}{%
6``蕴涵''3和3，某种意义上是这样吧} \\
～ & \ldots{} & 223 & \multicolumn{3}{l@{}}{%
2+2不是3} \\
\gebfont{∃} & \ldots{} & 333 & \multicolumn{3}{l@{}}{%
``\gebfont{∃}''看上去像个``3''} \\
\gebfont{∀} & \ldots{} & 626 & \multicolumn{3}{l@{}}{%
与a相反，而且还是6-2-6的``形象''} \\
: & \ldots{} & 636 & \multicolumn{3}{l@{}}{%
两个点，两个六} \\
标点 & \ldots{} & 611 & \multicolumn{3}{l@{}}{%
序列6-1-1的``形象''似乎意味着结束} \\
\end{longtable}

每个TNT符号都与由1、2、3和6组成的一个三位数配对，配对原则只为有助记忆。我将称这些三位数为哥德尔密码子，简称密码子。注意我没有为b，c，d及e配密码子。我们是在用简朴的TNT。这么做隐藏着的动机到第十六章你就知道了。至于最末尾关于标点的密码子，我会在第十四章解释。

现在我们可为任何TNT的串或规则重新穿扮了。作为一个例子，下面是公理1的两种记法，老的在新的下面：

\begin{longtable}[]{@{}llllllll@{}}
\toprule\noalign{}
\endhead
\bottomrule\noalign{}
\endlastfoot
626, & 262, & 636, & 223, & 123, & 262, & 111, & 666 \\
\gebfont{∀} & a & ： & ～ & S & a & = & 0 \\
\end{longtable}

正好，在一个大数中每三个数字插一个逗号，这种标准的约定与我们的密码子巧合了，这使得它们变得很``易读''。

下面是用新记法写的分离规则：

规则：若x及212x633y213都是定理，则y是定理。

最后，我们写一下上一章的一个完整的推导过程，同时用简朴TNT记法及我们的新记法写：

\begin{longtable}[]{@{}llllllllllllllllllllllll@{}}
\toprule\noalign{}
\endhead
\bottomrule\noalign{}
\endlastfoot
626, & 262, & 636, & 626, & 262, & 163, & 636, & 362, & 262, & 112, & 123, & 262, & 163, & 323, & 111, & 123, & 362, & 262, & 112, & 262, & 163, & 323 & \multirow{2}{=}{公理3} & \\
\gebfont{∀} & a & ： & \gebfont{∀} & a & \textquotesingle{} & ： & （ & a & + & S & a & \textquotesingle{} & ） & = & S & （ & a & + & a & \textquotesingle{} & ） & & \\
626, & 262, & 163, & 636, & 362, & 123, & 666, & 112, & 123, & 262, & 163, & 323, & 111, & 123, & 362, & 123, & 666, & 112, & 262, & 163, & 323 & \multicolumn{3}{l@{}}{%
特称} \\
\gebfont{∀} & a & ' & ： & （ & S & 0 & + & S & a & ' & ） & = & S & （ & S & 0 & + & a & \textquotesingle{} & ） & & & \\
362, & 123, & 666, & 112, & 123, & 666, & 323, & 111, & 123, & 362, & 123, & 666, & 112, & 666, & 323 & \multicolumn{6}{l}{%
} & \multicolumn{2}{l}{%
特称} & \\
（ & S & 0 & + & S & 0 & ） & = & S & （ & S & 0 & + & 0 & ） & \multicolumn{8}{l}{%
} & \\
626, & 262, & 636, & 362, & 262, & 112, & 666, & 323, & 111, & 262 & \multicolumn{11}{l}{%
} & \multicolumn{2}{l}{%
公理2} & \\
\gebfont{∀} & a & ： & （ & a & + & 0 & ） & = & a & \multicolumn{13}{l}{%
} & \\
362, & 123, & 666, & 112, & 666, & 323, & 111, & 123, & 666 & \multicolumn{12}{l}{%
} & \multicolumn{2}{l}{%
特称} & \\
（ & S & 0 & + & 0 & ） & = & S & 0 & \multicolumn{14}{l}{%
} & \\
123, & 362, & 123, & 666, & 112, & 666, & 323, & 111, & 123, & 123, & 666 & \multicolumn{9}{l}{%
} & \multicolumn{3}{l}{%
插入``123''} & \\
S & （ & S & 0 & + & 0 & ） & = & S & S & 0 & \multicolumn{12}{l}{%
} & \\
362, & 123, & 666, & 112, & 123, & 666, & 323, & 111, & 123, & 123, & 666 & \multicolumn{10}{l}{%
} & \multicolumn{2}{l}{%
传递} & \\
（ & S & 0 & + & S & 0 & ） & = & S & S & 0 & \multicolumn{12}{l}{%
} & \\
\end{longtable}

注意我把``加S''换成了``插入`123'\,''，因为后者才是目前合法的印符操作。

这种新记法给人一种相当奇怪的感觉。你完全失去了对意义的把握，但如果你长期与之相处，你就可以像读TNT记号一样很容易读用这种记法写的串了。你会一眼就看出哪个是良构的，哪个是非良构的。由于这完全是靠视觉完成的，你会把这当作一种印符操作------但同时，在这种记法下挑出良构的公式实际上就是挑出一类特定的整数，这一类整数当然也有算术的刻划方式了。

现在，对所有推理规则``算术化''该上到议事日程中了。现在的情况是，它们还都是些印符规则。但请注意：按照中心命题，一个印符规则实际上等价于一个算术规则。插入和移动十进制表示的数中的数字是种算术的操作，同时也是可以印符地完成的。正如在尾部拼上一个0完全等同于乘以10，每个规则都是对杂乱的算术操作进行描述的一种浓缩方式。所以，我们甚至无需寻找等价的算术规则，因为某种意义上所有的规则都已经算术化了！

TNT数：数的一个递归可枚举集合

这样看来，前面对于定理``362, 123, 666, 112, 123, 666, 323, 111, 123, 123, 666''的推导就是一系列高度曲折的数论变换，每个变换都作用在一个或多个作为输入的数之上，给出输出的数，而后者像从前一样称为可产生的数，或者更具体些，称为TNT数。有些算术规则使一个旧的TNT数以某种特定方式增大，得到一个新的TNT数；有些是减小TNT数；还有些规则是取两个旧的TNT数，分别进行一番处理，然后把两者结合起来得出一个新的TNT数------如此下去，等等，等等。而且我们不是只从一个已知TNT数起步，我们有五个初始的TNT数------当然是分别来自各个（简朴版本的）公理了。对TNT的算术化实际上与对WJU系统的算术化极其相似，只不过是有更多的规则与公理，而且显明地写出那些算术等价物会是非常烦人的------与顿悟更是毫不沾边。如果你弄懂了对于WJU系统是怎么做的，你就会发现这里的过程和那里毫无疑问是很相似的。

这种对TNT的``哥德尔化''引出了一个新的数论谓词：

a是个TNT数。

举个例子，从前面的推导我们知道362, 123, 666, 112, 123, 666, 323, 111, 123, 123, 666是个TNT数，而另一方面，大概123, 666, 111, 666不是个TNT数。

在我们看来，这个新的数论谓词通过TNT的某个带一个变元（就说是a吧）的串是可表示的。我们可以在前面加一个弯号，得到的串该是表示一个互补的命题：

a不是个TNT数。

如果我们把这第二个串中出现的所有a都替换成与123, 666, 111, 666相应的TNT数码------其中恰好含有123, 666, 111, 666个S，太长了以至于没法直接写出来------我们就会得到一个像无朋一样的TNT串，它可以从两个层次上来解释。首先，这个串说的是

123, 666, 111, 666不是个TNT数。

但因为有一个连结TNT数与TNT定理的同构，这个串还有第二层意义，就是：

S0=0不是TNT的定理。

TNT试图吞掉自己

这个意想不到的暖昧的双重性证实了TNT中有谈论其它TNT串的符号串。换句话说，我们在外面用于谈论TNT的元语言，多少是部分地在TNT内部遭到了模仿。并非TNT碰巧有这么个特点。这种事情之所以发生，是因为任何形式系统的建构都可以在N（数论）中得到反映。TNT有这个特点是不可避免的，正如唱机播放唱片时会引起震颤一样。表面上似乎震颤该是来自外部世界------比如乱蹦乱跳的小孩或皮球，但声音的产生就会有作用------这不可避免------而这副作用转回来恰恰就作用于产生它们的机制。这不是碰巧，这是逃脱不掉的副作用。这是由唱机的本质所决定的。因而，是数论的形式化的本质导致其元语言嵌在其自身中。

我们可以把这一看法冠以``数理逻辑的中心法则''这一名称，并用一个含两个步骤的图解来描述：

TNT \gebfont{⇒} N \gebfont{⇒} 元TNT

用话来说就是：TNT的一个串有一个N中的解释，而N中的一个陈述可以有一个第二意义，即作为关于TNT的陈述。

G：以编码方式谈论自身的符号串

这已经相当引人入胜了，但才只是故事的前一半。故事的其余部分涉及到对自指的浓缩，我们现在所处的阶段，相当于乌龟意识到可以造一张唱片，使那台唱机播放时震碎自己------现在的问题是：``给定一个唱机，你怎样实际琢磨出唱片上的东西？''这是件棘手的事情。

我们要找到一个TNT串------我们将称之为``G''------它是关于它自己的，确切地说它的被动意义之一是一个关于G的句子。具体说，那个被动意义将会是

``G不是TNT的定理。''

我该马上补充说明：G也有一个作为数论陈述的被动意义。像无朋一样，它容许从（至少）两个角度来解释。重要的是每个被动意义都是有效的，也是有用的，并且不会以任何方式使另一个被动意义罩上疑云。（一个播放唱片的唱机会引起自身及唱片的震颤这一事实，从任何角度看都不能抹煞那些震颤是乐音这一事实！）

G的存在就是导致TNT不完全的原因

创造出G的那个天才的方法，以及一些与TNT相关的重要概念，我们将在第十三和十四章里给出并讨论。至于现在，我们只是提前瞥一眼找到TNT中一件自指的东西会有什么后果，这会是挺肤浅的，但也挺有趣。谁知道呢，也许会爆炸的！某种意义上讲的确如此。我们来把注意力集中到下面这个明摆出来的问题：

G是否是TNT的定理？

让我们在这个问题上形成我们自己的观点，而不要依赖G关于它自己的观点。毕竟，G可能并不比一个禅师更理解其自身。G可能像无朋一样表达了一个假理。G可能像WU一样不是个定理。我们无需信赖TNT的每一个可能的串------只信赖它的定理。现在让我们在这一点上尽可能地用推理的威力来澄清问题。

我们要做我们通常都做的假定：TNT包容的推理方法是有效的，因而TNT决不会把假理当作定理。换句话说，任何在TNT中成为定理的东西都表示了一个真理。这样，若G是个定理，它表示的就是真理，即：``G不是定理''。它的自指性尽其全力打击了我们。由于充当了定理，G就不得不是个假理。根据我们的假定，TNT不会把假理当作定理，于是我们被迫得出结论说G不是个定理。这是可以的，然而它还是留给了我们一个并非不严重问题。知道了G不是定理之后，我们就得承认G表示了一条真理。这个局面揭示出TNT没能达到我们的期望------我们发现了一个符号串，它表示了真理，然而却不是个定理。惊异中别忘了G还有一个算术的解释这一事实------我们因此能把我们的发现小结如下：

发现了一个TNT的串，它毫不含糊地表示了关于自然数的某个算术性质的一个陈述；而且，通过在系统外面进行的推理，我们确定了那个陈述为真，还确定了那个串不是个TNT的定理。于是，如果我们问TNT那陈述是否是真的，TNT将既不说``是''也不说``不''。

《一首无的奉献》中乌龟的那个串是不是G的类似物？还不是。乌龟那个串的类似物是～G。这是为什么呢？好，让我们来看看～G说的是什么。它必定是与G说的相反。G说，``G不是TNT的定理''，所以～G必定是说``G是个定理''。我们可以重述G与～G如下：

G：``我不是个（TNT的）定理。''

～G：``我的否定是个（TNT的）定理。''

是～G与乌龟的串相平行，因为那个串不是说它自己，而是谈论乌龟最初提供给阿基里斯的那个串------上面多了一个弯（或少了一个弯，全在你怎么看它）。

无门来说最后一句话

无门也跨入了不可判定性这个无底洞的神秘之门，而且他对这个问题的理解比任何其他人都不差。这体现在他论及赵州的``无''时那首简洁的颂中：

狗子佛性，全提正令。

才涉有无，丧身失命！

\begin{enumerate}
\item
  \phantomsection\label{Chapter09.xhtml_B_1}
\item
  \phantomsection\label{Chapter09.xhtml_B_2}
\item
  \phantomsection\label{Chapter09.xhtml_B_3}
\item
  \phantomsection\label{Chapter09.xhtml_B_4}
\item
  \phantomsection\label{Chapter09.xhtml_B_5}
\item
  \phantomsection\label{Chapter09.xhtml_B_6}
\item
  \phantomsection\label{Chapter09.xhtml_B_7}
\item
  \phantomsection\label{Chapter09.xhtml_B_8}
\item
  \phantomsection\label{Chapter09.xhtml_B_9}
\item
  \phantomsection\label{Chapter09.xhtml_B_10}
\item
  \phantomsection\label{Chapter09.xhtml_B_11}
\item
  \phantomsection\label{Chapter09.xhtml_B_12}
\item
  \phantomsection\label{Chapter09.xhtml_B_13}
\end{enumerate}

\phantomsection\label{Part2.xhtml}{}

下篇

\chapter{异集璧EGB}

\pandocbounded{\includegraphics[keepaspectratio]{./media/OEBPS/Images/Part2_1.jpg}}

\phantomsection\label{Dialog10.xhtml}{}

\chapter{前奏曲\ldots{}}

阿基里斯和乌龟来到他们的朋友螃蟹家里，见到了螃蟹的朋友食蚁兽。介绍了一番之后，四个朋友坐下来喝茶。

\dialogname{乌龟：}老蟹，我们给你带来了一点东西。

\dialogname{螃蟹：}真太费心了，你们干嘛这么客气。

\dialogname{乌龟：}只是表示一下我们的敬意。阿基，把它拿给老蟹吧。

\dialogname{阿基里斯：}好的。老蟹，谨以此表达最美好的祝愿。愿你能喜欢它。

（阿基里斯把礼物递给螃蟹，那是个齐整漂亮的小包，方方的，很薄。螃蟹开始拆它。）

\dialogname{食蚁兽：}会是什么呢？

\pandocbounded{\includegraphics[keepaspectratio]{./media/OEBPS/Images/img_054.jpg}}

图54．莫比乌斯带Ⅱ

艾舍尔作（木刻，1963）

\dialogname{螃蟹：}马上就知道了。\kaitismall{（拆开之后，拿出了礼物。）}两张唱片！真是太好了！可是没有标签，嗯嗯------龟兄，这又是你的一个``特别礼物''吗？

\dialogname{乌龟：}如果你指的是破坏唱机的唱片，那这次不是。不过这唱片的确是特别录制的，全世界独一无二。说实在的，从来还没有人听过呢------当然除了巴赫当时演奏的时候。

\dialogname{螃蟹：}巴赫当时演奏的时候？你到底指什么？

\dialogname{阿基里斯：}噢，老蟹，等龟兄告诉了你这些唱片到底是怎么回事，你会高兴死的。

\dialogname{乌龟：}继续说继续说，阿基，告诉他怎么回事。

\dialogname{阿基里斯：}由我来说？这可不是件容易的事！我还是查查我的笔记吧。\kaitismall{（取出一张小卡片，清了清嗓子。）}啊哼。你们有兴趣听听数学上一个令人瞩目的新结果吗？你之所以能得到这些唱片，都应归功于这一结果。

\dialogname{螃蟹：}我的唱片是来自什么数学结果？其奇怪！好吧，你已经逗起了我的好奇心，我得听个究竟。

\dialogname{阿基里斯：}好吧。\kaitismall{（停顿了一小会儿，呷了一口茶，然后才开始说起来。）}你们听说过费马那个臭名昭著的``最后定理''吗？

\dialogname{食蚁兽：}我说不好\ldots\ldots 听起来怪熟的，可我想不起来了。

\pandocbounded{\includegraphics[keepaspectratio]{./media/OEBPS/Images/img_055.png}}

图55．彼埃尔·德·费马

\dialogname{阿基里斯：}内容非常简单。彼埃尔·德·费马，这位专业律师、业余数学家，在读一本丢番图的经典教本《算术》时，看到了这样一个方程式：

a\textsuperscript{2}+b\textsuperscript{2}=c\textsuperscript{2}

他立刻意识到有无穷多个a、b、c能满足这个方程，于是在书的边空上写下了下面这段臭名远扬的批注：

方程

a\textsuperscript{n}+b\textsuperscript{n}=c\textsuperscript{n}

其中a、b、c及n取正整数，仅当n=2时有解（此时将有无穷多的三元组a、b、c满足这个方程）；但n＞2时没有解。我已找到了一个精彩的证明，可是不凑巧，这里的边空太小了，我无法写下它。

\hfill\break

从三百多年前的那天起，数学家们绞尽脑汁，试图做到：或者证明费马的断言，从而维护费马的名誉------因为由于一些人怀疑他并没有真的找到他所说的那个证明，费马的名声多少有些不佳------或者否定费马的断言，找一个反例：给出四个整数a、b、c和n，其中n＞2，满足那个方程。直到最近，这两个方向上的努力都是失败的。诚然，这个定理已在许多具体的n值上被验证是成立的------即那些小于125000的n值。

\dialogname{食蚁兽：}既然没有一个真正的证明，是不是该称之为``猜想''而非``定理''？

\dialogname{阿基里斯：}严格说来，你是对的。但习惯上一直这么叫。

\dialogname{螃蟹：}有什么人最终解决了这个著名的问题吗？

\dialogname{阿基里斯：}你说着了！事实上，正是龟兄解决的，像以往一样，是靠了绝妙的一招。他不仅发现了费马最后定理的证明（因而使``定理''这个名称成为合法的，也保住了费马的名誉），还找到了一个反例，说明那些持怀疑态度的人具有良好的直觉！

\dialogname{螃蟹：}我的天！真是个革命性的发现。

\dialogname{食蚁兽：}你别让我们悬在这儿啊，到底是些什么神奇的整数满足了费马方程？我尤其感兴趣的是n的值。

\dialogname{阿基里斯：}噢，坏了！真对不住大家。简直难以相信，我竟然把写着那些值的那张巨大的纸忘在家里了！太可惜了，那张纸大得没法随身携带。我倒真希望我带来了，现在就给你们看看。不过有一件事我确实记得，也许对你们能有帮助------那个n的值是不出现在π的连分数展开式中的唯一一个正整数。

\dialogname{螃蟹：}噢，真遗憾你没把它带到这儿来。不过没什么理由怀疑你的话。

\dialogname{食蚁兽：}无所谓，写成十进制的n毫无用处。阿基不是刚刚告诉了我们怎么找到它吗。好了，龟兄，在你取得这个划时代的发现之际，请接受我衷心的祝贺！

\dialogname{乌龟：}谢谢。可我觉得比结果本身更重要的是它直接导致的实际应用。

\dialogname{螃蟹：}我正盼着知道有什么应用呢，因为我一直觉得数论是数学的王后------数学中最纯的分支------数学中没有应用的一个分支！

\dialogname{乌龟：}不只你一个人有这种想法。不过说真的，要想给出一个一劳永逸的论断，说明数学中的某一分支------或者甚至某一个定理------什么时候以及怎样在数学之外产生重要影响，那几乎是不可能的。这种事情往往没法预言------我们目前的情况就是这种现象的一个极好例子。

\dialogname{阿基里斯：}龟兄那个一箭双雕的结果创造了音响恢复领域中突破性的成果。

\dialogname{食蚁兽：}什么是音响恢复？

\dialogname{阿基里斯：}这个名称就解释了它的全部内容：从极其复杂的信息中恢复出音响信息。音响恢复的一项典型任务是，当一块石头投入湖水中之后，从湖面上散开的波纹中恢复出石头入水时的声音。

\dialogname{螃蟹：}啊唷，那几乎是不可能的呀！

\dialogname{阿基里斯：}可能的。这很像一个人的大脑所做的事：当另一个人的声带发出某种声音后，这个人的耳鼓将振动传送到耳涡中的神经纤维，他的大脑就据此恢复出那种声音。

\dialogname{螃蟹：}我明白了。可我还是弄不懂数论是从哪儿掺和进来的，或者说这一切和我的新唱片都有什么关系？

\dialogname{阿基里斯：}有关系。在音响恢复的数学中，提出了许多与某种丢番图方程解的数目有关的问题。几年来龟兄一直在努力寻求一种方法，能对目前大气里所有分子的运动进行计算，从中恢复出巴赫两百多年前演奏拨弦古钢琴时的声响，

\dialogname{食蚁兽：}那肯定不可能！那些声音已经不可挽回地消失了，永远消失了！

\dialogname{阿基里斯：}无知的人就是这么想问题\ldots\ldots 龟兄在这个问题上花了好几年的心血，终于意识到这全都取决于方程

a\textsuperscript{n}+b\textsuperscript{n}=c\textsuperscript{n}

有多少正整数的解，其中n＞2。

\dialogname{乌龟：}当然，我可以解释到底这个方程是怎么提出来的，不过我敢肯定这会让你们不耐烦。

\dialogname{阿基里斯：}结果就是，音响恢复理论预言，巴赫的声音可以从大气中所有分子的运动中恢复出来，条件是或者那个方程至少有一个解------

\dialogname{螃蟹：}太棒了！

\dialogname{食蚁兽：}妙极了！

\dialogname{乌龟：}谁能想得到！

\dialogname{阿基里斯：}刚才没说完，我说的是``只要或者存在一个解，或者证明没有解！''于是，龟兄谨慎地开始了工作，从问题的两个方面同时入手。结果是，反例的发现是找出证明的关键组成部分，因此得到一个结果后直接就导致另一个结果。

\dialogname{螃蟹：}那怎么会呢？

\dialogname{乌龟：}是这样：我指出了费马最后定理的任何一种证明------假如有这么个证明------其结构的组织可以用一个漂亮的公式描述，而这个公式凑巧有赖于某个方程的解的值。我惊异地发现后一个方程原来就是费马方程。这是形式与内容之间关系的一个有趣巧合。所以我找到反例后，要做的不过是从这个数出发构造出那个方程没有解的证明。想想看，出奇的简单。我难以想象为什么以前从未有人发现这个结果。

\dialogname{阿基里斯：}由于数学上的这个出乎意料的巨大成功，龟兄因而能实现他梦想已久的音响恢复了。老蟹得到这件礼物标志着所有抽象的工作都已化为可见的现实。

\dialogname{螃蟹：}难道这是巴赫演奏他自己的羽管键琴作品的录音？可不要蒙我！

\dialogname{阿基里斯：}我很抱歉，这可不是蒙你，这是事实！这两张唱片是一套，录的是约翰·塞巴斯第安·巴赫亲自演奏他的《平均律钢琴曲集》的全部作品。一张唱片是上集，另一张是下集，也就是说，每张唱片都有前奏曲及赋格曲共24首------每个大调及每个小调都有前奏曲及赋格曲各一首。

\dialogname{螃蟹：}这珍贵无比的唱片我一定要听听，马上就放，绝对的！哎，我该怎么谢你们二位呢？

\dialogname{乌龟：}我们已经很满足了，你为我们准备了这么好的茶。

（螃蟹从套中抽出了一张唱片，放了起来。一位羽管键琴演奏家以难及置信的高超技巧进行演奏的乐声充满了屋子，声音的保真度好得难以想象。人们甚至可以听见------或者是想象到？------巴赫一边演奏一边低吟时那轻柔的嗓音\ldots\ldots）

\dialogname{螃蟹：}你们谁要看着总谱听吗？我正好有一本《平均律钢琴曲集》，是独一无二的版本，我的一位老师在其中用文字构成了许多复杂的插图，他恰好还是位难得的书法家。

\dialogname{乌龟：}能读总谱可太好了。

（螃蟹走到他那漂亮的书橱前，打开门，取出两大册谱子。）

\dialogname{螃蟹：}给你，龟兄。我一直没能真正弄懂这个版本中所有那些精美的插图。也许你的礼物会进一步促使我理解它们。

\dialogname{乌龟：}我真希望能这样。

\dialogname{食蚁兽：}你们注意没有，在这些曲子中前奏曲总是为继后的赋格恰到好处地确定了基调。

\dialogname{螃蟹：}是的。虽然很难用语言说清楚，那两个曲子之间的确总是有着某种微妙的关系。即使前奏曲与其赋格没有共同的旋律主题，仍然会有某种两者都具有的无形的抽象特性，把它们紧密地联在一起。

\dialogname{乌龟：}前奏曲与赋格之间短暂的休止是极富于戏剧性的------在这里，单音的赋格主题即将奏出，稍后，这个主题与自己交织，形成越来越复杂的奇异优美的和声层次。

\dialogname{阿基里斯：}我懂你的意思。有好多前奏曲与赋格我还不太懂，那沉默构成的短暂插曲总让我激动不已，而我总是在那一段间歇里揣度老巴赫下面的意图。比如说，我总是想知道赋格的速度将是什么：快板还是慢板？6/8的，还是4/4的？三声部的，还是五声部的------或者四声部？然后，第一声部开始了\ldots\ldots 多么美妙的时刻！

\dialogname{螃蟹：}啊，真的，我还记得我年轻时那些逝去的日子，那时候每支新的前奏曲与赋格都使我激动，它们那新奇与优美让我兴奋不已，我常常会感受到不期而至的惊讶。

\dialogname{阿基里斯：}现在呢？那些激动都没有了吧？

\dialogname{螃蟹：}由于熟悉而消失了。任何激动不都是这样吗？不过在这种熟悉里面包含着深度，算是对消失了的激动的一种补偿。比方说，我发现总有我以前没注意到的东西。

\dialogname{阿基里斯：}有一些你以前未曾注意的主题出现？

\dialogname{螃蟹：}也许吧------尤其是当它逆转着藏在其他的几个声部中，或者从底层突然涌现出来。另外还有那些了不起的变调，充满了不可思议的奇迹，使人百听不厌，禁不住想知道老巴赫是怎么想出来的。

\dialogname{阿基里斯：}最初听到《平均律钢琴曲集》的时候，我有一阵子喜欢得都要痴颠了。那都是过去的事了，现在你说里面还有可期待的，我真是高兴------虽然这也让我感到悲哀，这个阶段对于我来说已是一去不复返了。

\dialogname{螃蟹：}噢，你不必担心你那种痴颠会彻底泯灭。年轻时那种痴颠的一个美妙之处就是它总能起死回生，恰恰就在你觉得它终于泯灭了的时候。需要的仅仅是外界恰到好处的触发，

\dialogname{阿基里斯：}真的吗？你说具体点好吗？

\dialogname{螃蟹：}比如说通过另一个人的耳朵来听，对于这另一个人来说听这支曲子是一次全新的体验------比如说这另一个是你，阿基。兴奋不知怎么传了过来，我就又一次感到激动了。

\dialogname{阿基里斯：}这可真有意思。那激动在你身体里休眠了，而你自己却无法从你的潜意识中把它勾出来。

\dialogname{螃蟹：}正是这样。在我的大脑结构中，那种激动的潜在的复苏以某种未知的方式``编了码''，我无法随意唤醒它，必须等着偶然的环境来触发。

\dialogname{阿基里斯：}我有一个关于赋格的问题，真有点不好意思开口。对我这样一个刚刚起步的赋格听众，也许你们哪位赋格欣赏专家能给我一些帮助\ldots\ldots？

\dialogname{乌龟：}要是确实能对你有所帮助，我当然愿意提供我这一点点关于赋格的知识。

\dialogname{阿基里斯：}谢谢。让我找一个合适的角度来提这个问题。你们对艾舍尔的画《魔带和立方架》熟悉吗？

\pandocbounded{\includegraphics[keepaspectratio]{./media/OEBPS/Images/img_056.jpg}}

图56．魔带和立方架

艾舍尔作（蚀版画，1957）

\dialogname{乌龟：}画的是些环状的带子，上面有小泡一样的东西，你若觉得它们是小包，它们似乎就变成了小坑------反过来也一样，是吗？

\dialogname{阿基里斯：}没错。

\dialogname{螃蟹：}我记得那幅画。那些小泡泡总是随着你看它们的角度而在凸与凹之间变来变去。没办法把它们同时看成凸的和凹的------人的大脑不知怎么不允许那样。在人们看那些泡泡时，有两种互相排斥的方式。

\dialogname{阿基里斯：}就是这样，我似乎发现我听赋格时，也有类似的两种方式。那就是：或者在一段时间里只跟着一个声部，或者，听所有声部合在一起的总体效果，不去从中区分出某个声部。我曾试过两者兼顾，结果非常沮丧，两个方式中的任何一个都会封闭住另一个。我的确没有能力跟着一个声部听下去，同时还能听总体效果。我发现我不经意地在两种方式之间跳来跳去。我并没有要这样，这多少是自发的。

\dialogname{食蚁兽：}就像你看那个魔带的时候一样，是吗？

\dialogname{阿基里斯：}是的。我想知道的是\ldots\ldots 是不是我谈起这两个方式，恰恰表明我无疑是个无知的、毫无经验的听众，甚至完全不可能抓住那种超出了我的知识范围的更深的感受方式？

\dialogname{乌龟：}不，完全不是。当然我只能说说我自己，我发现我也是在两种方式之间换来换去，虽然并没有下意识地努力控制哪种方式应该占上风。我不知道在座的哪位是否也有类似的体验。螃蟹：当然有。这是个很撩人的现象，这使你觉得那赋格的本质在你周围飘来飘去，你无法全都把握住，因为你不能让自己同时使用两个方式。

\dialogname{食蚁兽：}赋格就是这么有趣，它的每个声部自身都是一首乐曲，因而一支赋格可以被当作同时演奏的、建立在同一个主题上的几支不同乐曲的总合。怎么听取决于听众（或他的下意识），或者当作一个整体，或者当作彼此独立并且彼此和谐的几个部分。

\dialogname{阿基里斯：}你说那几个部分是``独立的''，这不太准确。它们之间肯定有某种契合，否则把它们放在一起只会是些杂乱无章的声响------这与事实差得太远啦。

\dialogname{食蚁兽：}更好的一种说法是：如果你单独听某个声部，你会发现它似乎是自足的。它可以独立存在，这就是为什么我说它是独立的。不过你也是对的，你指出了这些自足的各具意义的部分是秩序井然地彼此溶合在一起的，它们形成一个优美的整体。写出好赋格正是来自这种能力：使听众觉得有几支不同的旋律，每支旋律都好像是由于其自身的美而被写出来的，然而放到一起形成一个总体时，一点也不觉得勉强。我们所遇到的两难现象------即把一支赋格当作一个总体去听与只听它的某个声部------只是一种非常普通的两难现象的特例。许多种从较低层次搭起来的结构都具有这种两难现象。

\dialogname{阿基里斯：}噢，真的？你是说我那两种方式会有更一般的表现形式，适用于听赋格之外的情况，是吗？

\dialogname{食蚁兽：}绝对正确。

\dialogname{阿基里斯：}那怎么会呢。我猜想这与某种交替现象有关，一种把某个东西作为整体感受和作为部分的组合来感受时出现的现象。不过我只在听赋格时遇到过这种现象。

\dialogname{乌龟：}噢，天哪，都来瞧啊！我跟着音乐刚翻到这页，就见到这幅极妙的插图，正好对着这首赋格总谱的第一页。

\dialogname{螃蟹：}我以前从未见过这幅图，你为什么不让大家都看看？

（乌龟把起那本书传给大家。每个人都以自己的特有方式看着那幅图------这个离得远远的，那个凑到近前，每个人都这样或那样地歪着脑袋，困惑不解。最后，又传回到乌龟手里，他全神贯注地看着。）

\dialogname{阿基里斯：}啊，我觉得这支前奏曲马上就要结束了。不知道听了接下来的那支赋格之后，我对下面这个问题是否会有更好的理解：``以什么方式去听赋格是正确的------作为一个整体，还是部分的总和？''

\dialogname{乌龟：}仔细听仔细听，会有的！

\hfill\break

前奏曲结束了。出现了一段静默，然后\ldots\ldots{}

{[}紧接下段{]}

\phantomsection\label{Chapter10.xhtml}{}

第十章

\chapter{描述的层次和计算机系统}

描述的层次

哥德尔串G和巴赫的赋格曲都有这样一个性质：能在不同的层次上被理解。我们都熟悉这种事情。有时我们也会被搞糊涂，而另一些时候还是能毫无困难地掌握它们。例如，我们都知道人体是由大量的细胞所组成（大约有二十五万亿个），因此我们所做的一切原则上都能在细胞水平上被描述。或者说它们甚至可以在分子水平上被描述。大多数人都把这作为一个事实而接受了。我们去找医生看病，医生对我们的观察层次就要低于我们自我观察的层次。我们可以一边读着关于DNA和``遗传工程''的文章，一边呷着咖啡。我们似乎调和了这样一些关于我们自身的极不相同的图像，其实我们只不过是不把它们联系在一起。我们几乎无法把关于我们自身的微观描述和我们的自我感觉联系起来，因此我们很可能是把关于我们自身的不同表示存储在我们心智的不同``格子''中了。我们几乎从不在这些概念之间转来转去，纳闷``这两种完全不同的东西怎么可能是同一个我呢？''

在看电视的时候，屏幕上一连串画面显示着卓别林在跑来跑去，这时，我们知道事实上我们看到的不是一个人，而是平面上闪烁的一片光点。我们知道这一点，但这是埋藏在思想深处的。关于屏幕上的东西，我们有两个尖锐对立的表示，但这不会给我们造成麻烦。我们可以对其中之一置之不理，而去注意另一个------这是我们大家都用的办法。哪一个``更真实些''？这取决于你是一个人，还是一只猫、一台计算机、或一架电视机。

组块化与棋术

人工智能研究中的一个重大问题，就是要指出如何跨越这两种描述间的鸿沟，即如何构造一个系统，使它可以接收一个层次上的描述，然后从中生成另一个层次上的描述。通过探索如何给计算机编制一个好的下棋程序，这条鸿沟在人工智能中明确地表现出来了。以前人们认为------在50年代和60年代------让机器下一手好棋的关键，是使它在可能的着法序列构成的分支网络上，比任何棋术大师所能超前搜索的步数都多。可是，当这一目标逐渐达到后，计算机的棋艺并没有突飞猛进，超过人类高手的水平。事实上，一个高手在和当今最好的下棋程序对局时仍可以稳操胜券。\footnote{}

这种现象的原因实际上已讨论了许多年。在40年代，荷兰心理学家阿德里安·德·格鲁特就研究了下棋时新手和大师是怎样感知一个棋局的。用那些刻板的术语来说，他的研究结果指出大师们是以``组块''来感知棋子的分布的。有一种比直接描述每处棋子的位置层次更高的棋局描述，而大师就是以某种方式生成了这样一个棋局的心智表象。这一点已被下述事实所证实：两个人都用五秒钟时间审视棋盘之后，大师能很快地重新摆出一盘实际对局中的某个局面，而新手则要慢得多。下面这一事实是非常说明问题的：再现一个棋局时，大师们的错误涉及到把整整一组棋子放在错误的位置上，这几乎不改变对局的战略态势，但这在新手的眼里却是完全不同局面。一个决定性的事实是，在做同一个实验时，如果把棋子乱放在棋盘中的各个位置上，而不是取材于一个实际对局那么在重现这样一个随机局面的过程中，大师比新手也强不了多少。

结论是这样的：在正常的对局过程中，会有某种多次再现的局面------也就是某种模式------而大师正是对这种高层次的模式十分敏感。他和新手在不同的层次上思考，他们的概念集是不同的。几乎出乎所有人意料的发现是，大师在超前搜索时很少比新手走得更远------更有甚者，大师往往只检查屈指可数的几种可能着法！关键在于，他感知棋局的方式像一个过滤器：在观察局势时他实际上看不见坏棋------正像业余棋手在观察局势时看不见不合规则的棋步一样。任何懂一点棋术的人都已经对他的感知进行了组织，斜着走车、象走日字这样的着法不会出现在脑子里。类似地，大师级的棋手在观察棋局时已经建立了更高层次的组织，因此对他们来说，坏棋一般不会被想出来，就像对大多数人来说不合规则的棋不会被想出来一样。这可以称作对由各种可能性组成的巨大的树形结构所做的``隐式修剪''。相反，``显式修剪''是指先考虑一步棋，粗略地检查一番后，决定不再继续做进一步的检验。

这种区分同样适合于其它智力活动------例如作数学题。一个有才华的数学家往往不像庸才们所做的那样，考虑并尝试所有的错误途径，以达到所需的定理。相反，他一下就``嗅出''了有希望的途径，然后立即进行下去。

仅仅依赖于超前搜索的计算机下棋程序还没有被教会在一个更高的层次上思考，它的战略只是凭一股蛮劲进行超前搜索，企图压倒各种类型的对手。但这种办法并不灵。也许有一天，一个蛮劲足够大的超前搜索程序确实战胜了人类最好的棋手------但那也只是一个很小的智慧之果。因为我们已经揭示出了：智能紧密地依赖于为复杂的对象------如棋盘、电视屏幕、书页或画面------构造高层描述的能力。

相似的层次

通常，我们在理解一个情景时并不需要同时在头脑中容纳一个以上的层次，而且对同一个系统的不同描述一般在概念上相距甚远，就像我们前面提到过的一样。因此同时保存它们并不困难，只须把它们保存在心智世界中的不同区域里就行了。真正造成混乱的是这样一种情况：同一个系统允许两种以上不同层次的描述，而这些描述在某些方面又是彼此相似的。那时我们会发现，在考虑这种系统时，我们难以避免把层次搞混，很容易完全迷失方向。

毫无疑问，这种情况会发生。在我们考虑我们自己的心理的时候------例如，当我们试图理解人们从事各种活动的动机时------就是如此。在人的心智结构中有许多层次------当然对这个系统我们还没有很好地理解。但是，仍有上百种相互竞争的理论在说明人类行为的原因。每种理论都基于某种潜在的假设，关系到在寻找各种心理的``力''时，应当深入到这种层次结构的哪一层。由于目前我们差不多是用同一种语言描述所有的心智层次，这显然就造成了大量的层次混淆和许多错误的理论。例如，我们谈论``驱力''------但不知道这些驱力在人的心智结构中出自何处。我不想就此进一步讨论，只想说明我们进行自我认识时的混乱与下述事实有关：我们是由许多层次构成的，同时我们用重叠的语言在所有这些层次上描述我们自己。

计算机系统

在另一个领域中也有这样的情况：对同一个系统有许多层次的描述共存，而且这些层次在概念上彼此非常接近。我这里指的是计算机系统。当一个计算机程序在运行的时候，可以在若干个层次上观察它。在每一层上，描述都是用计算机科学的语言给出的，这使得所有这些描述都多多少少地彼此相似------然而不同的层次上得到的见解仍然极其不同。在最低层，描述会复杂得就像对电视屏幕上光点的描述一样。但就某些目的而言，这是最重要的视角。在最高层，描述极大程度地组块化，给人一种完全不同的感觉，尽管事实上许多概念都是最低层和最高层所共有的。高层描述中的组块类似于棋术大师的组块，也类似于屏幕图像的组块化描述：它们以简略的形式概括了若干在低层看来是不同的东西（见图57）。我们还是先别把事情搞得太抽象，先看看有关计算机的一些具体事实吧。我们来简述一下从最低层看计算机系统是什么样的。最低层？严格地说并不是，因为我没打算谈及基本粒子------但这里要谈到的是我们所愿意考虑的最低层。

\pandocbounded{\includegraphics[keepaspectratio]{./media/OEBPS/Images/img_057.png}}

图57．``组块化''的想法

把一组对象当作一个``块''后，可以把它们看成一个单位。其边界有点像细胞膜或国境线：它使得其内部的那一组对象构成一个独立的单位。根据不同的需要，可以考虑组块的内部结构，也可以不考虑。

计算机的基本成分有``存储器''、``中央处理器''（CPU）以及一些``输入输出（I/O）设备''。首先我们来谈谈存储器。存储器被分为若干称为``字''的物理上分离的单元。为了描述得具体一些，我们假设存储器共有65536个字（这是个特殊的数，2的16次方）。一个字可以进一步分为我们当作计算机科学的原子的东西------``位''。

一个典型的字中大概有36位左右。\footnote{}

\pandocbounded{
\includegraphics[keepaspectratio]{./media/OEBPS/Images/img_Chapter10_1.jpg}}

------一个36位的字------

从物理上看，一个位不过是一种磁性``开关''，它可以处于两种状态。你不妨把这两种状态称为``开''和``关''，或``X''和``O''，或``1''和``0''\ldots\ldots 第三种是最常用的。这种表示很方便，但也可能使人误解，以为计算机说到底是存储数字的。这是错的。三十六个位组成的集合不一定要被想象成数字，正像一角钱不一定要被想象成一根冰棍的价格。钱可以有各种不同的用途，这取决你怎样使用它。同样，存储器中的一个字也可以有不同的功能。当然，有时这36位确实是表示了一个二进制数。另一些时候它们也可能是表示一个电视屏幕上的36个点，或文章中的一个词。存储器中的一个字应当如何理解，这完全取决于这个字在程序中所扮演的角色。当然，它也可能扮演多种角色------像卡农曲中的一个音符一样。

指令和数据

有一种对于字的解释我还没有提到，这就是``指令''。存储器中的字里面不仅可以放被处理的数据，还可以放一些作用于数据之上的程序。有那么一组有限的固定的操作，它们可以被中央处理器------CPU------所执行，而一个字的一部分，通常是前面的几位，可以被解释成要被执行的指令类型的名字。一个解释成指令的字的其余部分代表什么呢？最常见的是，它们说明要对存储器中的哪些字进行作用。换句话说，剩下的位构成了一个指向存储器中其它字的``指针''。在存储器中，每个字都有各自的位置，像大街上的房子一样。字的位置称为``地址''。存储器里可以有一条``街''或几条``街''------它们被称为``页''。因此，要寻找一个字的地址，必须知道它的页号（如果存储器是分页的）以及它在该页中的位置，所以指令的``指针''部分就是存储器中某个字（或某些字）的地址的数字化表示。对于指针是没有限制的，因此一条指令甚至可以``指向''它自己，这样当它被执行时，将造成对自身的某些修改。

在任一给定时刻，计算机怎样才能知道该执行什么指令呢？这记录在CPU里。有一个特殊的指针，指向要被解释成指令的下一个字（即存储着它的地址）。CPU把那个字从存储器中提取出来，通过电子手段把它复制到属于CPU自身的一个特殊的字中（CPU中的字通常不叫``字''，而叫``寄存器''）。然后CPU就执行这条指令。这条指令可以要求完成任意一种类型的操作。典型的操作包括：

ADD{[}加{]}：把指令中所指出的字加到一个寄存器上。（在这种情况下，所指出的字显然被解释成一个数字。）

PRINT{[}打印{]}：把指令所指出的字作为字母打印出来。（在这种情况下，这个字显然不是解释成一个数字，而是解释成一串字母。）

JUMP{[}转移{]}：转移到指令中所指出的字。（在这种情况下，告诉CPU把那个特定的字解释成它的下一条指令。）

除非指令明确要求其它作法，否则CPU将取出下一个字，并把它解释成指令。换句话说，CPU假定自己应当像个邮递员那样，沿着``大街''走下去，依次把一个个字解释成指令。但是这种顺序工作方式可以被像JUMP（转移）这样的指令所打破，还有一些指令也有这种效果。

机器语言之别于汇编语言

上面所说的只是``机器语言''的一个粗略框架。在这种语言中，所能执行的操作类型构成了一个不可扩展的有限集合。因此所有的程序，不管多大、多复杂，都必须用这几种类型的指令构成。看一个用机器语言写成的程序，差不多可以和一个原子一个原子地看一个DNA分子相比！翻回去看看图41，那个DNA分子的核苷酸序列------想想每个核苷酸都包含二三十个原子------再想象一下为一个微小的病毒（更不要说为人了！）一个原子一个原子地写出DNA------你就会体会到，用机器语言写一个复杂的程序是什么滋味，你也会体会到，在只能与程序的机器语言描述打交道的条件下去掌握程序的进展情况，那是一种什么滋味。

但必须提到的是，在计算机刚出现的时候，程序设计是在比机器语言更低的层次上（还有更低的层次！）完成的------即把导线彼此联接，因此恰当的操作都是通过``硬性联接''而实现的。这用现代标准来看简直原始得可怕，甚至想象一下都觉得苦不堪言。但毫无疑问，初次进行这种尝试的人们一定和现代计算机的先驱者们一样曾经感到非常振奋\ldots\ldots{}

现在我们要上升到对程序进行描述的层次体系中较高的一层。这就是``汇编语言''层。汇编语言和机器语言之间差别并不大，事实上，这一改变很平缓。在汇编语言的指令和机器语言的指令之间实质上存在着一一对应。汇编语言的设计思想是要把各个机器语言的指令``组块化''，因此当你要用一条指令把一个数加到另一个数之上时，不必再写一个若干位组成的串``010111000''，而只写ADD（加）就行，然后也不必以二进制表示的形式给出地址，而是用一个``名字''指示存储器中的字。因此，用汇编语言写的程序很像用机器语言写的程序，只是容易被人读懂了。可以把用机器语言写的程序比作以晦涩的哥德尔配数法写成的一个TNT推导，而把用汇编语言写的同一程序比作以原始的TNT表示法写成的与之同构的TNT推导，后者当然比前者易懂多了。或者，回到关于DNA的想象上来，我们可以把机器语言和汇编语言之间的差别比作下面两种方法间的差别：一种是麻烦地一个原子一个原子地说明每个核苷酸，另一种是在说明一个核苷酸时仅仅给出它的名字（即``A''、``G''、``C''或``T''）。通过这种非常简单的``组块化''操作可以节约大量的劳动，虽然观念上并没有很大的改变。

翻译程序的程序

也许，汇编语言的主要意义不在于它与机器语言间并不显著的差别，而是这样一个关键的想法：程序完全可以在不同的层次上编制！想一想：硬件是被设计来``理解''机器语言写出的程序------即位序列------而不是面向字母和十进制数字的。把一个用汇编语言写的程序送进硬件会怎么样呢？这就像你想让一个细胞接受一片纸，上面的核苷酸序列是用字母写的，而不是作为一种化学构成物。细胞要一片纸有什么用呢？计算机要汇编语言写出的程序有什么用呢？

关键在这里：人们可以用机器语言写一个``翻译程序''。这个程序叫``汇编程序''，它可以接受易记的指令名、十进制数字及其它容易被程序员记住的方便的缩写，然后把它们转换成单调然而严格的位序列。用汇编语言书写的程序在被``汇编''（即翻译）之后，就可以``运行''了------或者说它的机器语言等价物可以运行了。但这只是个说法问题。哪个层次的程序在运行？如果你说是机器语言写出的程序在运行，那是绝不会错的，因为任何程序的运行都必须通过硬件------但根据汇编语言来考虑运行中的程序同样是很合理的。例如，你完全可以说``现在CPU正在执行JUMP（转移）指令''，而不是说``现在CPU正在执行`111010000'指令''。一个钢琴家在弹奏音符G-E-B-EGB时，也是在弹奏E小调和弦的一个琶音。没有理由排斥一个以高层次观点进行的事物描述。因此，可以把汇编语言写出的程序想象成与机器语言写出的程序在同时运行。我们对CPU正在做的事情有两种描述方式。

更高级的语言、编译程序和解释程序

靠计算机自己把程序从高层次翻译到低层次，这个强有力的想法在层次体系的更上一层得到了进一步贯彻。50年代初，人们用汇编语言编了几年程序之后，就认识到了一些有特点的结构往往反复出现在各种程序中。正像下棋一样，似乎存在某些基本的``定势''，它们会自然地出现在人们想要构造``算法''的时候------所谓``算法''是指人们对所要完成的操作过程的精确描述。换句话说，算法似乎有某些高层次的成分，用它们来描述操作过程，比用局限性很大的机器语言或汇编语言要容易得多，也漂亮得多。在典型情况下，一个算法的高层成分不只包含一两条机器语言指令，而是包含整整一组，并且这些指令在存储器中也不一定是相邻的。这样的一个成分可以在一个更高层的语言中用单个的项来表示------这就是一个组块。

除去标准组块------即新发现的那些可以构成所有算法的成分------之外，人们还认识到几乎所有的程序都包括有更大的组块------或者叫超组块。不同的程序有不同的超组块，这取决于程序要完成哪些高层次的任务。我们在第五章里已经讨论过超组块，用的是它们通常的称呼：``子程序''和``过程''。显而易见，对任何程序设计语言的一个强有力的扩充就是赋予它们下述能力：用已知的实体定义新的高层实体，然后用名字调用它们。这就把形成组块的运算直接建立在语言之中了。这样，就不再有一组固定的指令，要求所有的程序都直接用这组指令来编制，而是使程序员能构造自己的模块，每个模块都有自己的名字，而且可被用在程序内的任何地方，就像它们本来就是这种语言的成分一样。当然，这并没有超出下述事实：下到机器语言层，所有的事情仍是被同一组已有的机器语言指令所完成的，但它们对于处在高层的程序员来说不是明显可见的，它们是隐含的。

基于这种想法的新语言被称为``编译语言''。其中一种很早也很漂亮的叫做``Algol''，是``算法语言''的英文缩写。和汇编语言的情况不同，在Algol的语句和机器语言的指令之间不再有直接的一一对应。无疑，仍然有一种从Algol到机器语言的映射关系，但这要比汇编语言和机器语言间的关系``杂乱''得多。粗略地说，把Algol程序翻译成机器语言，有点像根据初等代数课本上的一道应用题列出方程式（事实上，从应用题转换到方程式要复杂得多，但这一比喻有助于体会把高级语言翻译成低级语言的过程中所必须完成的那些``清理''工作）。在50年代中期，称为``编译程序''的一些成功的程序被编制出来了，其功能就是完成从高级语言到机器语言的翻译。

``解释程序''也被发明出来了。像编译程序一样，解释程序也能把高级语言翻译成机器语言，但不是先把全部语句翻译完，然后执行机器代码，而是读一行然后立即执行它。这样做的优点是，使用者不必等写完一个程序后再使用解释程序。他可以一边一行一行地编，一边对它们进行检验。因此，如果把编译程序比作书面发言翻译，那么解释程序就可以被比作同声翻译。一种非常重要又很迷人的计算机语言是Lisp（``表处理''的英文缩写）。它是约翰·麦卡锡发明的，时间和Algol的发明时间差不多。从那以后，Lisp在人工智能研究人员中广泛地流行开了。

在解释程序和编译程序的工作方式之间有个很有趣的差别。编译程序先接收输入（如一个写好的Algol程序），然后产生输出（一个长长的机器语言指令序列）。至此，编译程序的任务就完成了，其输出就交给计算机去运行。与此相反，解释程序在程序员一条条地输入Lisp语句时一直在运行，每条语句马上被``就地执行''。但这并不是说每条语句是先翻译再执行的，那样的话一个解释程序就只是个一行一行进行处理的编译程序了。在解释程序中，下列操作是交织在一起的：读新的一行、``理解''它、执行它，它们同时进行。

解释程序的基本想法如下：每次新的一行Lisp语句输入，解释程序就设法处理它。这就是说解释程序开始工作，而且其中某些（机器语言的指令）被执行。当然，确切地说，执行哪一条指令，是取决于这个输入的Lisp语句的。在解释程序中有许多JUMP（转移）指令，因此新的一行Lisp可能使控制权以一种复杂的方式四处传递------向前，向后，再向前，等等。这样，每个Lisp语句都在解释程序中被转换成一条``路径''，沿着这条路径走就能得到所需的结果。

有时不妨把Lisp语句只看成一条条的数据，它们被连续地送入一个不停地运行的机器语言程序（即Lisp解释程序）。这样想往往是有帮助的，由此你会得到关于高级语言写出的程序和执行它的机器之间的关系的另一幅图像。

自举

当然，编译程序自身也是个程序，也必须用某种语言来书写。第一个编译程序是用汇编语言来书写的，而不是用机器语言，这样就可以充分利用已经建立在机器语言之上的初步成果。图58对这些复杂的概念进行了总结。

\pandocbounded{\includegraphics[keepaspectratio]{./media/OEBPS/Images/img_058.png}}

图58．汇编程序、编译程序与计算机语言的层次

汇编程序和编译程序都是用于向机器语言进行翻译的程序。这用直线表示。它们自身也是程序，开始时也需用某种语言写成。波浪线表示编译程序可以用汇编语言书写，汇编程序可以用机器语言书写。

随着复杂性的增长，人们认识到，一个部分完成了的编译程序可用于对其它部分进行编译。换句话说，一旦编译程序的核心部分写成了，这个小编译程序就能把一个大一点的编译程序译成机器语言------依此办理就能翻译一个更大的编译程序，直到最后得出整个所期望的编译程序。这个过程有个形象生动的名字：``自举''------其缘由是显然的（想象一个抓住鞋带想把自己举起来的人就行了）。这有点像一个刚刚能顺利地使用自己母语的幼儿，从那时开始，他的词汇量和流畅程度将飞速增长，因为他可以利用语言来取得新的语言。

运行着的程序的各种描述层次

编译语言一般说来并不反映运行由它所写出的程序的机器结构。这是它优于高度专门化的汇编语言和机器语言的主要特点之一。当然，在一个用编译语言书写的程序被翻译成机器语言之后，得到什么样的程序是依赖于具体机器的。因此，在描述一个正在执行的程序时，既可以独立于机器，也可以依赖于机器。正好像提到一本书中的某段时，既可以根据其谈论的内容（独立于出版社），也可以根据其页码（依赖于出版社）。

程序正常运行时，你怎么描述它或想象它的活动都没有太大关系。只有在出了毛病时，能够在不同的层次上思考才变得重要起来。例如，如果程序在某时要执行除法指令，而除数是零，程序将暂停运行，并把问题通知使用者，报告程序中出错的位置。但是，给出说明的层次往往低于程序员书写该程序的层次。下面是对一个嘎然而止的程序的三个平行的描述：

\begin{itemize}
\tightlist
\item
  机器语言层：``程序执行停止于1110010101110111位置''；
\item
  汇编语言层：``程序执行停止于遇到DIV（除）指令时''；
\item
  编译语言层：``程序执行停止于为代数表达式`(A+B)/Z'赋值的过程中''。
\end{itemize}

对于系统程序员（指那些编写编译程序、解释程序、汇编程序和其它一些供大家使用的程序的人）来说，最大的问题之一就是想办法编写查错子程序，以便给那些程序出错的用户提供高层而非低层的问题描述信息。一个有趣的颠倒是，如果一个遗传``程序''出毛病，（即出现一个突变），这时``错误''只能在高层次上显示给人们------即出现于表现型层次，而非遗传型层次。实际上，现代生物学把突变作为了解遗传过程的主要窗口之一，原因就是对这些过程可以在多个层次上进行探索。

微程序设计和操作系统

对于现代计算机系统来说，在层次结构中还有一些其它层次。例如，在有些系统中------通常是所谓``微型计算机''------其机器语言指令甚至比``把存储器中的一个数加到寄存器中的一个数上''那样的指令还要基本。使用者可以根据自己的意愿，决定需要哪些常规机器层次的指令来编程序，并且用机器提供的那些``微指令''来为所需的指令编制``微程序''。然后他所设计的这些``高级机器语言''的指令可以被做进电路中，从而成为硬件的一部分。当然不一定非要这样做。这种微程序设计使用户在一定程度上深入到通常的机器语言层之下。其结果之一就是一个厂家的一台计算机可以被``硬性联接''（通过微程序设计）成和同一个厂家------甚至另一个厂家------的另一台计算机具有相同的机器语言指令集。这个微程序化的计算机被称为在对另一台计算机进行``仿真''。

另一个层次是``操作系统''，它处于用机器语言写的程序和使用者进行程序设计的高级层次之间。操作系统本身是个程序，其功能是防止用户直接使用``裸机''（这就为系统提供了保护），并且使程序员摆脱许多非常琐碎繁杂的问题，如读入程序、调用翻译程序、运行翻译后的程序、把输出在适当的时候引向适当的通道、把控制权转交给下一个用户等等。如果多个用户同时与一个CPU``交谈''，那么操作系统将负责以某种有秩序的方式把系统的注意力从一个用户转向另一个用户。操作系统实在太复杂了，我只想通过下面的类比给一些提示。

请想想第一个电话系统。亚历山大·格雷厄姆·贝尔可以和另一间屋里的助手通电话。声音的电子传输就相当于一台没有操作系统的裸机：用电子计算！再看一个现代电话系统。你可以选择和哪一台电话接通。不仅如此，系统还可以同时控制多个不同的电话呼叫。你可以直拨、通过接线员接通、召开电话会议。你可能听到占线的嘟嘟声，也可能听到一个信号，警告你通话时间过长。你可以安装一台交换机，把附近的一批电话都联起来，如此等等。上面这个单子很惊人，你可以想想这是多么灵活，特别是和以前被当作奇迹的``裸''电话相比。现在复杂的操作系统就是为用户及其程序完成类似的通讯管理和层次转换操作的。实际上，在大脑中一定也发生着某种类似的事件：同时控制多种刺激；决定哪一种优先于其它种，优先权保持多久；紧急情况和偶发事件造成的突然``中断''，如此等等。

方便用户并保护系统

一个复杂的计算机系统中的许多层次，具有一种为用户提供``缓冲''的综合效用，使他可以不考虑许多低层活动。这些活动往往与他毫不相干。一个飞机乘客通常不想了解油箱里还有多少燃料、风速是多少、需要准备多少份饭、目的地附近的空中交通情况如何------这些事情已被交给航空公司系统中处于不同层次的工作人员来处理，而乘客只要求从一处到达另一处。同样，只有当出了什么毛病------比如他的行李没到，乘客才会意识到他以前不知道的那个使人迷乱的层次系统。

计算机是过于灵活还是过于刻板？

让人能尽可能自然地和计算机进行通讯，这一直是促使人们提高层次的主要目标之一。当然，和用机器语言写成的低层结构相比，用编译语言写成的高层结构更接近于人一般思考时所用的概念。但是，在这种使通讯更加容易的努力过程中，``自然''这个要求的一个方面被完全漠视了。事实上人际通讯中的约束远不像在人机通讯中那样刻板。例如，在寻找最好的表达方式时，我们常常说出无意义的句子片段：我们可能在说一句话的中间咳嗽；我们彼此打断对方的话；我们使用模棱两可的描述和``不规范''的句型；我们杜撰词组、歪曲意义------但我们的大部分信息一般仍能传递过去。用程序设计语言，一般说来是有着非常严格的句法的，你必须每时每刻遵守，而且没有模棱两可的词或结构。有趣的是，程序设计语言中也允许有对应于咳嗽的成分（即非本质性的或关系不大的注解），但要求它们必须以一个关键字作为开始信号（如COMMENT，即``注解''），然后终止于另一个关键字（如一个分号）。带有讽刺意味的是，即使是向灵活性作出的这样一个微弱姿态，也带有它自己的小陷阱：如果一个分号（或无论哪种表示注解结束的关键字）在一个注解内部被使用，翻译程序就会把这个分号理解成注解结束的标志，接着就是一场灾难。

假设程序中定义了一个叫``INSIGHT''（洞察）的过程，然后调用了它十七次，但第十八次由于程序员的大意，它被写成了``INSIHGT''。编译程序将徘徊不前，并刻板地打印出一个冷酷无情的错误通知，说它从来没听说过什么``INSIHGT''。编译程序发现了这样一个错误之后，它通常试图继续工作下去，但由于它没有洞察力，因此并没有理解程序员想要干什么。事实上，它可能会把这个错写的词设想成完全不同的东西，并且在那个错误的假设下工作下去。这时，一长串错误信息将雨点似地撒向程序的其余部分，因为编译程序------而非程序员------已经完全被弄糊涂了。可以想象，如果一个同声翻译正在把汉语译成俄语，突然在汉语中听到一个英语句子，然后就把后面的汉语都当成英语来翻译，这会造成多大的混乱。编译程序常常这样可怜地迷失方向。That\textquotesingle s life.

或许这听上去像是在谴责计算机，但实际上没有这个意思。从某种意义上说，这种情况是不可避免的。一旦你停下来想一下大多数人用计算机做什么，你将发现是要完成一些非常确定并精密的工作，这些工作太复杂，不适合由人来完成。如果计算机想要成为可靠的，那么它应当不带一点模棱两可地理解人希望它做的是什么。而且它必须既不多也不少地按照明确指示给它的去做。如果在为程序员提供方便的过程中，有一个程序的目的就是``猜测''程序员想做的事情，那么，下述情况是完全可以想象的：程序员设法传达了任务，但机器完全误解了。因此重要的是，对高层次的程序来说，在要求人用着合适的同时，仍应要求精确，没有歧义。

猜出程序员的意图

现在有可能设计一种容许某种不精确性的程序设计语言，以及一个能把它翻译到低层次上的翻译程序。一种可能的办法是让为这种程序设计语言服务的翻译程序设法弄懂那些``在语言的规则之外''所作的事情。但如果一种语言允许某些``犯规''，那么这种犯规已不再是真的犯规，因为它们已经被包容于规则之内了！如果一个程序员知道他可以出现某些种笔误，那么，他就会有意地利用语言这种特点，因为他知道他实际上是在刻板的语言规则范围内工作，尽管表面上看不是这样。换句话说，如果用户了解翻译程序为了他的方便而设置的所有灵活性，那么他也就知道了哪些界限是他所不能逾越的，因此在他看来，翻译程序仍然是刻板、不灵活的，虽然它和该语言的早期版本------它们没有体现``自动补偿人的错误''------相比是使人有了更多的自由。

对于那种``有弹性''的语言，可能出现的情况大概有两种：（1）用户了解构造在语言及其解释程序中的灵活性；（2）用户不了解这些灵活性。在前一种情况下，用户仍然可以用这种语言来精确地表达他的程序，因为他能预测计算机将怎样解释他用这种语言写的程序。在后一种情况下，解释过程具有一些隐蔽的特性，使得某些事件成为不可预测的（从不知道翻译程序内部工作方式的用户的角度看）。这可能导致对程序作出完全错误的解释，因此，在主要是出于速度和可行性的考虑而使用计算机的环境中，这样的语言是不合适的。

现在事实上还有第三种可能：（3）用户了解构造在语言及其翻译程序内部的各种灵活性，但由于它们为数众多，并以一种十分复杂的方式相互作用，致使用户无法断定他的程序会被怎样解释。这可能完全适用于翻译程序的设计者。他当然比别人更了解其内部情况------但他仍可能无法预测它对一个给定类型的不常见的结构将作出怎样的反应。

目前，在人工智能的研究中，主要的领域之一就是``自动程序设计''。这涉及到更高层次语言的开发------这些语言的翻译程序极其复杂，它们能做一些非常不容易的工作，例如：对例子进行概括、纠正某些笔误或语法错误、设法弄懂有歧义的描述的意思、利用一个基本用户模型来猜出用户的意图、对不清楚的事情进行提问、直接使用自然语言等等。我们希望在可靠性和灵活性之间走钢丝。

人工智能的进展即语言的进展

在计算机科学（特别是人工智能）的进展和新语言的开发之间有着非常紧密的相互联系。在过去的十年中出现了一个明显的趋势，即把新发现吸收到新语言之中。理解和构造智能的关键之一就存在于语言的不断开发和改进之中，用这些语言可以描述符号的处理过程。时至今日，已经有大约三四十种专门为人工智能研究所用的试验性语言被开发出来了。需要强调的是，应当认识到：任何可以用上述语言之一书写的程序，在原则上也能用更低层的语言来编制，但那需要人付出较大的努力，而且编出来的程序可能长得使人无法把握。每个较高的层次都没有扩展计算机的潜能，计算机的全部潜能已经存在于其机器语言指令集之中了。高层次语言中的新概念只不过是用它们的性质提示了一些方向和前景。

所有可能的程序所组成的``空间''是如此巨大，以致于没人能知道哪些东西是可能的。每种高层次语言都自然地适用于探索``程序空间''的某些特定区域，因此程序员在使用那种语言时，就被引入了程序空间的那些区域。语言并没有强迫他编写某类特定的程序，但语言为他做某些事情提供了便利。接近一个概念，一次轻微的推动，这些常常都是一个重大发现所需要的------而且这也是促使更高层的语言出现的动力。

使用不同的语言编程序就像使用不同的调式作曲，尤其是如果你是在钢琴键盘上作曲，那就更像了。如果你学过或写过不同调式的曲子，你会发现每个调式都有自己独特的感情色彩。同样，某种指法组合的使用方式对一个调式来说很自然，但对另一个调式可能就很别扭。因此，你所选择的调式对你就具有引导作用。甚至像升C和降D这样的等音调式都会以某种方式给人以不同的感受。这显示了一个记号系统会对最终产品产生相当大的影响。

\pandocbounded{\includegraphics[keepaspectratio]{./media/OEBPS/Images/img_059.png}}

图59．智能是个层次结构

为了构造智能程序，需要建造一个硬件和软件的层次结构，以解除在最低层观察一切所带来的痛苦。对同一个过程，不同层次的描述会是很不一样的，只有顶层是充分地组块化了，因此容易被人所理解。{[}摘自帕·亨·温斯顿，《人工智能》{]}

图59显示了一幅人工智能的``分层''图景，以机器元件（如晶体管）为底，以``智能程序''为顶。该图摘自帕特里克·亨利·温斯顿所著的《人工智能》一书，它代表了几乎为所有人工智能工作者所共有的关于人工智能的构想。虽然我同意人工智能应当以某种方式分层，但我并不认为通过这样少的几层就能达到智能程序。在机器语言层和真正达到智能的层次之间，我认为可能还需有十几层（甚至几十层）。每一个新的层次都基于下面一层，同时也扩展了下一层提供的灵活性。它们将是什么样的，我们现在做梦也想不出\ldots\ldots{}

妄想狂和操作系统

计算机系统中各层次间的相似性会使人经历一种奇特的层次混淆。我有一次在观看两个朋友------都是计算机方面的新手------在一个终端上玩``帕里''程序。``帕里''是个名声不好的程序，它以一种极初等的方式模拟一个妄想症患者。它能用英语吐出一些陈词滥调，其中的词汇是从程序里的一个大词汇库中选出来的。它的行为挺逼真，因为它能够判断库存的词语中，哪些在用来回答人输入的英文句子时听起来更合理。

有一次，机器的反应时间很长------``帕里''要过很久才给答复------我向朋友们解释说，这可能是由于分时系统负担太重了。我告诉他们，只要打入一个特殊的``控制''字符，就可以进入操作系统，查出有多少用户在使用机器，而这在运行``帕里''时是看不到的。我的一个朋友按下了这个控制字符，转眼之间，某些关于操作系统状态的内部数据把``帕里''在屏幕上显示的一些词冲掉了。``帕里''对此一无所知：它的``知识''仅仅限于赛马和赌博------不是关于操作系统、终端和特殊控制字符的。但对我的朋友们来说，``帕里''和操作系统都不过是``计算机''------一个神秘、遥远、难以名状的实体，它可以对他们输入的信息作出回答。因此就完全可以理解为什么我那朋友这个时候会冒失地用自然语言打入下列问题：``为什么你把屏幕上的东西冲掉了？''我的朋友们不清楚``帕里''对于它下面的操作系统一无所知。``你''知道关于``你自己''的一切，这种想法在处理人际关系时为大家所熟知，因此很自然地会推广到计算机上------它毕竟还是具有足够智能的，它都能和他们用自然语言``交谈''！他们那个问题就像是问一个人：``你今天为什么制造了这样少的红血球？''人们不了解他们身体的那个层次------所谓``操作系统层次''。

这种层次混淆的主要原因是，人和计算机系统的所有层次之间的通讯都是在同一个终端的同一个屏幕上进行的。虽然说我的朋友似乎有点过于天真了，但若是一个复杂系统的几个层次同时出现在一个屏幕上，即使是有经验的计算机工作者也常常会犯类似的错误。他们忘了是在和``谁''谈话，于是输入了某些在那个层次上毫无意义的命令。因此，似乎应当要求系统本身能对层次进行区分------依据命令在哪一层能``有意义''来对它进行解释。不幸的是，这种解释要求系统具有大量的常识，还要有关于程序员的总体意图的完备知识------这两者都需要更多的人工智能，而这是目前还不具有的。

软件和硬件之间的界限

人们还常常会被某些层次上的灵活性和另一些层次上的刻板性给搞糊涂。例如，有些计算机具有很好的文本编辑系统，可以允许文本片段从一种格式中``倾注''到另一种格式中，就像液体可以从一个容器中被倾注到另一个容器中一样。窄页可以变成宽页，反过来也可以。它们既然有这么大的本事，人们就会期望从一种字体变成另一种------如从宋体变成楷体------也该是同样的容易。当然如果屏幕上只能显示一种字体，那这种转换就不可能了。也可能这种转换对屏幕是可行的，但在打印机上无法打印------还可能出现相反的情况。和计算机打交道久了，人常常被宠坏了，总觉得什么事情都能编程序解决：打印机不该刻板得只有一种或有限的几种字符集------字形应当是可以由用户指定的！但是一旦达到了这种灵活性，人们又会因打印机不能使用多种颜色的墨水而烦恼，或指责它不能接受任何形状和大小的纸张，或抱怨它不能在被打碎后再把自己装起来\ldots\ldots{}

麻烦的事情在于，用第五章的话来说，所有这些灵活性都必须有个``终了''。在所有灵活性的下面必须是一个硬件层，而它是没有灵活性的。它可以藏得很深，上面的各个层次又提供了大量的灵活性，以致于没有多少用户能感到硬件的限制------但这种限制的存在是不可避免的。

软件和硬件间有区别，这是公认的，但到底是什么呢？就是程序和机器之间的区别------即长长的复杂的指令序列和执行它们的物理机器之间的区别。我喜欢把软件想象成``你可以通过电话线传送的任何东西''。一架钢琴是硬件，而乐谱是软件。一部电话机是硬件，而电话号码是软件。这种区分是很有用的，但也不是总能分得这样清楚。

我们人类同样有``软件''和``硬件''两个方面，它们的差别对我们来说是根深蒂固的。我们已经习惯于自己生理上的刻板性：我们事实上不能凭愿望治愈自己的疾病，或者长出任何颜色的头发------只举这两个简单的例子就够了。但是，我们能够为我们的心智``重新编程序''，这样我们就可以在一个新的概念框架中行动了。我们心智的惊人的灵活性似乎和下述观念水火不相容：我们的大脑一定是由具有固定规则的硬件构成的，不能被重新编程序。我们无法使我们的神经元兴奋得再快些或再慢些，我们无法重建它们之间的联系，我们无法重新设计一个神经元的内部结构，我们无法对硬件进行任何选择------但是，我们仍能控制自己的思维方式。

不过，显然我们的思维的某些方面是在我们的控制范围之外的。我们无法单凭愿望使自己聪明起来；我们无法以自己想要达到的速度去学会一门新的语言；我们无法使自己的思维速度比实际情况更快一些；我们无法使自己同时思考好几件事情；如此等等。这些基本的自我认识往往是不言而喻的，就像我们意识到周围有空气存在一样。我们从不真的费神去考察导致我们心智的这些``缺陷''的东西是什么：那就是我们的大脑组织。设想一些途径来调和心智的软件和大脑的硬件，这就是本书的一个主要目的。

中间层次和天气

我们已经看到，在计算机系统中有一些定义得很严格的层次，在每个层次上都能对一个正在运行的程序的操作进行描述。因此，不存在单一的低层次和单一的高层次------而是各种程度的低层次和高层次。中间层次的存在是否就是那些具有高低层次之分的系统的一般特征？一个例子，可以考虑这样一个系统：其``硬件''是地球的大气层（不是很硬，不过没关系），而其``软件''是天气情况。同时追寻所有分子的运动大概是一种对天气的非常低层次的``理解''方式，这很像在机器语言层观察一个巨大而又复杂的程序。显然这是超出人的理解能力的。但我们仍能用一些人类特有的方式来观察和描述天气现象。我们关于天气的组块化观点是基于层次很高的现象的，如雨、雾、雪、飓风、冷锋、季节、气压、信风、急流、积雨云、雷暴、逆温层等等。所有这些现象中涉及到的分子数量都达到了天文数字，它们以某种方式协同动作，造成了大尺度的天气趋势。这有点像读用高级语言描述的天气。

是否有某种东西类似于用中间层次语言来观察天气，比如像汇编语言那样？例如，是否存在非常小的局部范围内的``小型风暴''，就像人们偶尔见到的那种能搅起一股至多一米来宽尘土的小旋风那样？一个局部地区的阵风是否就是一种中间层次的组块，在一个高层天气现象的形成过程中发挥了作用？还是说我们没有切实可行的方法把关于这种现象的知识结合起来，以构造一个对天气的更容易理解的解释？

又有两个问题出现在我的脑海里。第一个问题：``是不是我们在自己的尺度上感知到的天气现象------一股龙卷风、一场旱灾等------只不过是中间层次的现象，是某种巨大、缓慢的现象的一部分？''如果是这样，那么真正的高层次天气现象应当是全球性的，而且采用地质学的时间尺度。冰河时代应当算是一个高层次天气事件。第二个问题：``是不是有这样的中层次天气现象：它们尚未被人所感知，但是一旦人们感知到它们，就会对天气的现状产生更深入的理解？''

从龙卷风到夸克

上面最后一种设想听起来有点离奇，但绝非不着边际。我们只须考察一下``硬科学''当中最硬的一个------物理学------就能发现这样的一些特殊系统，它们可以用相互作用的``成分''来解释，而这些成分本身却是不可见的。在物理学中，正像其它任何学科一样，一个``系统''是一组相互作用的``成分''。在我们了解的大多数系统中，这些成分在相互作用中保持着自身的特性，因此我们仍能在系统内部观察到这些成分。例如，当一队橄榄球运动员聚在一起的时候，每个运动员都保持着自身的独立性------他们并没有融化进某个复合实体，从而消失了自己的独立性。但是------这是很重要的------在他们的大脑中进行着某些过程。这些过程是被球队这种环境所诱发的，否则它们不会发生。因此，在某种微弱的程度上可以说，这些队员在成为一个更大的系统（即球队）的一部分时，他们自身的个性已经有所改变。这种系统被称为``几乎可分解系统''（这个名称来自赫·西蒙的文章《复杂性的结构》）。这样的系统包含相互作用较弱的模块，每个模块在相互作用过程中保持自己独有的个性，但变得有些不同于它在系统之外时的状况，以此保证了系统的聚合性行为。物理学所研究的系统通常都是这种类型的。例如，可以认为原子中包含一个原子核，它所带的正电荷把若干电子俘获到``轨道''上，或者说使它们处于约束状态。被约束的电子很像自由电子，尽管它们处在一个复合对象的内部。

物理学中所研究的另一些系统和比较简单明确的原子恰成一鲜明对照。这些系统中包含着极强的相互作用，结果是其成分都被系统所吞没，丧失了它们的部分或全部个性。其中一个例子就是原子核。原子核通常被描述成``由质子和中子所组成''。但把作为其成分的粒子拉到一起的力实在太强了，因此这些粒子的存在方式完全不同于其``自由''形态（即它们在原子核外时所处的形态），而且事实上一个原子核的活动在很多方面都像单个的粒子，而不像一组相互作用的粒子。当一个原子核分裂时，常常释放出质子和中子，但也会释放出其它粒子，如π介子和γ射线都是通常会产生的。是所有这些不同的粒子在原子核分裂前就都实实在在地存在于其中，还是说它们只不过是原子核分裂时迸溅出的``火花''？或许回答这样一个问题的尝试本身就是没有意义的。在核子物理学的层次上，是存储了产生``火花''的潜力，还是存储了实在的粒子，这两者的差别不是很明显的。

因此，原子核就是这样一个系统：它的成分虽然在其内部是不可见的，但可以被拉出来变为可见的。但是，还有一些更反常的情况，例如把质子和中子看成系统就会如此。质子和中子都被设想成是由三个``夸克''构成的，而夸克是一种假想的粒子，它们可以三三两两结合成许多已知的基本粒子。但是，夸克间的相互作用实在太强了，致使它们不仅不能在质子和中子内部被观察到，而且甚至根本就拉不出来！这样，虽然夸克在理论上有助于理解质子和中子的某些性质，它们自身的存在可能永远不会被单独地确认。在这里我们得到了``几乎可分解''系统的反义词------这种系统是``几乎不可分解''的。但不可思议的是，关于质子和中子（以及其它粒子）的一个基于夸克的理论具有可观的解释力。使用``夸克模型''，可以很好地定量说明许多涉及到假设由夸克构成的粒子的实验结果。

超导：一个关于重正化的``悖论''

在第五章中，我们已经讨论过被重正的粒子从它们的裸核中产生的过程，即通过虚粒子递归复合地相互作用。一个被重正的粒子既可以被看成是这个复杂的数学建构，也可以看成单一的实在团块。对粒子的这种描述方式所导致的最奇特也最有戏剧性的结果之一，就是为著名的``超导''现象所提供的解释。超导就是指在极低的温度下，电子可以在某些固体中无阻抗地流动。

据发现，固体中的电子是由它们自己与一些称为``声子''的奇特的振动量子之间的相互作用所重正的（声子自己也被重正了！）。这些重正后的电子称为``极化子''。计算表明，在极低的温度下，两个自旋相反的极化子将开始互相吸引，而且确实能以某种确定的方式约束在一起。在适当的条件下，所有载流极化子将两两结合，形成``库珀对''。具有讽刺意味的是，这种结合的产生完全是由于电子------成对的极化子的裸核------相互之间的电斥力。与电子相反，库珀对之间既不相吸也不相斥，结果一个库珀对可以自由地通过一块金属，就像这块金属是真空一样。如果你把对这样一块金属的数学描述从以极化子为基本单位转换成以库珀对为基本单位，你就能得到一组明显地简化了的方程式。这种数学简单性就是物理学家的认识方式，他们知道形成库珀对这样的``组块''对观察超导来说是一种很自然的方式。

这里我们就有若干个层次的粒子：库珀对自身、构成它的两个自旋方向相反的极化子、构成极化子的电子和声子，然后，在电子内部，还有虚拟的光子和正电子，如此等等。我们可以观察某个层次，看到在那里出现的现象，而这些现象可以通过对下面几个层次的理解而得到解释。

``隔　　离''

类似地，同时也是幸运地，人们不必了解关于夸克的全部情况，仍能对那些大概是由它们构成的粒子进行深入的理解。这样，一个核物理学家可以使用基于质子和中子的理论，而不必考虑夸克理论及其竞争对手。核物理学家有一幅关于质子和中子的组块化图像------一个从低层理论中推导出的描述，但这并不需要了解低层理论。同样，一个原子物理学家有一幅从原子核理论中推导出的关于原子核的组块化图像。然后一个化学家有一幅关于电子及其轨道的组块化图像，并且构造了关于小分子的理论。这一理论可以以组块化的形式被分子生物学家所接收，他们有关于小分子如何挂在一起的直觉，但其技术专长是在超大分子及其相互作用的领域之中。然后细胞生物学家又有一幅组块化的图像，其中的基本单位是分子生物学家所思考的，而细胞生物学家只是设法用这些基本单位说明细胞相互作用的方式。问题很清楚。每个层次都是在某种意义下被``隔离''于下面的层次的。这个生动的说法也是西蒙所提供的，这使我们想起潜水艇建造水密舱的方式：如果一个部分损坏了，水会涌进来，但只要关闭舱门，损坏的舱室就被隔离开了，这样可以防止乱子的蔓延。

在科学的层次结构中，层次间总是有些``渗漏''的，因此化学家无法完全不考虑低层的物理学，生物学家也无法完全不考虑化学，但是，两个相距较远的层次间几乎是没有渗漏的。这就解释了为什么人们可以直觉地理解别人，而无须懂得夸克模型、原子核结构、电子轨道的性质、化学键、蛋白质构造、细胞中的细胞器、细胞间的通讯方式、人体各种器官的生理学、或器官间复杂的相互作用。一个人所需的一切不过是一个关于最高层活动方式的模型。众所周知，这样的模型是非常实际也非常有效的。

组块化和确定性之间的利弊权衡

然而，组块化模型可能有个很大的缺点：它通常不具有精确的预测力。这就是说，通过使用组块化模型，我们不必力不从心地把人看作夸克（或某种更低层的东西）的集合，但显然这种模型只能让我们或然地估计别人的感觉、别人对我们的言行的反应，如此等等。简而言之，在使用组块化的高层模型时，我们为简单性牺牲了确定性。尽管我们不能确定人们将会对一个笑话作何反应，我们说这个笑话时仍预期着他们的某些行为，诸如笑或者不笑------而不是，比如说，爬到最近的一根旗杆上去。（但禅宗大师可能会这样干！）一个组块化模型是：定义一个``空间''，并预期着行为会落入其中，而且描述了行为落在该空间的不同区域的概率。

``计算机只能做你告诉它去做的事''

上述想法既能用于复合的物理系统，也能用于计算机程序。有一种老生常谈的说法，即``计算机只能做你告诉它去做的事。''这在某种意义下是对的，但并没有说到点子上：你不能预先知道你告诉计算机去做的事会导致什么结果，因此它的行为可能像一个人那样令人迷惑、使人惊讶、无法预测。你一般事先知道它的输出会落在哪个空间里，但你不知道其细节。例如，你可以写一个程序来计算π的前一百万位。你的计算机会以比你快得多的速度打出前四十位------计算机超过了它的程序员，但这不会导致悖论。你事先知道输出会出现在哪个空间------即0到9之间的数字空间------这就是说，你有一个关于程序行为的组块化模型，但假设你连剩下的也都知道了，那你也就不会去写这个程序了。

说这句老话过时了，还有另外一层意思。它涉及到这样一个事实：你编程序的语言层次越高，你越无法精确地知道你让计算机所干的事！一层一层的翻译把一个复杂程序的``用户界面''和实际的机器指令分隔开了。在你思考和编程序的层次上，你的语句可能更类似于陈述句和建议，而不太类似于祈使句和命令。一般说来，一个高层语句的输入所诱发的内部过程对你来说是不可见的，正好像在吃饭的时候，你用不着明确意识到由此触发的消化过程。

``计算机只能做人告诉它去做的事''，这是由洛芙莱丝命妇在她有名的回忆录中首先提出的。这个观点一直很流行，而且往往联系于另一个观点：``计算机不能思维''。关于后者，我们将在后面的章节中再来讨论，那时我们会更老练些。

两类系统

在下述两类由许多部分构成的系统之间存在着一种重要的区别。在一类系统中，某些部分的行为趋向于抵消另一些部分的行为，结果是低层发生的事件影响不大，因为大多数事件都将导致类似的高层行为。这种系统的一个例子就是煤气罐，它里面的分子以非常复杂的微观方式彼此碰撞，但总的效果，在宏观上看，它是一个非常平静、稳定的系统，具有一定的温度、压力和容量。而在另一类系统中，单个低层事件的效果将被放大成一个巨大的高层后果。克郎棋就是这样的一个系统，其中棋子间的碰撞对整个局势的变化会产生重要的影响。

计算机是这两类系统的一个复杂的结合。它有像导线这样的构件，其行为方式是可以准确预测的，其导电过程服从欧姆定律这样一个非常精确的组块化定律，这类似于煤气罐里的煤气所服从的定律。因为导线的行为取决于统计效应，其中成千上万的随机事件彼此抵消，导致了整体行为的可预测性。计算机中也有像打印机这样的宏观构件，它们的行为完全决定于电流的精确模式。打印机的打印结果可不是被大量相互抵消的微观效应造成的。事实上，对大多数计算机程序来说，程序中每一位的值都对打印的结果起着关键性的作用。如果任何一位被改变，输出也会大幅度地改变。

那些仅仅由``可靠的''子系统组成的系统------即其子系统的行为都可以从组块化的描述中可靠地被预测出来------在我们的日常生活中发挥着无法估量的重要作用，因为它是稳定性的柱石。我们可以相信墙不会倒、人行道今天和昨天通向同一个地方、太阳会发光、钟表会指示正确的时间，如此等等。这种系统的组块化模型实际上是完全确定的。当然，另一类系统在我们的生活中也起了很大作用，它们的行为是可变的，取决于某些我们无法直接观测的微观参数------而且这些参数往往为数众多。对这样的系统来说，我们的组块化模型必须以操作空间的形式出现，并且能对结果落入空间的不同区域的可能性给出概率估计。

正像我已经指出的那样，煤气罐是一个可靠系统，其原因是许多相互抵消的效应最终遵从精确的、确定性的物理定律。这种定律是组块化的，它们把罐里的煤气看作一个整体，不管它的组成部分。而且，对煤气的微观和宏观描述使用的是完全不同的词汇。前者需要说明其中每一个分子的位置和速度，而后者只需要说明三个新引进的量：温度、压力和容量，并且前两个量甚至根本没有微观对应物。把这三个参数联系在一起的简单数学关系------pV=cT，其中c是个常数------是一条既依赖于低层现象，又独立于低层现象的定律。说得再明白一点，这条定律可以从主宰着分子层次的定律中被推导出来，在这个意义上它依赖于低层。另一方面，这条定律允许你完全不考虑低层，如果你愿意的话。在这个意义上它又独立于低层。

重要的是，要认识到高层定律是不能用低层描述的词汇表里的词来叙述的。``压力''和``温度''是一些新词，它们所代表的经验是不能在低层次上传达的。我们人是直接感知温度和压力的，我们的构造就是如此，因此我们能发现这一定律也就不足为奇了。但如果有一种生物仅把气体看作纯理论的数学结构，那它们在发现这一定律时就必须具有一种综合出新概念的能力。

旁效现象

在本章快要结束的时候，我想讲一段与一个复杂系统有关的往事。一天，我在和两个系统程序员谈论我所用的计算机。他们说操作系统在应付三十五个以内的用户时并不感到困难，但在达到三十五个用户左右时，反应时间突然一下子长了起来，慢得好像你都可以退出系统回家去等着。我开玩笑说：``这好办------看看操作系统里哪个地方存着`35'这个数，然后把它改成`60'不就行了！''大家都笑了。当然，问题在于根本没有这种地方。那么这个临界值------35个用户------是从哪里来的呢？答案是：它是整个系统组织的一个可以看见的结果------一个``旁效现象''。

类似地，你可以去问一个短跑运动员：``你能用10秒5跑100米，那10秒5这个数字存储在你身体里的什么地方？''显然，哪里也没有。这个时间是他的身体构造和他的反应时间的一个结果。在他跑的时候，有成千上万个因素在相互作用。这个时间是完全可以再现的，但它却没有存储在他体内的任何地方。它分布在他身体的所有细胞中，只有在奔跑的行动中才能表现出来。

旁效现象随处可见。我们都知道在围棋里，``两个眼可以活''。这不是写在规则中的，但这又是运用规则的结果。对人脑来说，``上当受骗''就是这样一种性质。你容易上当吗？是不是你脑子里有个``上当中心''？如果是这样的，能不能找个神经外科医生，打开你的头骨做一个精巧的手术，以后你就会少上当了？如果你相信这种观点，那你可就太容易上当了，或许你应当考虑去做这样一个手术。

心智之别于大脑

在后面几章中，我们将讨论大脑，检验一下是否可以理解大脑的顶层------心智------而不必先去理解那些它既依赖又不依赖的较低层次。是否存在这样一些思维定律：它们可以被``隔离''于那些控制着微观的脑细胞活动的低层定律？是否能把心智从大脑上``撇出来''------就像把油从汤的表面上撇出来那样------然后移植到其它系统中去？或许根本就不可能把思维过程分解到整齐的、模块化的子系统中去？大脑是更像一个原子，还是更像一个重正后的电子，或一个原子核、一个中子、一个夸克？自我意识是一种旁效现象吗？为了理解心智，我们是否一定要深入到神经细胞的层次？

\begin{enumerate}
\item
  \phantomsection\label{Chapter10.xhtml_B_1}
\item
  \phantomsection\label{Chapter10.xhtml_B_2}
\end{enumerate}

\phantomsection\label{Dialog11.xhtml}{}

\chapter{\ldots 蚂蚁赋格}

\ldots\ldots 然后，那首赋格曲的四个声部一个接一个地插进来。

\dialogname{阿基里斯：}我知道你们大家肯定不会相信：这个问题的答案其实近在眼前，就藏在这张图里。全部答案就是一个字------但是个无比重要的字：``无''！

\dialogname{螃蟹：}我知道你们大家肯定不会相信：这个问题的答案其实近在眼前，就藏在这张图里。全部答案就是一个词------但是个无比重要的词：``整体论''！

\dialogname{阿基里斯：}你只消盯着它瞧上一会儿，就一定能看出来。这简直同青天白日一样清楚，这张图的意思是``无''，不是``整体论''！

\dialogname{螃蟹：}对不起，我的眼睛可不是一般地好。请再看一看，然后告诉我这张图是不是有我认为它有的那种意思！

\dialogname{食蚁兽：}我知道你们大家肯定不会相信：这个问题的答案其实近在眼前，就藏在这张图里。全部答案就是一个词------但是个无比重要的词：``简化论''！

\dialogname{螃蟹：}你只消盯着它瞧上一会儿，就一定能看出来。这简直同青天白日一样清楚，这张图的意思是``整体论''，不是``简化论''！

\dialogname{阿基里斯：}又一个上当受骗的！不是什么``整体论''，也不是什么``简化论''，``无''才是这张图的意思，这再清楚不过了。

\dialogname{食蚁兽：}对不起，我的眼睛可不是一般地棒。请再看一看，然后告诉我这张图是不是有我认为它有的那种意思！

\dialogname{阿基里斯：}你们没看出来吗？上面的图案是由四条杠杠组成的，每一条杠杠都是汉字的一个笔划。

\dialogname{螃蟹：}它是由四条杠杠组成的，这点你说对了，可这些杠杠是怎么回事你却没搞对。那两条横着的杠杠分别是由两个同样的词``整体论''组成的，那条斜的和那条弯的杠杠也是由许许多多的这个词组成的，只是字要小些。这两部分的字为什么大小不一样，我不清楚，但我知道我所看到的是``整体论''，这跟青天白日一样清楚。我不明白除此之外你怎么还能看出些别的东西来。

\dialogname{食蚁兽：}它是由四条杠杠组成的，这点你说对了，可这些杠杠是怎么回事你却没搞对。那两条横着的杠杠分别是由许许多多个同样的词``简化论''组成的，那条斜的和那条弯的杠杠也分别是由三个这个词组成的，只是字要大些。这两部分的字为什么大小不一样，我不清楚，但我知道我所看到的是``简化论''，这跟青天白日一样清楚。我不明白除此之外你怎么还能看出些别的东西来。

\dialogname{阿基里斯：}我现在知道是怎么回事了。你们二位各自看到的那两个词，要么是对方看到的那个词的构成成份，要么是由对方看到的词所组成。那两条横着的杠杠确实是两个``整体论''，可这两个词全都是由用更小的字写成的许许多多个``简化论''组成的。而那两条或斜或弯的杠杠，呈一种互补的形式，则是三个``简化论''，而这三个``简化论''全是由用更小的字写成的许许多多个``整体论''组成的。这一下问题就解决了，而你们俩刚才那通愚蠢的争吵，全都是只见树木不见森林。你们明白了吗？要想正确理解这张图，就得超越``是`整体论'还是`简化论'\,''这一问题，这张画的意思应该是``无''，像你们俩这样在刚才那个问题上争来争去是毫无益处的。

\dialogname{螃蟹：}你所讲的这张图的构成我现在也看出来了，阿基，不过你说的什么``超越''之类的话，我一点也不明白。

\pandocbounded{\includegraphics[keepaspectratio]{./media/OEBPS/Images/img_060.png}}

图60．``无之图''

刘友庄绘

\dialogname{食蚁兽：}你所讲的这张图的构成我现在也看出来了，阿基，不过你说的什么``无''之类的话，我一点也不明白。

\dialogname{阿基里斯：}要是你们愿意先答应我一个要求，告诉我``整体论''、``简化论''这些怪词儿都是什么意思，我就给你们解释我刚才的话。

\dialogname{螃蟹：}整体论是你能在这个世界里找得到的最普通不过的事。它是说``整体大于各个部分的总合。''只要是精神正常的人就不会反对整体论。

\dialogname{食蚁兽：}简化论是你能在这个世界里找得到的最普通不过的事。它是说``如果你理解了一个整体的各个部分，以及把这些部分`整和'起来的机制，那么你就能够理解这个整体。''只要是精神正常的人就不会反对简化论。

\dialogname{螃蟹：}我就反对简化论。我要你给我讲讲，比方说，如何用简化论的原理来理解大脑。任何对大脑的简化论解释都不可避免地在说明``大脑所具有的意识究竟从何而来''这一问题时显得无能为力。

\dialogname{食蚁兽：}我就反对整体论。我要你给我讲讲，比方说，对蚁群的整体论描述如何比对该蚁群内蚂蚁的分别描述、它们各自的角色，以及相互间的关系等，更能说明事情的本来面目。任何对蚁群的整体论解释都不可避免地在说明``蚁群所具有的意识究竟从何而来''这一问题时显得无能为力。

\dialogname{阿基里斯：}打住吧！我刚才的那番话难道是为了引起一通新的争论吗！好啦，我现在明白你们在争论什么了，我看我对``无''的解释一定会对你们有很大的帮助。你们知道，``无''是古代禅宗对某些问题的回答，用来废问这个问题。我们现在的问题似乎是：``世界应以整体论原理来理解还是应以简化论原理来理解？''\,``无''这一答案把这一问题的前提，即``二者必居其一''这一规定给否定了。通过废问这个问题，揭示出更广泛的真理：存在着某种更大的领域，使得整体论与简化论的解释都适合。

\dialogname{食蚁兽：}荒谬绝伦！你这个``无''跟火车的``呜呜''声一样无聊。我不要听禅宗的什么胡言乱语。

\dialogname{螃蟹：}荒唐透顶！你这个``无''跟轮船的``呜呜''声一样无聊。我不要听禅宗的什么乱语胡言。

\dialogname{阿基里斯：}哦，天哪！我们现在谁也说不服谁。你为什么一直这么奇怪地保持沉默，龟兄？这叫我很不自在。你肯定能帮我们解脱这种困境，是吧？

\dialogname{乌龟：}我知道你们大家肯定不会相信：这个问题的答案其实近在眼前，就藏在这张图里。全部答案就是一个字------但是个无比重要的字：``无''！

（就在他说这话的时候，那首赋格曲的第四个声部也加入进来了，正好比第一声部低一个八度。）

\dialogname{阿基里斯：}噢，龟兄，这次你可叫我失望了。我刚才还以为像你这样一位事事都比我看得透的人，准能解决这个窘人的难题呢------可很明显，你并不比我看的更远。也好，我想我应该为这次能跟龟兄看的一样远而高兴才是。

\dialogname{乌龟：}对不起，我的眼睛可不是一般地尖。请再看一看，然后告诉我这张图是不是有我认为它有的那种意思。

\dialogname{阿基里斯：}当然有！你只不过重复了我的发现。

\dialogname{乌龟：}也许在这张图中，``无''所在的地方比你所想象的层次要低，阿基------低一个八度（形象点说）。不过我看我们可以把这个争论抽象化。我想听听阐述得更清楚些的整体论和简化论的观点，然后我们才会有根据下判断。比方说，我很愿意听一听对蚁群的简化论描述。

\dialogname{螃蟹：}也许食蚁兽大夫能告诉你一些他这方面的经验。他毕竟是这一领域里的专家。

\dialogname{乌龟：}我确信我们有很多东西要向您请教，食大夫。您能从简化论的角度对我们谈谈有关蚁群方面的事吗？

\dialogname{食蚁兽：}愿意效劳。正像老蟹对您说过的那样，我的职业使我对蚁群有着比较深入的了解。

\dialogname{阿基里斯：}可以想象！食蚁兽这一职业看来同作一名蚁群专家是一回事！

\dialogname{食蚁兽：}对不起。``食蚁兽''不是我的职业，而是我的种属。我的职业是蚁群外科医生。我的专长是用外科切除手术来治疗蚁群的神经错乱。

\dialogname{阿基里斯：}噢，是这样。可你所说的蚁群的``神经错乱''是什么意思？

\dialogname{食蚁兽：}我的大多数病人患的都是失语症。你知道，一群群的蚂蚁不得不每天乱找些词语来应付各种出现的场合，这确实是个悲剧。我是想改善这一状况，办法是，唔------清除------蚁群中出毛病的部分。这种手术有时相当复杂，要想做这种手术需要预先几年的研究功夫。

\dialogname{阿基里斯：}但是------那些患者在患失语症之前是具备语言能力的吗？

\dialogname{食蚁兽：}对。

\dialogname{阿基里斯：}可蚁群本来是不具备这种能力的，所以我有点糊涂。

\dialogname{螃蟹：}上周你没在这儿真是太糟糕了，阿基，那次食大夫和马姨到我这里来作客，我真应该把你邀来。

\dialogname{阿基里斯：}马姨是你的亲姨吗，老蟹？

\dialogname{螃蟹：}哦，不，其实她谁的姨也不是。

\dialogname{食蚁兽：}但是那可怜虫一定要每个人都叫她马姨，即使陌生人也得这样。这是她那些可爱的怪癖中的一个。

\dialogname{螃蟹：}是这样，马姨是个很怪的人，不过她可真是个活宝。上周我没有邀你过来见见她，真是太可惜了。

\dialogname{食蚁兽：}她确实是我有幸结识的蚁群中受教育最好的一个。我们俩一连好几个晚上就十分广泛的问题进行了长时间的交谈。

\dialogname{阿基里斯：}我过去一直以为食蚁兽是蚂蚁的克星，而不是蚁中知识分子的保护神。

\dialogname{食蚁兽：}啊，当然，这两者是很不一致的。我跟蚁群保持着非常良好的关系。我吃的是蚂蚁，而不是蚁群------这样做对于我和蚁群双方都是不勉为其难的。

\dialogname{阿基里斯：}怎么可能------

\dialogname{乌龟：}怎么可能------

\dialogname{阿基里斯：}------吃掉蚁群中的蚂蚁却又造福于这个蚁群？

\dialogname{螃蟹：}怎么可能------

\dialogname{乌龟：}------火烧森林却又造福于这片森林？

\dialogname{食蚁兽：}怎么可能------

\dialogname{螃蟹：}------砍掉树枝却又造福于这棵大树？

\dialogname{食蚁兽：}------给阿基里斯剃头却又造福于他？

\dialogname{乌龟：}也许你们大家讨论得太专心了，巴赫这首赋格曲里刚出现的那个漂亮的提前进入你们谁都没有注意到。

\dialogname{阿基里斯：}什么是提前进入？

\dialogname{乌龟：}哦，抱歉得很，我还以为你知道这个词呢。这个术语的意思是同一个主题在不同的声部中一次又一次地重复插入，这些插入彼此间的时间间隔非常短。

\dialogname{阿基里斯：}要是我听的赋格足够多的话，我很快就会弄懂这一切，无需别人的提示，我自己就能把它们指出来。

\dialogname{乌龟：}请原谅，朋友们。很抱歉把你们的讨论打断了。刚才食大夫正要说明吃蚂蚁是怎样同作蚁群的朋友这件事完美无缺地协调起来的。

\dialogname{阿基里斯：}嗯，我现在模模糊糊地有点明白了，有规律地吞吃些数量有限的蚂蚁会改善整个蚁群的健康状况，这事也许是可能的------可真正叫人糊涂的是跟蚁群交谈这种鬼话。这是不可能的。蚁群只是一伙单个儿的蚂蚁，只会没头没脑地跑来跑去，忙着觅食做窝什么的。

\dialogname{食蚁兽：}如果你固执地只见树木不见森林，那你可以这么说，阿基。实际上，蚁群被当作一个整体来看时，是些良定义的单位，具有自己的特点，有时这种特点还包括对语言的掌握。

\dialogname{阿基里斯：}我觉得很难想象我一个人在森林中间大喊几声，随后就能从蚁群那儿听到什么回答。

\dialogname{食蚁兽：}不开窍儿的家伙！蚁群可不是这样运用语言的。它们不是大声说出来，而是写下来。你知不知道众多的蚂蚁是如何排成一串东游西串的？

\dialogname{阿基里斯：}哦，知道------常常是直冲着我那应有尽有的厨房杀奔过去，一头钻进我的桃子酱里。

\dialogname{食蚁兽：}其实，这种串有的就用编码的形式包含着信息。要是你懂得那个系统，你就可以像读一本书一样读懂它们说了些什么。

\dialogname{阿基里斯：}了不起。你能把你的意思再反传给它们吗？

\dialogname{食蚁兽：}没问题。我跟马姨就是用这种方式交谈了数小时之久的。我拿着一根棍往湿地上画道道儿，然后看着蚂蚁们弄明白我的这些道道儿。不一会儿，有的地方开始组合成新的列串。我非常喜欢看这些列串的形成过程。在它们组合时，我就预测它们下一步会怎么办（不过我猜错的时候比猜对的时候要多）。当串的排列完成时，我就知道了马姨在想些什么，接着就该由我作出回答了。

\dialogname{阿基里斯：}这个蚁群中一定有一些聪明过人的蚂蚁，我敢肯定。

\dialogname{食蚁兽：}我看你依然没有弄明白这里面的层次区别。就像你不会把一棵树错当成一片森林一样，你在这里也不应该把一只蚂蚁当作整个蚁群。你明白吗，马姨中的所有蚂蚁要多笨就有多笨，打死它们也不会说话！

\dialogname{阿基里斯：}那么，进行交流的能力从何而来？一定是存在蚁群里的某个地方！我不明白既然马姨能一连几个小时跟你大谈特谈，妙语连珠，那些蚂蚁怎么会全都蠢笨不堪呢？

\dialogname{乌龟：}在我看来，这种情况跟由神经元组成的人脑的构成没什么不同。当然没有人会认为单个的脑细胞本身是有智能的，并用它来解释人为什么能进行有智慧的谈话。

\dialogname{阿基里斯：}哦，当然不会。关于脑细胞的问题，我完全明白你的意思。只是\ldots\ldots 风马蚁不相及呀。我是说，蚂蚁到处逛来逛去，完全是随意的，无非是时不时地弄点食物的残渣\ldots\ldots 它们想干什么就干什么，在这种随心所欲里面，我一点也看不出它们的行为作为一个整体，能体现出什么有条有理的东西来------特别是那种为交谈所必需的大脑行为中所具有的条理。

\dialogname{螃蟹：}在我看来，蚂蚁的这种随心所欲是局限在一定的范围之内的。比方说，它们可以随意游逛，互相擦拂，捡起小东西，排列成串，等等。但是它们从没有从这个小世界中，这个它们所在的蚂蚁系统中，跨出过一步。它们永远也不会这么做，因为它们没有作这种想象的智力。所以蚂蚁是种非常可靠的工具，在这种意义上说，你可以依靠它们用某种方法来从事某项任务。

\dialogname{阿基里斯：}不过，即使如此，它们在其限度之内也还是自由的，它们只是随随便便地转来转去，毫无条理，一点也不考虑作为更高一层存在的思想机能，可食大夫却宣称这些蚂蚁都只是这种更高一层存在的组成部分。

\dialogname{食蚁兽：}嗯，不过还有一点你没能看出来，阿基------统计学中的合规律性。

\dialogname{阿基里斯：}这是什么？

\dialogname{食蚁兽：}比方说，作为个体来看，即使每个蚂蚁的活动路线似乎是很随意的，然而就大量的蚂蚁来说，从这种乱糟糟的状态中，还是能看出存在着一些总的趋向。

\dialogname{阿基里斯：}哦，你的意思我明白了。事实上，蚂蚁组串是这种现象的一个极好的例子。一方面，任何一只单个蚂蚁的活动你都是无法预知的，而另一方面，串本身似乎仍然是界说良好的、稳定的。这当然意味着单个蚂蚁并不是完全随意地跑来跑去。

\dialogname{食蚁兽：}完全正确，阿基。蚂蚁之间存在着某种程度的交流，这种信息传递恰好足以使它们避免完全随意的活动。靠着这种最低限度的交流。蚂蚁们可以互相提醒，以使它们意识到它们不是孤单的，而是正在同队友进行合作。这样就能调集大量的蚂蚁，用这种方式互相支援，去从事随便什么活动------例如组串------只是所需时间不限。我对脑外科手术模模糊糊的认识使得我坚信神经元的发射情形也是这样。一大堆神经元发射脉冲，以使另一个神经元发射，不是这样吗，老蟹？

\dialogname{螃蟹：}肯定是这样。以阿基里斯大脑中的神经元为例吧。每个神经元都从与它的输入线连在一起的那些神经元接收信号，如果输入的总和在某一时刻超过了临界点，这个神经元就会发射，并把它自己的输出脉冲传递到别的神经元，而这后一神经元也将发射------并沿着这条线一直传送下去。神经讯号沿着阿基里斯的路径毫不延迟地迅猛掠过，那架势比飞燕俯冲捕食蚊蚋还有过之无不及。每段曲折、每个回转，都由阿基里斯大脑中的神经结构预先排定，直至输入感觉的讯号产生了影响。

\dialogname{阿基里斯：}在正常情况下，我觉得我想些什么是受我自己控制的------可按你这种说法就整个儿颠倒了，你的意思好像是说``我''只是所有这些神经结构、自然规律控制的某种有机体的副产品，往坏里说则是由我那被歪曲了的感觉所造成的某种人为概念。换句话说，你让我觉得我并不知道我是谁------或者是什么------如果我还是个什么的话。

\dialogname{乌龟：}随着我们讨论的继续，你会越来越明白的。但是，食大夫------你用这种类比来说明什么呢？

\dialogname{食蚁兽：}我以前就注意到，在这两个彼此非常不同的系统中存在着某种相一致的东西。现在，我对这一点的认识更进一步了。事情似乎是------以排列成串为例来说------只有当蚂蚁的数量达到某一临界数量时，有条理的蚁群现象才会出现。如果出现了某种指向某一目标的活动，这也许是某一地带少数蚂蚁的行为，那么就会发生下面两种情况中的一种：或者是在一个短暂而热闹的开头之后夭折------

\dialogname{阿基里斯：}因为没有足够多的蚂蚁把这一活动进行下去？

\dialogname{食蚁兽：}一点不错。另一种情况是蚂蚁的数量达到了临界指标，那样一来那项活动就会像滚雪球一样，把越来越多的蚂蚁裹挟进来。在这后一种情况下，一个为同一目标而工作的完整的``蚁队''就出现了。这一目标可能是要排成串，或是为了收集食物，或许还可能是为了造窝。尽管这种结构规模很小，又十分简单，可它会在更大的规模上产生十分复杂的后果。

\dialogname{阿基里斯：}我现在已经能从你的描述里掌握``紊乱中的有序''这个观念了，可这离具有交谈能力还差得远呢。气体分子的随意碰撞同样也是混乱中的有序------这种混乱有三种参数可以对它进行描述：这就是体积、压强和温度。这同具有理解世界或谈论世界的能力的要求，还差得远呢！

\dialogname{食蚁兽：}这显示出在对蚁群行为的解释与对某一容器中气体行为的解释之间，存在着某种非常有趣的区别。解释气体的行为只要计算一下其分子运动的统计规律就行了。除了考虑这种气体本身以外，无需讨论别的什么比分子更高一层的结构因素。而在另一种情况下，对于蚁群，除非你弄清了结构的各个层次，否则你便无法理解------哪怕一点儿------这个蚁群的活动。

\dialogname{阿基里斯：}你的意思我明白了。在气体的研究里，你可以从最低的层次------分子------一下子跳到最高的层次------气体本身。不存在结构上的中间层次。那么，蚁群中是怎样产生作为中间层次的有条理活动的呢？

\dialogname{食蚁兽：}这跟蚁群内存在着好几种不同种类别的蚂蚁有关。

\dialogname{阿基里斯：}哦，对。我想我听说过这一点。这种现象被称作``种姓制''，对吗？

\dialogname{食蚁兽：}完全正确。除了蚁王，还有雄蚁，事实上，养家糊口的事它们是不管的，此外还有------

\dialogname{阿基里斯：}当然还有战士------光荣的反共战士！

\dialogname{螃蟹：}嗯\ldots\ldots 这恐怕不对，阿基。蚁群的内部是很共产主义化的，那它们的战士干嘛要反共？我说的对吗，食大夫？

\dialogname{食蚁兽：}是的，你对蚁群的说法是对的，老蟹，它们确实建立在某种共产主义的原则上。至于那些战士，阿基想得有点天真。其实，这些所谓``战士''几乎根本不会打仗。它们是些长着大脑袋的行动迟缓、动作笨拙的家伙，它们长着一副利颌，用来啃咬东西，不过它们可谈不上什么``光荣的''。正像在真正的共产主义社会里那样，工蚁才被称作是光荣的。大多数的家务都是由它们来做的，比如觅食、巡猎、带幼虫等。大多数的战争也是由它们去打的。

\dialogname{阿基里斯：}嗬。这事真荒唐。战士不打仗！

\dialogname{食蚁兽：}嗯，我刚说过的，它们根本不是什么战士。工蚁才是真正的战士，而那些所谓的战士------其实应该叫兵蚁------都是些又懒又呆的家伙。

\dialogname{阿基里斯：}哦，太可耻了！哼，我要是只蚂蚁，一定给它们规定些纪律！让那些呆脑壳开开窍儿！

\dialogname{乌龟：}你要是只蚂蚁！你怎么会是只蚂蚁！根本没法把你的大脑映射成蚂蚁的大脑，我看在这个问题上绞尽脑汁是不会有什么结果的。把你的大脑映射到蚁群中才是个更合理的命题\ldots\ldots 不过我们先别把话题岔开。让食大夫继续对蚂蚁组织中更高层次的种姓制和它们的角色做精彩的描述吧。

\dialogname{食蚁兽：}好吧。一个蚁群中有各式各样的任务需要完成，而单个的蚂蚁则发展了各种专长。蚂蚁的专长通常随蚂蚁年龄的变化而变化。当然，这也依赖于它们所在的种姓。任何一时刻的任何一个蚁群里，所有类别的蚂蚁都同时存在。在某些地方也许某种种姓的蚂蚁很稀少，而在别的地方也许很稠密。

\dialogname{螃蟹：}某一特定种姓或专长的蚂蚁，其稠密度是随意的吗？或者换句话说，在某一地区某一类的蚂蚁分布得很集中，在另一地区不那么集中，是有什么原因的吗？

\dialogname{食蚁兽：}你这个问题提得很好，这对理解蚁群是如何思想的有着至关重要的意义。实际上，一个蚁群内各种姓的分布是很严格的，是长期进化的结果。正是这种分布规律使蚁群达到了能同我进行交谈的复杂程度。

\dialogname{阿基里斯：}在我看来，蚂蚁们来来去去的不断活动使分布的严格性完全不可能了。蚂蚁们那种随意的活动会很快地破坏掉分布的严格性，正像某种气体中分子的固定格局由于各方面分子的随意碰撞，一刻也不能存在一样。

\dialogname{食蚁兽：}在蚁群里，情况正相反。事实上，恰恰是蚁群中蚂蚁们这种来呀去呀永不停息的活动，使得种姓分布能适应不断变化的情况，因而也就保持了合适的种姓分布。你明白吗，种姓分布并不是始终如一、严格不变的，而是必须要不断变化，从而以某种方式反映出蚁群所处的现实世界的环境。正是蚁群里的这种活动使种姓分布不断得到更新，使蚁群能同所面临的现时环境相适应。

\dialogname{乌龟：}能举个例子吗？

\dialogname{食蚁兽：}当然可以。当我，一只食蚁兽，来拜访马姨时，那些呆头呆脑的蚂蚁因为闻到了我的气味，全都吓得张皇失措------当然，这就是说，它们全都从我来之前它们所在的地方四散奔逃了。

\dialogname{阿基里斯：}不过这可以理解，你是蚁群的凶敌嘛。

\dialogname{食蚁兽：}哦，不。我必须重申，我根本不是蚁群的敌人，我是马姨最对劲儿的伙伴。而马姨则是最对我口味儿的姨。说蚁群中所有单个的蚂蚁都很怕我，这我同意------可这完全是另一码事。无论如何，你能理解了吧：由于我的到来所引起的蚂蚁的活动，完全改变了蚂蚁的内部分布。

\dialogname{阿基里斯：}很明白了。

\dialogname{食蚁兽：}这就是我刚才说过的更新。新的分布反映出了我的存在。这种从过去的状态向新状态的改变，可以被说成是使蚁群增添了``一条知识''。

\dialogname{阿基里斯：}你怎么能把一个蚁群内不同类型蚂蚁的分布说成是``一条知识''？

\dialogname{食蚁兽：}这里面有个至关重要的问题，需要详细解释。你知道吗，接下来的问题是你打算如何描述种姓分布。如果你继续用低层次的单位------单个的蚂蚁------来思考问题，那你就会只见树木，不见森林。这种层次太微观了，当你从微观的角度进行思考时，你必然会漏掉更大规模上的现象。你不得不用更恰当的高层次上的构架来描述种姓分布------只有这样，种姓分布如何能编码出许多条知识这一问题才有意义。

\dialogname{阿基里斯：}那么，那种用来描述蚁群现存状态的恰当单位，你是怎么找到的？

\dialogname{食蚁兽：}好吧，让我们从最基本的说起。当蚂蚁想做成某件事时，它们就组成各种小``蚁队''，也就是凑在一起去干一件事。就像我前面提到过的那样，蚂蚁们的小团伙总是在不断地形成和解散。那些实际上能存在一段时间的就是各种蚁队，它们没有散伙的原因就是因为有事情需要它们去做。

\dialogname{阿基里斯：}可前面你说过如果一群蚂蚁的规模超过了某一临界点，这个群体就会集聚在一起。现在你又说如果有什么事可做，这个群体就会集聚在一起。

\dialogname{食蚁兽：}这两种说法意义是一样的。比方说，在觅食时，如果某个闲逛的蚂蚁在某个地方发现了一些食物，并想把这个喜讯传给别的蚂蚁们，那么，作出反应的蚂蚁数量将会同食物样品的大小相一致------太少的数量将不会吸引足以超过临界点数量的蚂蚁。这也正是我所谓的无事可做的意思------食物少得不值得去重视。

\dialogname{阿基里斯：}我明白了。我应该假定这些``蚁队''是结构中的一个层次，处在单个蚂蚁的层次和蚁群的层次之间。

\dialogname{食蚁兽：}正是这样。有一种类型特别的蚁队，我把它称作``信号''------结构中所有的更高级的层次都建立在信号上。事实上，一切比信号更高级的实体都只是由行动协调一致的信号所组成的集合。高层次上的队其组成单位不是蚂蚁，而是些低层次上的队。这样一直找下去，最后就到了最低层次的队------也就是所谓的信号------在这之下，才是蚂蚁。

\dialogname{阿基里斯：}它们为什么叫``信号''这样一个容易引起联想的名字？

\dialogname{食蚁兽：}这名字是指它们的功能。信号的作用是把具有各种专长的蚂蚁输送到蚁群中需要的部分。因此信号的典型活动是这样的：信号的产生是由于蚂蚁数量超过了维持一个信号的存在所必需的临界点，随后这个信号便通过蚁群向远处传递，在某些时候，信号便多多少少分解还原成为其单个的组成部分，并丢下它们不管。

\dialogname{阿基里斯：}听起来有点像波浪，从远处挟裹来一些海星、海草，把它们孤零零地拋到了岸上，丢下它们不管。

\dialogname{食蚁兽：}在某种意义上这种类比是对的，蚁队确实抛下了一些它从远处带来的东西，可海浪还是要回到海里的，而在信号这里，就没有类似的载体，因为信号是由蚂蚁组成的。

\dialogname{乌龟：}我觉得一个信号只是在途经该蚁群中原来需要这类蚂蚁的地方时，才会分解掉。

\dialogname{食蚁兽：}那自然。

\dialogname{阿基里斯：}那自然？在我看来，某个信号总能到达原来需要它的那个地方这种事，并非那么显而易见。即使它传播的方向正确，它怎么能判断出它应该在哪儿解散？它怎么知道它已经到了该去的地方了？

\dialogname{食蚁兽：}这是个极端重要的问题，因为这涉及到对有目的的行为------或者说是似乎有目的的行为------从信号角度所作的解释。从对它们的描述里，人们倾向于把信号的行为说成是旨在满足需要，从而把它称作``有目的''的。但你也可以从别的角度看这一问题。

\dialogname{阿基里斯：}哎，等等。要么这种行为是有目的的，要么不是。我不明白怎么能又是又不是。

\dialogname{食蚁兽：}我来给你解释一下我看问题的方式，看你是否同意。一个信号一旦形成，这个信号本身并不考虑它应该往哪个方向传播。在这里起决定作用的是种姓分布。是它决定着信号在蚁群中的运动，以及这个信号能持续多长时间、在哪里``解散''。

\dialogname{阿基里斯：}这么说一切都依赖于种姓分布，是吗？

\dialogname{食蚁兽：}正是这样。让我先来描述一下正在传播中的一个信号。当信号发出的时候，组成这一信号的蚂蚁们或是通过直接的接触，或是依靠气味的传达，同位置上与它们相邻的蚂蚁进行合作，以使信号通过。这种接触与气味报道了某个地方出现了需求，诸如筑窝、哺育等这类事情。只要这种局部的要求还得不到满足，这个信号就会一直保持下去。但是一旦出现一批蚂蚁能够满足那一需求，信号便会解散，而奔赴来救急的蚁队就在现场全面铺开，进行工作。现在，你明白蚁群中的种姓分布如何发挥各队的指导作用是怎么回事了吗？

\dialogname{阿基里斯：}很明白了。

\dialogname{食蚁兽：}这种观察事物的方法并不要给这种信号加上个什么目的，这一点你也明白了吧？

\dialogname{阿基里斯：}是的。的确，我开始学会从两种互不相同的角度来看这一问题了。从蚂蚁的眼睛所看到的角度，信号是无目的的。信号中普通的蚂蚁只是在蚁群附近逛来逛去，没有什么特定的目的，直到发现了什么它觉得还可以的东西后，才停下来。它的队友们常常也同它意见一致，于是这时蚁队才算完成了任务，原地解散，有组织的集合就不复存在了，而只剩下一些散兵游勇。没有什么计划啦、事先的预测啦这一类东西，也不需要什么侦察来决定正确的方向。然而从蚁群的角度看，蚁队的行为却是对以种姓分布这种语言写成的信息所作的反应。所以从这一角度看，这种行为显得很像是有一定的目的。

\dialogname{螃蟹：}如果种姓分布完全是随意的，那会怎么样？那些信号还会结聚和解散吗？

\dialogname{食蚁兽：}当然啦。不过这个蚁群将不会保持多久，这是因为它的种姓分布是无意义的。

\dialogname{螃蟹：}这正是我想弄明白的。蚁群的存在是由于它的种姓分布有一定的意义，而这一意义是在其整体水平上产生的。在比整体水平低的层次上，是看不到这种意义的。你除非在这种高层次上进行解释，否则就无法读懂它的意思。

\dialogname{食蚁兽：}我明白你的观点，不过我看你把问题想得太狭隘了。

\dialogname{螃蟹：}怎么？

\dialogname{食蚁兽：}蚁群几十亿年来一直在经历着严酷的进化选择。只有少数一些机制被选择出来，而大部分都被淘汰了。结果就形成了我们现在正在讨论的支配着蚁群的这一套机制。如果你能在电影上观察这一完整的过程------当然，这比实际进程要快上差不多十亿倍------你就会觉得各种机制的出现完全是对外界压力的自然反应，就像开水冒泡是对外在的热源的自然反应一样。我想你不会觉得开水冒泡有什么``意义''和``目的''------不过，也许你会？

\dialogname{螃蟹：}不，但是------

\dialogname{食蚁兽：}这就是我的观点。无论泡沫有多大，其产生都可以追溯到分子水平，你可以不管什么``高层次的规律''。蚁群和蚁队的道理也是一样的。通过从进化这一开阔的角度考虑这一问题，你就可以排除有意义有目的的蚁群这一想法。这样它们就变得可有可无了。

\dialogname{阿基里斯：}哎，食大夫，你不是刚告诉我们你跟马姨交谈过吗？现在看来你要否认她能交谈或思想啦。

\dialogname{食蚁兽：}我没有自相矛盾，阿基。你知道，同别的人一样，用这么巨大的时间单位去观察事物，对我来说也是很困难的，所以我觉得改换一下角度会更方便些。当我撇开进化过程而对此时此地的现象进行观察的时候，就又回到了那种目的论式的语言：即种姓分布的意义和信号的有目的性。我不仅在考虑蚁群时采用这种角度，而且当我考虑我自己和别人的大脑时，也是如此。不过，如果必要的话，我也能采用别的角度，把意义排除于这些系统之外，只是这多少要费点劲儿。

\dialogname{螃蟹：}进化的确能创造奇迹，正像你没法儿从创造奇迹的魔术师那里预知他又要从袖子里抖落出什么东西一样，你也无法猜测进化下一次将会弄出什么玩艺儿来。比方说，要是出现这样的事，我一点儿也不会感到吃惊：如果在理论上可能，有两个或更多的``信号''彼此交错，每一个都不知道另一个也是一个信号------每一个都把另一个仅仅当作是一般状态的蚂蚁群。

\dialogname{食蚁兽：}这不仅是理论上可能的，事实中也是常有的！

\dialogname{阿基里斯：}哟\ldots\ldots 我脑子里冒出来一个什么怪念头啊。我在想象蚂蚁们朝四个不同的方向运动，有黑的，有灰的，互相交错，共同组成一个有序的图案，就像------就像------

\dialogname{乌龟：}你是说像一首赋格曲？

\dialogname{阿基里斯：}对------正是！一首蚂蚁赋格。

\dialogname{螃蟹：}很有趣的念头，阿基。不过，说到开水倒让我想起茶来了。谁想再要点儿？

\dialogname{阿基里斯：}我再要一杯，老蟹。

\dialogname{螃蟹：}行啊。

\dialogname{阿基里斯：}你认为可以把这样一首``蚂蚁赋格''分解成彼此不同的可见的``声部''，是吗？这可太难了，我是说，如果我要------

\dialogname{乌龟：}我不要，谢谢。

\dialogname{阿基里斯：}------在一首赋格曲中------

\dialogname{食蚁兽：}要是不太麻烦的话------

\pandocbounded{\includegraphics[keepaspectratio]{./media/OEBPS/Images/img_061.jpg}}

图61．``蚂蚁赋格''

艾舍尔作（木刻1953）

\dialogname{阿基里斯：}------跟踪其中某一个声部------

\dialogname{食蚁兽：}------我也想来点儿，老蟹。

\dialogname{阿基里斯：}------而所有那些------

\dialogname{螃蟹：}区区小事。四杯茶------

\dialogname{乌龟：}三杯！

\dialogname{阿基里斯：}------声部都在同时进行。

\dialogname{螃蟹：}------转眼就得！

\dialogname{食蚁兽：}这个想法很有意思，阿基。不过没有人能让人信服地画出这种画来。

\dialogname{阿基里斯：}太令人遗憾了。

\dialogname{乌龟：}这个问题也许你能回答，食大夫。一个信号，从它形成到解散，总是由同样一群蚂蚁组成吗？中间不会换吗？

\dialogname{食蚁兽：}实际上，如果附近有同样种姓的蚂蚁，信号中的蚂蚁经常脱离信号而由别的同种姓的蚂蚁来替代。大多数情况下，当信号解体的时候，组成这一信号的蚂蚁已经全都不是当初形成这一信号时的蚂蚁了。

\dialogname{螃蟹：}我搞清楚是怎么一回事了：蚁群中的信号不断影响种姓分布，这一活动是对蚁群内部的要求所作出的反应，而蚁群的这种内部要求则又是蚁群对其所处外部环境的反应。因此，正如您所说的那样，食大夫，种姓的分布不断地调节、变化，最终反映着外部世界。

\dialogname{阿基里斯：}可结构中的那些中间层次是怎么回事？你刚才一直在说种姓分布最好应该以由别的蚁队组成的蚁队为单位来描述，而不应以蚂蚁或信号为单位来描述，那些组成蚁队的蚁队又是由一些别的蚁队组成的，同理下推，直至组成蚁队的成员是单个儿的蚂蚁，而不再是蚁队。你还说，要想理解种姓分布何以能被描述为某种有关外部世界的编码了的信息，这个观点是关键性的。

\dialogname{食蚁兽：}正是，我们现在正要谈这个问题。我喜欢把某些层次较高的蚁队称作``符号''。你要注意，这个词在这里的意思与它通常的意义有所不同。我所说的``符号''是指一个复杂系统的主动的子系统，这种子系统是由更低层次的主动的子系统构成的\ldots\ldots 因此，它们与被动的符号很不一样，那些被动的符号在系统之外。汉字中的笔划或单个儿的音符，都属于这种符号，它们是死的，有待于某个主动的系统去处理它们。

\dialogname{阿基里斯：}噢，这可真够复杂的。我一点儿也没想到蚁群具有如此抽象的结构。

\dialogname{食蚁兽：}是的，很了不起。不过，结构中的所有这些层次对于各种各类知识的储存都是必需的，有了这种知识储存，才能使这个机制变成``有智能的''，``有智能的''的意思在这里可以是这个词的任何合理的涵意。任何能运用语言的系统都基本上具有与此相同的一套内部层次。

\dialogname{阿基里斯：}先慢着点。你是不是暗示说我的大脑实际上也是由一群窜来窜去的蚂蚁组成的？

\dialogname{食蚁兽：}哦，不。你不要从字面意义上理解我的话。最低的层次是全然不同的。比方说，食蚁兽的大脑就不是由蚂蚁组成的。但是你如果从大脑的最低层上溯一两个层次的话，你就会发现这些层次的机制在别的具有同等智力的系统中------比如蚁群中------有与它完全一致的对应物。

\dialogname{乌龟：}这就是把你的大脑映射到蚁群中去，而不是映射到蚂蚁的大脑中去的依据，阿基。

\dialogname{阿基里斯：}对你的这份夸奖我领情了。不过这种映射是怎样进行的？比方说，我大脑中的什么东西对应于你所谓的信号这种低层次的蚁队呢？

\dialogname{食蚁兽：}哦，我对大脑的研究只是种业余爱好，所以无法一一细述它的神异之处。可是------如有不妥之处，老蟹，请不吝赐教------我猜测，在大脑中对应于蚁群里信号的那种东西是兴奋的神经元。或者，也可能是某种更大范围上的活动，比如神经元兴奋的模式什么的。

\dialogname{螃蟹：}我没有什么异议。不过你是否觉得，找出确切的对应物这一问题，即使有希望得到解决，相对于我们这场讨论的目的来说，也并不是什么十分紧要的事。在我看来，最重要的是知道确实存在着这种对应，即使我们此时此刻还不完全清楚这种对应指的是哪一部分也没有关系。我只想对你提出的问题中的一点追问一下，食大夫，因为这一问题涉及究竟在哪个层次上，人们可以断定开始产生了这种对应关系。你似乎认为一个信号在大脑中有其直接的对应部分，然而我觉得只是在你所谓的``主动的符号''这一层次上，才可能存在那种对应关系。

\dialogname{食蚁兽：}你的解释比我的更确切，老蟹。多亏你提出这个微妙的问题。

\dialogname{阿基里斯：}有什么事情只能由符号做而不能由信号做？

\dialogname{食蚁兽：}这情形有点像字与笔划之间的区别。字作为具有意义的单位，是由笔划组成的，而笔划本身却不具有意义。这也正是符号与信号的区别所在。这实际上是个很有用的类比，条件是你不要忘了字和笔划是被动的，而符号与信号却是主动的。

\dialogname{阿基里斯：}我一定会记住的，不过我不敢说我理解了为什么强调主动与被动两种性质的区别具有如此重要的意义。

\dialogname{食蚁兽：}原因在于你加于被动符号上的意义，例如一页书中的某个字，实际上来源于你大脑中与之相对应的主动符号所产生的意义。因此，只有当一个被动符号的意义与主动符号的意义发生关系时，这个被动符号的意义才能被正确理解。

\dialogname{阿基里斯：}明白了。可是究竟是什么赋予一个符号------确切地说，是主动的符号------以意义的呢？信号虽然也是一种很出色的东西，可你却说它不具有意义。

\dialogname{食蚁兽：}这与符号可以使其它符号被触发的方式有关。当某一符号成为主动时，它并非是孤立无援的。它可以通过某种媒介------在蚂蚁这里就是种姓分布------传播开来。

\dialogname{螃蟹：}当然，大脑里是没有种姓分布这类玩意儿的，在大脑中，与种姓分布相对应的是``大脑状态''。你必须描述出所有神经元的状态、所有它们交互联系的状态，以及使每个神经元触发的临界点。

\dialogname{食蚁兽：}一点不错。现在让我们把``种姓分布''和``大脑状态''都一并归于一个共同的名称之下，这个共同的名称就是``状态''。这样，这种状态就可以在某种更低或更高的层次上描述出来了。对蚁群状态进行低层次上的描述将会非常麻烦，这需要确定每个蚂蚁的位置、年龄、种姓，以及其他一些诸如此类的事情。这种详细的描述实际上无助于统观为什么该蚁群处于这种状态。另一方面，在更高层次上进行的描述则需要逐一查清哪些符号可以被哪些其他符号的组合所触发，以及触发的条件等等。

\dialogname{阿基里斯：}在信号或蚁队层次上进行的描述会怎么样呢？

\dialogname{食蚁兽：}这种描述所在的层次位于较低的层次和符号层次之间。对于在蚁群中某一特定位置上正在发生的事，这种描述能提供相当多的信息，虽然这要比一个蚂蚁一个蚂蚁的描述所含的信息少些，这是因为蚁队是由一团团的蚂蚁组成的。那种一个蚁队一个蚁队的描述就好像是一个蚂蚁一个蚂蚁描述的提要。然而，在蚁队层次上的描述必须加进一些在蚂蚁层次上所没有的内容------例如各蚁队之间的相互关系，散于各处的各种种姓的补充力量。这种新添的复杂性是你进行这种概括所必须付出的代价。

\dialogname{阿基里斯：}我觉得对不同层次上的描述进行比较是件非常有意思的事。最高层次上的描述似乎最具有解释力，在这一层描述中，你能得到关于蚁群的最直观的印象，然而，说来奇怪的是，这种描述却撇开了似乎是最重要的方面------蚂蚁。

\dialogname{食蚁兽：}不过你要明白，尽管你所看到的都不过是蚂蚁，可它们却全然不是最重要的方面。我们承认，如果没有这些蚂蚁，蚁群就不会存在。可是，与之相应的某种东西都是存在的------比如大脑------完全不依赖蚂蚁的存在而存在。因此，至少从一种高层次的观点上看，蚂蚁是可有可无的。

\dialogname{阿基里斯：}我敢说没有任何蚂蚁会热烈拥护你的这种理论。

\dialogname{食蚁兽：}不过，我从没有见过什么蚂蚁具有高层次的观点。

\dialogname{螃蟹：}你的这番描述太违背人们的直观了，食大夫。如果你所言不差，事情似乎就会是：为了了解整个儿结构，你必须省略掉构成这一结构的任何最基本的组成材料。

\dialogname{食蚁兽：}也许我打一个比方你们就能更清楚些了。请想象一下你面前摆着一本查尔斯·狄更斯的小说。

\dialogname{阿基里斯：}《匹克威克外传》------可以吗？

\dialogname{食蚁兽：}再好没有了！现在我们就来做这样一个游戏：你先想办法使每种笔划都具有意义，这样当你一个笔划一个笔划地读这本书时，整部《匹克威克外传》就都具有意义了。

\dialogname{阿基里斯：}嗯\ldots\ldots 你是说我每见到一个字，就比如``人''字吧，就要把它看作是两个各具某种一定意义的概念，并且是一个接一个地紧挨着，每种相同的笔划的意义总是一定的，没有任何变化，是吗？

\dialogname{食蚁兽：}正是这样。这样一来就只有``丿''概念、``\gebfont{㇏}''概念了------而且你每次碰到它们时，它们的意思都同前面它们出现时的意思一样。

\dialogname{阿基里斯：}嗬，这么一来，``阅读''《匹克威克外传》的活动简直就变成了一场难以言状的噩梦了。不管我赋予每个笔划什么意义，这一切都只不过是一场徒劳无益的玩弄无意义的游戏。

\dialogname{食蚁兽：}没错儿。单个儿的笔划与现实世界之间不存在什么自然的映射关系。只有在更高的层次上才会有这种自然的映射------即在字词与现实世界的各部分之间。你如果要叙述这本书的内容，你就不要涉及它的笔划层次。

\dialogname{阿基里斯：}当然不会！我只会叙述情节、人物这一类的东西。

\dialogname{食蚁兽：}这你就对了。你应该完全略去这本书的构成材料，即使这本书有赖于它们。它们是媒介，而不是信息。

\dialogname{阿基里斯：}明白了------那么，蚁群呢？

\dialogname{食蚁兽：}在蚁群这里，主动的信号代替了被动的笔划，主动的符号代替了被动的字词------不过原理是一样的。

\dialogname{阿基里斯：}你是说在信号与现实世界之间无法建立起映射关系？

\dialogname{食蚁兽：}你会发现你无法通过触发某个新的信号来使它具有什么意义。在比它更低的层次上------如蚂蚁层次------也同样办不到。只有在符号层次上，触发模式才有意义。比方说，你想象一下在我来拜访你时，你正在观察马姨。你无论观察得多仔细都行，不过，除了蚂蚁们的集聚、解散之外，你仍然什么也看不到。

\dialogname{阿基里斯：}肯定是这样。

\dialogname{食蚁兽：}而我在观察的时候，由于研究了它的高级层次，而不是低级层次，就会看到有好几种休眠的符号此时被唤醒了，并转译成下面的思想：``哦，食大夫这宝贝儿又来了------这真叫人高兴！''------或者是其他一些大意如此的话。

\dialogname{阿基里斯：}听起来正像我们四个刚才看``无之图''时那样，我们各自看到了不同的层次------或者至少说是我们三个\ldots\ldots{}

\dialogname{乌龟：}我在《平均律钢琴曲集》中翻到的那幅古怪的图跟我们现在讨论的问题之间有某种相类似的地方，这可真是个奇怪的巧合。

\dialogname{阿基里斯：}你认为这只是巧合吗？

\dialogname{乌龟：}当然。

\dialogname{食蚁兽：}哎，我估计你们现在能够理解马姨的那些从对符号的操作中产生的想法了。这些符号是由信号组成的，信号是由蚁队组成的，蚁队是由低层次的蚁队组成的，如此下推，直至单个儿蚂蚁。

\dialogname{阿基里斯：}你为什么把它称作``对符号的操作''？如果符号都是主动的，那么是由谁来操作？这个发号施令的是谁？

\dialogname{食蚁兽：}这得回到你前面提到过的有关目的的问题上来。说符号自身是主动的，这一点不错，可这些符号所进行的活动却不是绝对自由的。所有符号的活动都严格地受制于它们所处系统的状态。因此，是整个系统决定着它里面的符号如何彼此相互触发，所以，把整个儿系统说成是``发号施令的''是合情合理的。当符号被触发时，它所在系统的状态便慢慢发生转变，以适应新的情况。不过，其中有些特征是不随时间的变化而变化的。正是这个部分恒定、部分变化的系统才是那个发号施令的。我们可以给整个儿系统起个名字。比方说，马姨就可以被看做是那个操作符号的。同样的道理也适用于你，阿基。

\dialogname{阿基里斯：}对于我到底是谁这个问题，你的这种看法真是挺奇怪的。我不敢说我已经完全理解了，但我会好好考虑它们的。

\dialogname{乌龟：}在你思考你大脑中的那些符号时，去注意你大脑中此时产生的符号会非常有趣的。

\dialogname{阿基里斯：}这对我来说太复杂了。理解在符号水平上观察并解释蚁群的可能性，已经叫我很头疼了。我完全可以想象得出在蚂蚁层次上观察蚁群是怎么一回事；我也可以想象在信号层次上的观察，只不过多少要费点儿劲儿；可天知道在符号层次上观察蚁群是怎么回事！

\dialogname{食蚁兽：}只要经过长时间的练习就能学会。不过，谁要是达到了我现在的阶段，谁就能很容易地释读蚁群的最高层次了，就像你能从``无之图''中读出``无''一样。

\dialogname{阿基里斯：}真的吗？这一定特别不同寻常。

\dialogname{食蚁兽：}在某种意义上说是这样------不过，另一方面上讲，对你来说这并不陌生。

\dialogname{阿基里斯：}对我来说不陌生？你指什么？除了在蚂蚁层次上以外，我没有用任何其它方式观察过蚂蚁。

\dialogname{食蚁兽：}也许你没有。不过蚁群在很多方面同大脑没有什么不同。

\dialogname{阿基里斯：}可我也从没有看过或读过任何大脑呀。

\dialogname{食蚁兽：}你自己的大脑呢？你意识不到你自己的思想吗？这难道不是意识的本质吗？当你意识到你自己的思想时，你难道不是直接在符号层次上读你的大脑吗？

\dialogname{阿基里斯：}我从没有这么想过。你是说我绕过了所有低级层次，只看到最高层？

\dialogname{食蚁兽：}意识系统正是这样。它们只是在符号层次上感觉到自己，对更低的层次，如信号层次，则毫无意识。

\dialogname{阿基里斯：}这是不是说，在大脑中，主动的符号总是在调节自身，以反映出大脑的整个儿状态，这一切总是处在符号水平上？

\dialogname{食蚁兽：}正是这样。在一切有意识的系统中，符号表现着大脑的状态，而这些符号本身却正是它们所代表的大脑状态的一个部分。因为意识活动要求极大程度上的自我意识。

\dialogname{阿基里斯：}这可够古怪的。这就是说虽然我的大脑中每时每刻都在发生忙忙碌碌的活动，可我却只能以一种方式感受这种活动------就是说是在符号层次上。对于更低的层次，我完全没有感觉。这就像是在读一本狄更斯的小说时只通过直接的视知觉，而完全看不到一个个的笔划。我无法想象这么古怪的事竟能真的发生。

\dialogname{螃蟹：}可这种事在你看那幅``无之图''时确实发生了，你只读出了``无''而没有看出更低层次上的``整体论''与``简化论''。

\dialogname{阿基里斯：}你说的对------我忽略了低层次，只看到了最高层。我不知道我是否也同样忽略了我大脑中更低层次的意义，而只是看到了符号层。糟糕的是最高层并没有包含底层的所有信息，因此，只知道最高层并不能因而知道底层在说些什么。不过我觉得，希望最高层能把底层的一切信息都编成码，这是一种很幼稚的想法------底层信息可能无法渗透上来。``无之图''是可能有的最明白的图解：最高层只有``无''，它与低层没有任何关系！

\dialogname{螃蟹：}绝对正确。\kaitismall{（拿起``无之图''，凑近了研究。）}嗯\ldots\ldots 图中最小的那些笔画有点奇怪。它们歪歪扭扭的\ldots\ldots{}

\dialogname{食蚁兽：}让我看看。\kaitismall{（凑近了看着``无之图''。）}我看还有一层，我们都给漏掉了！

\dialogname{乌龟：}你这话可不代表别人，食大夫。

\dialogname{阿基里斯：}啊------不可能！让我瞧瞧。\kaitismall{（很仔细地看。）}我知道你们大家肯定不会相信：这张图的秘密其实近在眼前，就藏在深层里面。全部秘密就是一个字，重复了一遍又一遍------但是个无比重要的字：``无''！谁想得到！跟最高层一样！我们大家谁也没能哪怕是猜到一丁点儿！

\dialogname{螃蟹：}要不是你，阿基，我们谁也不会注意到。

\dialogname{食蚁兽：}我不知道最高层与最低层之间的这种一致是否是出于巧合，还是某个创作者有意为之？

\dialogname{螃蟹：}谁能说得准呢？

\dialogname{乌龟：}因为我们一点儿也不知道这幅奇怪的图为什么是在螃蟹版的《平均律钢琴曲集》中，所以我看我们没法儿回答这个问题。

\dialogname{食蚁兽：}我们的讨论一直很活跃，可我还是愿意继续洗耳恭听这首又长又复杂的四声部赋格。真是太美了。

\dialogname{乌龟：}确实很美。再过一会儿，就要有持续音了。

\dialogname{阿基里斯：}持续音是指一段音乐逐渐慢下来，并时不时地停止在某一音符或和弦上，然后，在一段短暂的静默之后，又重新恢复正常速度，是吗？

\dialogname{乌龟：}不对，你这是在说``延长记号''------一种音乐上的分号。不知你注意没有，刚才那首前奏曲中就有这么一段？

\dialogname{阿基里斯：}我想我是错过了。

\dialogname{乌龟：}不过，你马上还会有机会听到一个延长记号的------事实上，在这首赋格临近结尾的地方，还会有这么两段的。

\dialogname{阿基里斯：}哦，太好了。你过会儿会把它们指出来，是吧？

\dialogname{乌龟：}只要你愿意。

\dialogname{阿基里斯：}那么，请告诉我，什么是持续音？

\dialogname{乌龟：}持续音就是在复调音乐中，所有声部中的一个声部（通常是最低的那个声部）停留在某一个音符上，而与此同时，其他的声部继续向前发展，不受影响。在这支赋格中，这个持续音停留在音符G上。仔细听，你会听到的。

\dialogname{食蚁兽：}有一天我去拜访马姨，碰到了一件事，让我想起你的那个建议。你建议当阿基里斯的大脑中的那些符号创造有关它们自身的思想时，应该去留心那些符号本身。

\dialogname{螃蟹：}快告诉我们你碰到什么事了。

\dialogname{食蚁兽：}马姨那天感到很孤独，很想同别人聊聊。因此她很大方地要我自己挑那些我所能见到的最鲜嫩的蚂蚁。\kaitismall{（她从不吝啬她的那些蚂蚁。）}

\dialogname{阿基里斯：}妙极------

\dialogname{食蚁兽：}这时我正在观察那些产生她的思想的符号，因为在那些符号中间有很多看去极其鲜嫩的蚂蚁。

\dialogname{阿基里斯：}妙极------

\dialogname{食蚁兽：}这样我就挑了些最肥嫩的，这些蚂蚁是我正在解释的高层次符号中的一部分，它们所在的那个符号是在说明这样一个意思：``自己随便挑那些能引起你食欲的蚂蚁吧。''

\dialogname{阿基里斯：}妙极------

\dialogname{食蚁兽：}这对它们来说很不幸，可对我来说运气极了，这些小虫子们一点也不懂它们作为一个集体在符号层次上对我说的话。

\dialogname{阿基里斯：}妙极------了！这真是个大怪圈。它们完全不知道自己所参加的活动。它们的行为可以被看作是更高层次上某种模式的一部分，但是它们肯定对此一无所知。嗨，真可怜------事实上也是绝妙的讽刺------它们对此毫无感觉。

\dialogname{螃蟹：}你说的对，龟兄------这是个妙极了的持续音。

\dialogname{食蚁兽：}我以前从没听到过，不过它太明显了，没有人会感觉不到的。给人的印象太深了。

\dialogname{阿基里斯：}什么？持续音已经出现过了吗？它要是这么明显，我怎么会没注意到？

\dialogname{乌龟：}也许你太陶醉于你说话的内容了，以致于你对此一无所知。嗨，真可怜------事实上也是绝妙的讽刺------你对此毫无感觉。

\dialogname{螃蟹：}告诉我，马姨是住在一个蚁穴中吗？

\dialogname{食蚁兽：}哦，她有一份很大的家产。那里曾属于别人，留下了一个令人伤感的故事。不管怎么说，她的庄园是非常大的。比起许多别的蚁群来，她住的算是很阔气了。

\dialogname{阿基里斯：}这与你前面跟我们讲的蚁群的共产主义性质怎能一致起来呢？在我看来，共产主义与住在一所豪华的大宅里是互相矛盾的！

\dialogname{食蚁兽：}共产主义是蚂蚁层次上的。在一个蚁群里，所有的蚂蚁都为共同的利益而工作，即使有时这对单个儿的蚂蚁有害。这些只是马姨结构的内在方面，就我所知，她甚至可能感觉不到这种内部的共产主义。人类大都对他们自己的神经元一无所知。事实上，作为有点神经质的生物，他们也非常满足于对自己的大脑一无所知这种状况。马姨也多少有些神经质。只要她一开始想她的那些蚂蚁，她就会马上变成热锅上的蚂蚁。所以，只要可能的话，她就竭力避免想它们。我实在怀疑她是否了解她内在结构中的共产主义社会。她自己是个坚定的自由主义信徒------你知道，她是完全自由放任的。所以，至少对我来说，她愿意住在一所豪华的庄园里是非常可以理解的。

\pandocbounded{\includegraphics[keepaspectratio]{./media/OEBPS/Images/img_062.png}}

图62．两个伟大名字的交织

刘皓明绘

\dialogname{乌龟：}在我浏览《平均律钢琴曲集》的这个可爱版本的时候，现在刚好翻过这页，我发现那两个延长记号中的一个马上就要出现了------你好好听听吧，阿基。

\dialogname{阿基里斯：}我会的，我会的。

\dialogname{乌龟：}对着这一页的地方还有一张古怪透顶的图。

\dialogname{螃蟹：}难道又有一张？

\dialogname{乌龟：}自己找吧。\kaitismall{（把总谱递给螃蟹。）}

\dialogname{螃蟹：}啊哈！是四组小字。让我看看------分别是``彼''、``德''、``巴''、``赫''，各自出现了许多次。我纳闷，为什么中间的字大，两边的字小？

\dialogname{食蚁兽：}能让我看看吗？

\dialogname{螃蟹：}当然。

\dialogname{食蚁兽：}我看你是只见树木，不见森林了。实际上，这是四个大字：``约''、``塞''、``费''、``马''，中间的字小，两边的字大。阿基里斯，你看呢？

\dialogname{阿基里斯：}让我瞧瞧。嗯，我看这些字是越往右边越小。

\dialogname{乌龟：}它们能组成有意义的词吗？

\dialogname{阿基里斯：}嗯\ldots\ldots``约''、``塞''、``巴''、``赫''。哦，我明白了。这是巴赫的姓名！

\dialogname{乌龟：}你这种看法真怪。我看这些字是越往右边越大，而且\ldots\ldots 组成了\ldots\ldots 一个人\ldots\ldots 的\ldots\ldots 姓\ldots\ldots 名，\kaitismall{（他逐渐变慢，最后几个字拉得特别长。然后是一阵短暂的静默。突然，他又正常起来，好像什么都没有发生过似的。）}------``彼·德·费马''。

\dialogname{阿基里斯：}哦，我敢肯定你满脑子里都是数学家费马。你看什么都是费马最后定理。

\dialogname{食蚁兽：}你刚才说的对，龟兄------我在这首赋格里刚听到一个迷人的小段延长记号。

\dialogname{螃蟹：}我也听到了。

\dialogname{阿基里斯：}你们是说除了我人人都听到了？我开始觉得我有点笨了。

\dialogname{乌龟：}哎哎，阿基------别沮丧。我敢肯定你不会错过赋格的最后延长记号------马上就会出现的。不过，让我们回到刚才的题目上来吧，食大夫，你提到的那个有关马姨庄园的前所有者的伤感故事是怎么回事？

\dialogname{食蚁兽：}那位前所有者是位出类拔萃的人物，是有史以来最富于创造性的蚁群之一。他的名字叫蚁翰·塞巴斯蚁安·蜚蚂，一位专业数学家、业余音乐家。

\dialogname{阿基里斯：}他可真是多才多艺啊！

\dialogname{食蚁兽：}在他的创造力的颠峰期，他惨遭最不是时候的猝死。一天，那是夏天里非常炎热的一天，他外出想去避避暑，正好赶上一场疾风暴雨------是那种一百年左右才会有一次的大暴雨------真叫是从天而降，把蚁·塞·蜚蚂淋了一个落汤蚂蚁。因为这场暴雨突如其来没有任何预兆，蚂蚁们完全晕头转向了。几十年建立起来的如此完美的复杂机体瞬间便毁于一旦。真是个悲剧。

\dialogname{阿基里斯：}你是说所有的蚂蚁全都淹死了，这当然就是说可怜的蚁·塞·蜚蚂完蛋了？

\dialogname{食蚁兽：}实际上并没有。蚂蚁们打算重新振兴，它们中幸存的每只蚂蚁都爬到飘浮在激流洪涛上的各种树枝木棍上。可是在洪水消退、蚂蚁们重新回到它们巢穴附近的地面上时，整个组织已不复存在了。种姓分布已经彻底毁坏了，这些蚂蚁已无力重建曾经有过的那个协调得如此完美的组织。它们就像破碎的镜子一样不能指望重圆了。我就像所有的媒婆月老一样，企图把可怜的蜚蚂重新撮合到一块儿。我满怀信心地放了些糖和奶酪，徒劳地希望蜚蚂也许会重新合成\ldots\ldots{}\kaitismall{（掏出一块手帕擦他的眼睛。）}

\dialogname{阿基里斯：}你真英勇！我以前以不知道食蚁兽们竟有这么一副慈悲心肠。

\dialogname{食蚁兽：}不过这根本无济于事。他去了，不可能再重建。可是随后发生了一件非常奇怪的事：在接下来的几个月里，那些过去是蚁·塞·蜚蚂成员的蚂蚁们慢慢儿又重新组合起来，并且建立了一个新的组织。因此马姨就诞生了。

\dialogname{螃蟹：}真了不起！马姨就是由组成蜚蚂的那些蚂蚁构成的？

\dialogname{食蚁兽：}对，最初是这样的，没错儿。而现在，有些老蚂蚁已经死掉了，被替代了。但是有一些从蚁·塞·蜚蚂时代留下来的遗老还活着。

\dialogname{螃蟹：}你偶尔也会在马姨身上找到某些明显的蚁·塞·蜚蚂的旧特征吗？

\dialogname{食蚁兽：}一点也找不到。他们毫无共同之处。而且在我看来，也没有任何理由应该这样。重新组织许多部分以构成一个``总和''，毕竟有着许多种彼此不同的方式。而马姨只不过是旧的各部分的一种新``总和''而已。注意，并没有多于那个总和------只是这个特定种类的总和。

\dialogname{乌龟：}说到总和，叫我想起了数论。在数论里有时可以把一个定理拆卸成它的各构成符号，再以一种新的顺序重新组织它们，从而得出一个新定理。

\dialogname{食蚁兽：}我从未听说过这种现象。不过我承认在这一领域里我完全是个外行。

\dialogname{阿基里斯：}我也没听说过------而我对这一学科是颇为精通的，也许我自己不该这么说。我怀疑龟兄不过是在打哈哈。如今我可太了解他了。

\dialogname{食蚁兽：}说到数论，又叫我想起了蚁·塞·蜚蚂，因为数论是他拿手的领域之一。事实上，他对数论作出过某些十分了不起的贡献。而马姨则在任何哪怕跟数学沾一丁点儿边的事情上都极其迟钝。此外，她在音乐方面的趣味不过平平而已，而塞巴斯蚁安在音乐方面极有天才。

\dialogname{阿基里斯：}我非常喜欢数论。你能不能给我们讲一讲塞巴斯蚁安做出的那些贡献的本质？

\pandocbounded{\includegraphics[keepaspectratio]{./media/OEBPS/Images/img_063.jpg}}

图63．一座蚁桥的照片

徙途中，兵蚁有时会用它们自己的身体搭起一座活桥。在这幅照片中就是这样一座桥。可以看到一个蚁群的工蚁们腿脚相连，它们的跗节肢钩在一起，顺着桥头形成许多不规则的链系统。一只共生的蠹鱼正在过桥，此时它处在中心部位。{[}引自爱德华·威尔逊，《昆虫社会》{]}

\dialogname{食蚁兽：}好吧。\kaitismall{（停顿了一小会儿，呷了一口茶，然后才开始说起来。）}你们听说过狒犸那个臭名昭著的``平均验刚劲猜想''吗？

\dialogname{阿基里斯：}我说不好\ldots\ldots 听起来怪熟的，可我想不起来了。

\dialogname{食蚁兽：}内容非常简单。溺爱儿·的·狒犸，这位专业数学家、业余律师，在读一本丢饭蠹的经典教本《算术》时，看到了这样一个方程式：

2\textsuperscript{a}+2\textsuperscript{b}=2\textsuperscript{c}

他立刻意识到有无穷多个a、b、c能满足这个方程，于是在书的边空上写下了下面这段臭名昭著的批注：

方程

n\textsuperscript{a}+n\textsuperscript{b}=n\textsuperscript{c}

其中a、b、c及n取正整数，仅当n=2时有解（此时将有无穷多的三元组a、b、c满足这个方程）；但n＞2时没有解。我已找到了一个精彩的证明，可是不凑巧，这个证明太小了，如果写在这块边空上几乎会看不清。

从三百多天前的那年起，数学家们绞尽脑汁，试图做到：或者证明狒犸的断言，从而维护狒犸的声誉------因为由于一些人怀疑他并没有真的找到他所说的那个证明，狒犸的名声多少有些不佳了------或者否定狒犸的断言，找出一个反例：给出四个整数a、b、c和n，其中n\textgreater2，满足那个方程。直到最近，这两个方向上的努力都是失败的。诚然，这个猜想已经在许多具体的n值上被验证是成立的------即那些小于12500的n值。但是没有人成功地证明它对所有的n值都成立------没有人，这就是说，直到蚁翰·塞巴斯蚁安·蜚蚂出现之前。是他发现了那个证明，从而保住了狒犸的名誉。现在它的名字是``蚁翰·塞巴斯蚁安的平均验刚劲猜想''。

\dialogname{阿基里斯：}如果正确的证明最终给出了，是不是应该把它称作``定理''而非``猜想''？

\dialogname{食蚁兽：}严格说来，你是对的，但习惯上一直这么叫。

\dialogname{乌龟：}塞巴斯蚁安写了些什么样的音乐？

\dialogname{食蚁兽：}他在作曲方面极有天赋。不幸的是，他的最伟大的作品被罩在一层神秘的网下，因为他从未发表过它。有人确信那整部作品都在他的心里，另一些人则不太客气，说他也许压根儿就没有作过这支曲子，只不过是吹吹牛而已。

\dialogname{阿基里斯：}这部巨作是什么性质的？

\dialogname{食蚁兽：}是一部庞大的前奏曲和赋格，其赋格将有二十四个声部，涉及二十四个不同的主题，每一主题各有一个大调和一个小调。

\dialogname{阿基里斯：}把二十四个声部的赋格作为一个整体来听一定很难！

\dialogname{螃蟹：}更别说创作这么一部曲子了！

\dialogname{食蚁兽：}但是我们所知道的有关它的一切就是塞巴斯蚁安对它的描述，他把这段描述写在了他的那本为管风琴而写的幂幂地加在一起的前奏曲和赋格曲的空白处。在他悲剧性的猝亡前，他写的最后几句话是：

我已写成了一部精彩的赋格。在里面我将调的24次幂与主题的24次幂加了起来，作出了一部具有24个声部幂的赋格。可是不凑巧，这里的边空太小了，我无法写下它。

这部未能实现的杰作就被命名为《蜚蚂的最后赋格》。

\dialogname{阿基里斯：}哦，这真是个叫人受不了的悲剧。

\dialogname{乌龟：}说起赋格，我们一直在听的这支赋格就要结束了。在接近结尾的时候，主题会出现一个奇怪的新转折。\kaitismall{（翻阅《平均律钢琴曲集》）}哎，这是什么？又是一幅插图------多吸引人呀！\kaitismall{（给螃蟹看。）}

\dialogname{螃蟹：}哎，这是什么？哦，我明白了：是``整化论''，是用大个儿字写的，先是缩小，然后又恢复到原来的尺寸。可这毫无意义，因为这不是一个词。嘿，怎么搞的，嘿，真是的！\kaitismall{（把它递给食蚁兽。）}

\dialogname{食蚁兽：}哎，这是什么？哦，我明白了：是``简体论''，是用小个儿字写的，先是变大，然后又缩小到原来的尺寸。可这毫无意义，因为这不是一个词。嘿，真是的，嘿，怎么搞的？\kaitismall{（把它递给阿基里斯。）}

\dialogname{阿基里斯：}我知道你们大家肯定不会相信：可事实上这幅图是由``整体论''一词组成的，方块儿字从左向右不断缩小，``体''字是草体，所以老蟹认错了。\kaitismall{（把它还给了乌龟。）}

\pandocbounded{\includegraphics[keepaspectratio]{./media/OEBPS/Images/img_064.png}}

图64．一支整体论-简化论的``螺旋桨''

刘皓明绘

\dialogname{乌龟：}我知道你们大家肯定不会相信，可事实上这幅图是由``简化论''一词组成的，方块儿字从左向右不断增大。``化''字是草体，所以食大夫认错了。

\dialogname{阿基里斯：}这次我终于听到主题的那个新转折了！真感谢你给我的指点，龟兄。我想我最终还是开始掌握听赋格的艺术了。

\phantomsection\label{Chapter11.xhtml}{}

第十一章

\chapter{大脑和思维}

关于思维的新观点

只是随着计算机的出现，人们才真正试图制造``能思维的''机器，并见识了思维这个主题上各种稀奇古怪的变奏。一些程序设计出来了，但若把它们的``思维''与人的思维相比，那就如同把鹦鹉的叽喳学舌与人说话相比一样。这种新发现出来的能力，即用具有与人类完全不同的、但却是人工构造的思维形式------或思维的近似物------来做实验的能力，突然之间揭示出了人类思维的特性、缺陷、威力、捉摸不定之处和发展变化之道。结果是近二十年来，我们在关于思维是什么、不是什么这一问题上，获得了一些新的观点。与此同时，脑的研究者们对于大脑硬件的小尺度及大尺度的探索也都大有收获。这种探索目前还不足以清楚地揭示出脑是如何处理概念的，但已使我们对思维过程所凭藉的生物机制有了一些认识。

在下面两章中，我们将努力把几个领域的研究成果结合起来，包括实现计算机智能的尝试、在活的动物脑上所做的精巧实验以及认知心理学家对人的思维过程的探索。《前奏曲，蚂蚁赋格》已经为此做了铺垫，现在我们来进一步展开其中的观点。

内涵与外延

思维必须依赖于在大脑硬件中对客观实在的表示。在前面几章中，我们已经开发了几个形式系统，数学实在的领域可以在那种符号体系中表示出来。用这样的形式系统作为大脑如何支配思想的模型，在多大程度上是合理的？

在pq系统以及后来的几个更复杂的系统中，我们已经看到了意义------在这个词的受到一定限制的意思下------是如何作为一种同构的结果而出现的。这种同构把印刷符号映射到数、运算和关系，把印刷符号串映射到陈述。在大脑中我们没有印刷符号，但我们有更好的东西：能动的成分，它们可以存贮、传送并从其它能动成分接收信息。也就是说，我们有主动的符号，而不是被动的印刷符号。在大脑中，规则就混在这些符号之内（而在纸上，符号是静止实体，规则是在我们的头脑里）。

要强调的一点是：从我们已经见到的那些形式系统所具有的相当严格的性质中不应产生这样的想法：符号与现实事物之间的同构是刻板的一一映射，就像木偶与牵着它的那根线一样。在TNT中，``50''这个概念可以有不同的符号表示公式，例如

((SSSSSSS0•SSSSSSS0+(S0•S0))

((SSSSS0•SSSSS0)+(SSSSS0•SSSSS0))

两者表示同一个数，这一点并不是自明的。你可能分别处理每个表达式，后来在某个时刻你偶然碰上了一个定理，不禁叫道``啊------它就是那个数！''

你的心智同样可能对同一个人有不同的描述，例如

我不久前送给一个在波兰的朋友的那本书的作者。

今晚在这个饭馆里同我和我的朋友聊天的这个陌生人。

两者表示同一个人，这一点并不是自明的。这两个描述可能互不相干地存在于你的脑海中。在这个夜晚的某个时刻你们可能偶然谈起了某个话题，由此导致你发现二者指示同一个人，不禁叫道``啊------你就是那个人！''

对于一个人的各种描述并非都要联系于某个存贮着他姓名的中心符号。描述本身就可以被构造和处理。我们可以通过构造一个描述来创造不存在的人；当发现两个描述表示同一实体时，我们可以使它们合为一体；当发现一个描述表示着两个事物时，我们可以把它分成两半------如此等等。这种``描述的演算''存在于思维的核心之中。它被称作是``内涵的''而非``外延的''，意思是描述可以``漂浮''着，而不落实在具体的已知客体上。思维的内涵性与其灵活性相关联，它使我们能够想象假设的世界，合并不同的描述，或把一个描述砍成分离的几段，等等。

假设一个借走了你自行车的朋友打来电话说，她被一辆汽车撞了，你的自行车彻底完蛋了，她也差点一命呜呼。你于是在脑海里想象出一系列场面。随着她不断地添加细节，这些场面越来越生动，到最后你简直就像亲眼见到了一样。然后她告诉你这一切只不过是一个玩笑，她和自行车全都安然无恙！很大程度上这其实无济于事：那个故事和那些场景的生动性丝毫不受影响，而且这段记忆将伴随你很长很长时间。以后，你甚至会由于这一印象的强烈影响而认为她骑车是靠不住的，而这些印象在你得知它们并不真实之后，本应已被抹掉了。幻想和事实在我们的头脑里紧密地混在一起，这是由于思维牵涉到复杂描述的构造和处理，而这些描述并非一定要束缚于真实的事件或物体。

一个关于世界的灵活的内涵表示是思维的全部所在。那么，像脑这样一个生理系统是怎样支持这样一个表示系统的呢？

大脑中的``蚂蚁''

\pandocbounded{\includegraphics[keepaspectratio]{./media/OEBPS/Images/img_065.png}}

图65．神经元示意图

摘自伍尔德里奇{[}D.Wooldridge{]}，《大脑的机制》{[}The Machinery of the Brain{]}，纽约：Me-Graw-Hil1，1963年版，第6页。

大脑中最为重要的细胞是神经细胞，或称神经元（见图65），它们大约有一百亿个（奇怪的是，胶质细胞的数目是神经元的十倍。据信，它们主要是对神经元这个主角起支持作用的，因此我们这里不讨论它们）。每个神经元具有若干个突触（``输入端口''）和一个轴突（``输出通道''）。输入和输出均为电化学流，即移动的离子。神经元的输入端口和输出通道之间是细胞体，``决定''就是在这里作出的。

一个神经元所面临的是这样一种决定------这种决定每秒要出现一千次------即是否``发射''。所谓``发射''就是沿其轴突释放离子，而这些离子最终将穿入一个或多个其它神经元的输入端口，致使它们作出同类决定。决定是以一种非常简单的方式作出的：如果所有输入的总和超过了一个确定的阈值，则发射；否则，不发射。有些输入可能是否定输入，它们对来自其它地方的肯定输入起抵消作用。无论如何，都是这种简单的加法在支配着心智的最低层。笛卡尔的著名论断``我思故我在''可以改写成``我思故我算''。

虽然作出决定的方式听起来很简单，但下述事实使问题复杂化了：一个神经元所具有的不同输入端口可以多达200000个，这意味着为决定该神经元的下一步行动需处理200000个被加数。决定一旦作出，一个离子脉冲将沿轴突射向其末端。但是，当离子到达轴突的末端前，可能会遇到一个或多个分叉点。在这种情况下，单一的输出脉冲在通过分叉点时将分裂，在到达末端时``它''已经变成了``它们''------而且它们可能在不同的时刻到达各自的目的地，原因是它们经过的轴突分支可能具有不同的长度和不同的阻抗。但重要的是，它们都始于一个出自细胞体的单一脉冲。在一个神经元发射后，它需要一个短暂的恢复时期才能再次发射，这是以毫秒来度量的，因此一个神经元的发射频率可达每秒上千次。

脑的大尺度结构

至此我们已经描述了大脑中的``蚂蚁''。而``队''或``信号''是怎样的呢？``符号''又是怎样的呢？我们看到：尽管单个神经元的输入很复杂，它却只能以一种非常基本的方式作出反应------发射，或不发射。这只具有很少量的信息。在传递加工大量信息的过程中，显然必须包括许多神经元。因此可能的猜测是：存在着由许多神经元构成的大尺度结构，它们在一个较高的层次上处理概念。这毫无疑问是对的，但最朴素的假设------每个不同的概念对应于一个固定的神经元群------却基本上是错的。

\pandocbounded{\includegraphics[keepaspectratio]{./media/OEBPS/Images/img_066.png}}

图66．人脑左视图

奇怪的是视觉区在头的后部。{[}摘自史蒂文·罗斯，《有意识的大脑》修订本，第50页{]}

人脑从解剖学的角度看可以分成若干彼此相区别的部分，如大脑、小脑和下丘脑（见图66）。大脑是人脑中最大的部分，它可分成左半球和右半球。每个半球的外面几毫米被分层的``壳''包裹着，这就是大脑皮层。从解剖学上看，大脑皮层的总量是人脑区别于低智能物种的脑的重要特征。我们不准备详细描述大脑的任何子器官，因为迄今为止在这些大尺度子器官和它们所负责的心智或物理活动之间只能建立起最粗略的对应。例如，已经知道语言主要是由两个半球之一来掌握的------事实上通常是左半球。我们还知道小脑的作用是发送一系列脉冲到肌肉，以控制运动。但是，这些区域是如何完成它们的功能的，这在很大程度上还是个谜。

脑之间的映射

说到这里，一个极其重要的问题出现了。如果思维的确是在脑中发生的，那么两个脑是如何彼此相区别的呢？我的脑和你的脑有什么不同呢？你的思想当然和我的不会完全一样，对其他任何人来说也如此。但我们的大脑都具有相同的解剖学划分。脑之间的这种一致性可以扩展到什么程度？能达到神经元层次上吗？如果你观察那些在思维的层次体系中处于足够低的位置上的动物------例如低等的蚯蚓，答案将是肯定的。下面的引文摘自神经生理学家大卫·休贝尔在一次会议上的讲话，这次会议的议题是与地球外的智能生物通讯。

我估计像蚯蚓这样的动物会拥有数以千计的神经细胞。非常有趣的是，我们可以在一条特定的蚯蚓身上标出一个特定的单个细胞，然后在同一品种的另一条蚯蚓身上找到一个与之相对应的细胞。\footnote{}

蚯蚓具有同构的脑，人们甚至可以说，``只存在一条蚯蚓。''

但是，随着你观察的对象在思维层次体系中的上升及其神经元个数的增加，这种不同个体的脑之间的一一映射关系很快就消失了------这支持了你的猜疑：不是只存在一个人！然而，一旦在比单个神经元大、又比脑的主要子器官小的尺度上把不同人的脑相比较，仍然可以发现显著的物理相似性。个体的精神差异是如何表现在脑的物理结构上的？上述事实对这个问题来说意味着什么？如果我们来观察我的神经元的相互联系，我们是否能发现各种结构，它们可以被确定为以编码的形式表示了我所知道的特殊事物、我所具有的特殊信念、我怀有的特殊愿望、恐惧和好恶？如果精神体验可以归因于大脑，那么知识和精神生活的其它方面是否也能类似地追溯到大脑中的特定位置，或大脑特定的物理子系统？这是我们在本章和下一章中要不断涉及的一个中心问题。

大脑过程的定位：一个谜

为了回答这个问题，神经病学家卡尔·拉施利从1920年左右开始进行了多年实验，试图发现老鼠在脑中的哪个地方存贮了它跑迷津的知识。史蒂文·罗斯在他的著作《有意识的大脑》中这样描述了拉施利的尝试与磨难：

拉施利试图确定记忆在皮层中的位置，为此，他首先训练老鼠跑迷津，然后切除皮层的不同区域。他让动物进行恢复，然后检查跑迷津技能的保持情况。出乎他意料的是，要找到一个迷津通路的记忆能力所对应的特定区域是不可能的。相反，所有进行了皮层局部切除的老鼠都受到某种损害，而损害的程度大致正比于皮层的切除量。皮层切除破坏了动物的运动和感觉能力，它们可能一瘸一拐，步态蹒跚，但总能设法穿过迷津。就记忆而论，皮层似乎是等势的，这就是说，所有的区域都有相同的可能效用。后来，在他发表于1950年的最后一篇论文``寻找记忆痕迹''中，拉施利沮丧地断定唯一的结论是记忆根本就是不可能的。\footnote{}

奇妙的是，在40年代末，几乎在拉施利进行他的最后的工作的同时，支持相反观点的证据在加拿大出现了。神经外科医生怀尔德·潘菲尔德那时正在检查进行了某种大脑手术后的患者的反应。这种手术是把电极插入患者被暴露了的大脑的各种部位中，然后用微弱的电脉冲刺激那些与电极相接触的神经元或神经元群。这些脉冲和来自其它神经元的脉冲很相似。潘菲尔德发现：对某些神经元的刺激会可靠地给患者带来特定的幻像和幻觉。这些人工诱发的感觉是各式各样的，从陌生而不可名状的恐惧感到嗡嗡声及颜色等等不一而足。最令人印象深刻的是，其中还包括早年生活中事件的完整过程，比如童年时代的一次生日晚会。能触发这种特定事件的位置只限于一个很小的区域------基本上集中于一个单个神经元。于是潘菲尔德的这些结果戏剧性地回击了拉施利的结论：说到底，这些结果似乎意味着特定的记忆是由局部区域负责的。

对此我们该怎样理解呢？一种可能的解释是：记忆是局部编码的，但这种局部编码在皮层的不同区域反反复复地进行着------这可能是进化中发展起来的防卫策略，用来对付在战斗中或在神经生理学家所进行的实验中可能出现的皮层损失。另一种解释是：记忆可以从分布于整个大脑的动态过程之中重建，但可以从局部点触发。这一理论是基于现代电话网络的观念之上的，在这种网络中，一个长途呼叫的传送途径无法事先预料，因为这是在呼叫发出时才选定的，取决于整个电话网的状况。毁坏网络的任何局部都不会阻塞呼叫，而仅仅使它们在传送时绕开损坏的地区。在这个意义下，任何呼叫都潜在地是不可定位的。然而，任何呼叫都恰好联结了两个特定的点，在这个意义下，任何呼叫都是可定位的。

视觉处理的特性

一些有关大脑过程定位的最有趣、最有意义的研究工作，是在近十五年内由大卫·休贝尔和托斯坦·威瑟尔在哈佛完成的。他们绘出了猫脑中的视觉通道，它起自视网膜中的神经元，沿它们那些导向头部后方的路径通过侧膝体的``中继站''，终止于大脑正后部的视觉皮层。首先，由于有拉施利的结果，存在有界说良好的神经通道这一点是不同寻常的。但更不同寻常的是位于通道不同阶段上的神经元的性质。

已经发现，视网膜神经元基本上是对比感觉器。更具体地说，它们的活动方式是这样的：每个视网膜神经元通常都以一个``常规速率''发射。当视网膜上它所在的部分被光照射时，它的发射速率或是加快或是减慢，甚至可能停止发射。但是，它这样做的条件是视网膜的周围部分受到的光照较少。这意味着存在两种神经元：``开中心''和``关中心''。开中心神经元的特征是：在它们敏感的圆形视网膜区域中，当中心亮外围暗时，它们的发射率会上升；关中心神经元的特征是：当中心暗外围亮时，它们发射得快一些。如果把一个开中心模式显示给一个关中心神经元，它的发射率会下降（反之亦然）。均匀的光照将不影响这两种视网膜神经元，它们仍继续以常规速率发射。

从视网膜开始，信号从这些神经元出发，通过视神经到达接近大脑中部的侧膝体。在那里，我们可以发现视网膜表面的一个直接映像。之所以这样说，是因为侧膝体神经元只能被落在视网膜特定区域的特定刺激所触发。在这个意义下，侧膝体消失了，它似乎仅仅是一个``中继站''，而非进一步的处理器，（虽然公平地说，对比敏感性在侧膝体中有所增强）。视网膜图像按照侧膝体中神经元的发射模式以一种直接的方式编码，虽然事实上那些神经元不像视网膜那样分布于一个二维平面上，而是分布在一个三维团块中。这就既保持了信息又把二维映射到三维，形成了一个同构。表示方式维数的变化可能还有更进一步的意义，这一点还没有充分揭示出来。不管怎么说，尚未得到解释的视觉阶段实在太多了，以致于对下列事实我们不应失望而应高兴------在某种意义上，我们已经把这一个阶段弄清楚了！

\pandocbounded{\includegraphics[keepaspectratio]{./media/OEBPS/Images/img_067.png}}

图67．某些神经元样本对模式的反应。

（a）这个负责边缘检测的神经元寻找左明右暗的垂直边缘。第一列显示了边缘的倾角与该神经元的关系，第二列说明边缘在域中的位置与这个特定神经元无关。（b）显示了超复杂细胞的更具选择性的反应：在这里仅当下行的舌状物通过域中间时反应最强。（c）假想的``祖母细胞''对各种随机刺激的反应。读者不妨思考一下对同一组刺激一个``章鱼细胞''将作何反应。

从侧膝体出发，信号向后传至视觉皮层。在这里某些新的加工出现了。视觉皮层细胞分为三类：简单的、复杂的和超复杂的。简单细胞的活动方式和视网膜细胞或侧膝体细胞相似：它们对视网膜特定区域中与背景形成对比的亮点或暗点作出反应。而复杂细胞则通常要从上百个其它细胞接收信息作为输入，识别视网膜上成特定倾角的明暗线条（见图67）。超复杂细胞对应于沿特定方向运动的角、线条甚至于``舌状物''（也见图67）。最后一种细胞是高度专门化的，因此它们有时被称为``高阶超复杂细胞''。

一个``袓母细胞''？

由于在视觉皮层中发现了能被复杂性不断增加的刺激所触发的细胞，有些人怀疑事情是否在向``一个细胞，一个概念''的方向发展------例如，你可能有一个``祖母细胞''，它当且仅当你的祖母进入视野时发射。这样一个``极端超复杂细胞''的多少带点幽默的例子不过是个玩笑，但是，可供选择的理论中哪个更合理，这并不是显而易见的。一种可能性是更大的神经网络被足够复杂的视觉刺激整体性地激活。当然，这些更大的多神经元单位的触发，可能以某种方式依赖于许多超复杂细胞发出的信号的整合。没人知道这是怎样完成的。正当我们似乎接近了``符号''怎样从``信号''中形成这一关键点时，线索中断了------一个吊人胃口的``且听下回分解''。但我们很快会回到这个故事，并努力再补充一些内容。

前面我提到过，人类的大脑之间在大的解剖学尺度上存在粗糙的同构，还提过蚯蚓的脑之间存在非常精密的神经元层次的同构。很有趣的是在猫、猴和人的视觉处理器官之间也有一个同构，其``粗细程度''介于粗糙和精密之间。这个同构是这样的：首先，三个物种都在脑的后部为视觉加工``特供''了一部分皮层，即视觉皮层。其次，在三者之中视觉皮层都分成三个子区域，称为皮层的17、18和19区域。这些区域仍具有普遍性，这就是说在三个物种的任何正常个体的脑中都可找到它们的位置。在每个区域中还可以再细分，一直到视觉皮层的``柱状''组织。柱状组织垂直于皮层表面，沿半径方向伸向内脑，其中绝大多数的视觉神经元间的联系是沿脑半径方向，即柱体的方向，而非建立于柱体之间。每个柱体映射到一小块特定的视网膜区域上。柱体的数目在不同的个体中是不同的，因此不可能找到``同一个柱体''。最后，在柱体内部，某些层主要由简单神经元构成，而另一些层主要由复杂神经元构成。（超复杂神经元往往在区域18和19中占数量优势，而简单和复杂神经元主要出现于区域17中）。看起来我们的分析细到这个程度就不会再有同构了。从这里直到单个神经元层次，每个猫、猴或人都有完全唯一的模式------有点像指纹和签名一样。

在猫脑和猴脑中的视觉加工之间，存在一个虽然次要但或许能说明问题的差别，这与把来自两眼的信息整合成单一的复合高层信号的时刻有关。研究结果表明，在猴脑中这一整合发生的时刻稍迟于猫脑，致使每个单眼信号都有稍长的独立加工时间。这并不值得惊讶，因为可以设想，一个物种在智能层次体系中所处的位置越高，其视觉系统所要处理的问题就越复杂，因此信号在最终得到一个``标签''前所要经过的预处理也就越来越多。对初生牛犊视觉能力的考察已经非常戏剧性地证实了这一点。小牛似乎一生下来就已经达到它将来所可能具有的视觉分辨力。它会胆怯地回避人或狗，但不怕别的牛。有可能它的整个视觉系统在出生前就被``硬性联接''了，而且只涉及较少的视觉皮层。与此相反，人的视觉系统深深地依赖于皮层，而且要经过若干年才能成熟。

汇集到神经模块

令人迷惑不解的是，在迄今为止关于脑组织的发现之中，并没有找到多少大尺度硬件和高层次软件之间的直接对应关系。举例来说，视觉皮层就是块大尺度硬件，它的软件用途是很清楚的------视觉信息加工，但至今发现的所有加工过程仍都是层次很低的，远没有接近对象识别在皮层中的定位。这意味着没有人知道复杂细胞和超复杂细胞的输出在哪里或怎样转换成对形状、位置、图像、面容等等有意识的识别。人们已在寻找证据，试图证明许多低层次的神经反应``汇集''成越来越少的高层反应，最终达到像所谓的祖母细胞一类的东西，或是前面提到过的某种神经元网络。显而易见，这个发现不会出于某种粗略的解剖学划分，而是需要进一步的微观分析。

替代祖母细胞的一种可能解释是，在这一``汇集''过程的终点处有一组神经元，比如说有数十个，当祖母进入视野后它们全都发射。对于每个不同的可识别的对象，都有唯一的一个神经元网络和一个``聚焦''于该网络的汇集过程。沿着类似的思路还可以提出更复杂的解释，包括可以用不同的方式而非固定方式激活的网络。这样的网络就是我们大脑中的``符号''。

但这种``汇集''是必须的吗？或许一个被辨认的物体就是以它在视觉皮层上的``印迹''来隐含地标识的------所谓``印迹''，是指来自简单、复杂和超复杂细胞的集体反应。或许大脑并不需要什么进一步的识别机制来识别一个特定形状。但是，这种理论将造成一系列问题。假设你在欣赏风景，这在你的视觉皮层上留下了印迹，但此后你怎样才能从中得到关于这个景象的词语描述呢？举例来说，法国后印象派画家爱德华·维亚尔的作品，往往需要几秒钟的审视，然后一个人的形象才突然呈现在你面前。大概只须片刻其印迹就能形成于视觉皮层上------但这幅画是在几秒钟后才被理解的。这只是一个例子，实际上这种现象很普遍------在识别时好像感到有什么东西在你的脑海里``结晶''了，这一过程不是发生在光线射在你视网膜上的时候，而是稍迟一些，要等到你智力的某些部分有机会作用于视网膜信号之后才行。

``结晶''这个比喻从统计力学的角度提供了一幅恰当的图像。在这个图像里，介质中的无数微观的互不相干的活动缓慢地造成局部区域的内在一致性，这一区域不断扩展，最终这无数微观事件会对其介质产生根本性的结构改造，使其从一个杂乱无章彼此独立的元素集合体，变为一个大的、一致的、内在联系紧密的结构。如果你把早期神经元活动看作独立的，而把它们这许多独立发射的最终结果看作是触发了一个界说良好的大神经元``模块''，那么``结晶''这个词似乎再贴切不过了。

汇集过程的另一个论据基于如下事实：使你感到你看见了同一个对象的景象可以有无数多个------例如你的祖母，她可能喜也可能愁，可能戴帽子也可能没戴，可能在明亮的花园里也可能在昏暗的车站中，可能近也可能远，可能侧对着你也可能面向着你，如此等等。这些场景会在视觉皮层上产生极其不同的印迹，但它们都能使你叫``喂，奶奶''。因此，一个汇集过程必定出现于接收到视觉印迹后，但又在词语发出前的某一时刻。有人可能主张这种汇集不是对祖母的感知的一部分，而仅仅是词语表达过程的一部分。但把过程做这样的划分看上去很不自然，原因是你可以在内部使用``这是奶奶''这一信息，同时又没有任何词语表达。当然，在整个视觉皮层上处理全部这些信息未免太笨拙了，实际上其中大部分都该扔掉，因为你根本不关心影子落在哪里或她罩衫上有多少扣子这类事情。

如果在理论上不承认存在有汇集过程，那么另一个困难是需要说明为什么对同一个印迹会有不同的解释------例如艾舍尔的画《凸与凹》（图23）。正如这样一个显而易见的事实；我们从电视屏幕上感知到的不仅仅是点，还有组块，而且，假定当一片点状印迹产生于视觉皮层上之后感知就已经发生，这似乎有些荒唐。一定是有某种汇集，其最终效果是触发某些特定的神经元模块，而每个模块都关联于场景中的概念------或称组块。

作为思维过程媒介的模块

于是我们可以作出这样的结论：对应于每个概念，都存在一个界说良好的可触发模块------由一小群神经元构成的模块------也就是前面设想过的那种``神经复合体''。这个理论有个问题------至少朴素地看是如此------它使人想到应该能够在大脑里确定这些模块的位置。迄今为止还没有做到这一点，而且存在一些反对定位化的证据，如拉施利的实验。但是，现在下结论未免为时过早。每个模块都可能有多个复本分布在各处，也可能模块在物理上相互重叠，这两种效应都倾向于使把神经元分成``包裹''的界限变得模糊不清。或许复合体像摞成一摞的纸张，偶尔会互相交叉；或许它们像互相缠绕的长蛇，在这里或那里像眼镜蛇的头那样伸出来；或许它们像蜘蛛网；或许它们是某种电路，信号在里面转来转去，比扑向蚊虫的燕子还难以捉摸。到底是什么谁也不知道。这些模块甚至可能是软件现象，而非硬件现象------这一点我们过一会儿再讨论。

这些假设性的神经元复合体可以使人想到许多问题，例如：

\begin{itemize}
\tightlist
\item
  它们是否延伸到了脑的内部区域，如中脑和下丘脑等处？
\item
  单独的一个神经元能否属于一个以上这种复合体？
\item
  单独一个神经元可以属于多少个这种复合体？
\item
  一个这种复合体可覆盖多少神经元？
\item
  对每个人来说，这些复合体是否大致相同？
\item
  在不同的人脑中，是否能在对应的位置找到对应的复合体？它们在每个人脑中的覆盖方式是否相同？
\end{itemize}

从哲学上看，所有这些问题中最重要的是：这些模块------例如一个祖母模块------的存在说明了什么？这是否使我们对自身的意识现象有所认识？还是说这和仅仅知道``大脑是由神经元和胶质细胞所组成''一样，仍使我们对意识是什么一无所知？也许你在读《蚂蚁赋格》时已经猜到，我的感觉是，要想理解意识现象，我们还得经过一个漫长的过程。需要完成的关键步骤是：对同一个脑的同一个状态来说，低层次的描述------面向神经元的------要变成高层次的描述------面向模块的。或者再一次重申《蚂蚁赋格》中带有启发性的术语，我们需要的是把脑状态的描述从信号层转向符号层。

活跃的符号

让我们从现在开始把这些假设性的神经复合体、神经模块、神经包、神经网络、多神经元单元------你可以随便叫它们什么，也不管它们的形状是一摞纸、园艺耙还是眼镜蛇、小雪花------统称作``符号''。用符号的说法来描述大脑状态，这在对话中已经提到过了。这样一个描述会是什么样的呢？把什么样的概念设想成为``符号化''了的才算是合理的？符号之间会有什么样的相互关系？这幅图景又为对意识的认识提供了什么见解呢？

首先要强调的是，符号既可能是休眠的，也可能是觉醒的（激活的）。活跃的符号就是那些已经被触发了的------即其中由外部刺激造成发射的神经元数目已经达到了阈值。由于一个符号可以通过许多不同方式来触发，因此它在觉醒状态下可以按许多不同的方式活动。这提示我们不应把符号看成一个固定的实体，而应看成一个可变的实体。因此，仅仅通过说``符号A、B、C\ldots\ldots N是所有活跃符号''来描述大脑的一个状态是不够的，我们还必须为每个活跃符号提供一组附加参数，以刻划该符号内部工作的某些方面。一个有趣的问题是，在每个符号中是否都有某些核心神经元，它们在符号被激活时一定是发射的？如果这样一个神经元核心团存在，我们可以称之为该符号的``恒定核心''。作下面这样的假设是很诱人的：每当你想到一个事物，例如瀑布，某种固定的神经过程就重复进行，这些过程无疑在细节上依当时的情景会有不同，但总是可靠地发生。不过，是否果真如此还不清楚。

那么一个符号在觉醒时做些什么呢？一个低层次的描述可能是：``它的许多神经元发射了''。但我们已经对此不感兴趣了。高层次的描述应消除一切涉及神经元的说法，而完全集中于符号。这样，一个关于符号如何是处于活跃状态而非休眠状态的高层次描述就会是``它发出消息------或称信号------其目的是试图唤醒------或触发------其它符号''。当然这些消息是被神经元以神经脉冲流的方式承载的------但我们应尽可能避免使用这种表达方式，因为这代表着一种低层次的观察事物方法，而我们希望能完全在高层次上进行描述。换句话说，我们希望思维过程可以被设想成与神经事件相隔离，正好像一座钟的行为可以与量子力学定律相隔离，或细胞生物学可以和关于夸克的定律相隔离一样。

但这样一幅高层次的图景有什么好处呢？为什么``符号A和B触发了符号C''比``神经元183通过612刺激了神经元25并导致了它的发射''要好一些？在《蚂蚁赋格》中回答了这个问题：这是因为符号把事物符号化了，而神经元不是如此。符号是概念的硬件实现。一组神经元触发另一个神经元并不对应于外部事件，而一些符号对某个符号的触发却以某种方式对应着现实世界------或某个假想世界------中的事件。符号间通过来回传送消息来保持联系，其方式使得符号的触发模式与在我们的世界中发生的------或在一个类似于我们的世界中可能发生的------大尺度事件十分相像。意义在这里出现的原因实质上恰如在pq系统中的情形一样------同构。只不过在这里这个同构已变得极其复杂、微妙、精巧、多面化，而且是内涵的。

而且，符号能够来回传递复杂的消息，这一点就足以使得神经元自身无法充当符号的角色。因为神经元仅用唯一的方式发送信息，不能带有选择性地使一个信号一会儿发往一个方向，一会儿又发往另一个方向，所以它完全不具有选择性触发能力，而这种能力是符号在模拟现实世界中的对象时所必须具有的。威尔逊在他的著作《昆虫社会》中，指出了蚁群内部消息传播的类似特点：

（大众传播）可定义为在群体中并非以一个个体传向另一个个体的方式所进行的信息传递。\footnote{}

看来把大脑想象成一个蚁群并不算太荒唐！

下一个问题------这也是一个极端重要的问题------与大脑中单个符号所代表的概念的性质和``尺寸''有关。关于符号的性质存在这样一些问题：到底是存在一个代表瀑布的一般观念的符号，还是各个具体的瀑布对应于不同的符号？或者上述两种方案同时都实现？关于符号的``尺寸''存在这样一些问题：是否存在一个代表一个完整的故事、一段旋律、或一个笑活的符号？或者，作为更大的可能，符号只相应于像词这种规模的概念，而更大尺度的，如短语和句子，可以表示成不同符号的同时或顺序的激活？

让我们考虑一下符号所表示的概念的规模这一问题。大部分表达成句子的思想都是由一些基本的、像原子似的成分所构成，而对这些成分我们一般是不作进一步的分析的。这大致就是词的规模------有时长些，有时短些。例如，名词``瀑布''、专有名词``尼亚加拉瀑布''、表示完成的虚词``了''、动词``赶上''以及更长的固定短语都是接近于原子化的。这都是一些典型的基本笔触，我们就是用它们来绘出更复杂概念的肖像，例如一部电影的情节、一座城市的风情、意识的性质等等。这些复杂的概念都不是单一的笔触。认为语言的笔触也就是思维的笔触，这看起来是合理的。因而符号所表示的概念大概也就是这种规模的。这样，一个符号大略对应于那些你可用一个词或一个固定短语所表示的东西，或者对应于那些你指定了一个专有名称的东西。而在大脑中表示更复杂的观念------例如恋爱中碰到的一个难题------则可能需要一个非常复杂的符号激活序列。

类与例

在思维中有一种一般性的区别：范畴与个体，或类与例（另外两个有时会用到的词是``型''与``例''）。初看起来似乎一个给定符号应当固有地或者是一个类符号，或者是一个例符号，但这么想是过于简单化了。实际上，大多数符号可以扮演二者中任何一个角色，这取决于它们的激活环境。例如下列序列：

⑴一种出版物

⑵一种报纸

⑶《中国日报》

⑷五月十八日的《中国日报》

⑸我的那份五月十八日《中国日报》

⑹我的那份五月十八日《中国日报》，在我第一次拿起它的时候（而不是几天后在我的炉子中燃烧时）

在这里，第二行到第五行都扮演着双重角色。其中，第四行是第三行对应的一般类的一个例，而第五行又是第四行的一个例。第六行是一个类的一种特殊例：一个``表现''。一个对象在其生存历史中相继所处的阶段是它的各个表现。有趣的是，不知道一个农庄中的母牛是否能感知到，在那个喂它们干草的快活农民的所有表现的背后有一个不变的个体。

原型原则

看起来，上面那个序列是关于一般性的一个层次结构。其顶层是一个非常一般的概念范畴，底层是某种在确定的时空位置上的低级的特定事物。但是，如果认为一个``类''必定总是非常一般和抽象的，这就未免局限性太大了。其原因是我们的思维采用了一种机智的原则，我们可以称其为``原型原则''：

{最具体的事件可以被用作一类事件的一个一般范例。}

我们都知道，特殊事件具有一种生动性，这使得它们可以被牢牢印在记忆中，以便后来被当作在某个方面与它们相似的其它事物的模型。因此，在每个具体事件中，都蕴含着全部相似事件组成的类的萌芽。一般性即寓于特殊性之中，这个思想具有深远的重要性。

至此我们自然会问：大脑中的符号是表示类，还是表示例？是否某些符号只表示类，而另一些符号只表示例？或者说一个符号可以交替完成类符号和例符号双重职责，这取决于它的哪一部分被激活？后一种理论看起来很有吸引力，有人可能会设想一个符号的``轻微的''激活可能代表一个类，而更深入的，或更复杂的激活可能包括更细致的内部神经发射模式，因而是表示一个例。但再一想，这样太古怪了。比如，这意味着如果以足够复杂的方式激活代表``出版物''的符号，你就能得到一个非常复杂的符号，它代表着在我的炉子中燃烧的一份特定的报纸。而且所有其它印刷品的所有其它可能的特例均可通过对代表``出版物''的这个符号的某种方式的激活来内在地表示。对``出版物''这个符号来说，这样一副担子似乎太重了。因此我们只能下结论说：例符号可能和类符号同时存在，而不仅仅是后者的激活方式。

从类中分离例

另一方面，例符号的许多性质常常是从它所隶属的类中继承来的。假设我告诉你我去看了个电影，你将开始``铸造''那个特定电影的一个新的例符号。但在缺乏更多信息的情况下，这个新的例符号不得不紧密地有赖于你关于``电影''的已有的类符号。你将不自觉地依靠许多关于该电影的预先假定------例如，它的放映时间在一小时至三小时之间，是在本地一家剧场上映的，内容是关于某些人的一个故事，等等。这些都作为与其它符号的可能联系（即潜在的触发关系）而被构造在类符号中，称为``缺席选择''。在任何新铸造的例符号中，缺席选择都可以很容易地被取消，但如果没有明确地这样做，它们就会被例符号从类符号中继承并保留下来。它们将采用``规范型''（或者说类符号）所提供的一些合理猜测为你对这个新的例------例如我所看的电影------的考虑提供一些初步的根据，直到它们被取消为止。

这就像一个既无主见也没经验的孩子------他完全依赖于父母的经验和见解，只会鹦鹉学舌。但渐渐地，随着与世界的其它部分的相互作用不断增多，这个孩子获得了他自己的独特经验，并不可避免地会离开父母。最后，孩子羽翼丰满，成了大人。一个新的例也能以同样的方式通过一段时间从它的``父母''类中分离出来，凭本身的资格成为一个类或一个原型。

为了形象化地描述这种分离过程，假设在某个周末的下午你打开了收音机，碰巧收听到两个``随机的''球队之间的一场橄榄球赛。开始你不知道其中任何一个队员的名字。此时解说员说道：``回文斯接住球以后，就一直往前，达阵得分了！''这时你所记住的只不过是某个队员接到了另一个队员的球，并达阵得分了。因此这种情况是类符号``橄榄球队员''被激活，同时以某种方式被并行激活的还有关于带球的符号。但后来回文斯又出现在一些关键的时刻，你开始专门为他建立一个新的例符号，或许以他的名字为中心。这个符号像个孩子似地依赖于类符号``橄榄球队员''：你关于回文斯的大部分想象都来自于你在符号``橄榄球队员''中存放的关于橄榄球队员的规范型。但渐渐地，随着你得到的信息的增多，``回文斯''这个符号变得越来越具有自主性，对需要并发激活其父母类符号的依赖性逐渐减少。只要回文斯做出一些精彩的动作，表现出众，上述过程可能仅在几分钟内就发生。但他的队友们仍可能全由类符号的激活所代表。最后，可能在几天之后，你在报刊的体育专栏中读到了某些文章，这时``脐带''断了，回文斯可以独立存在了。至此你知道了他的故乡、他的爱好等情况，你记住了他的相貌，如此等等。到了这时，回文斯不再仅仅被认作是一个橄榄球运动员，而是首先是一个人，他正好也是个橄榄球队员。``回文斯''这个例符号于是可以在其父母类符号（橄榄球队员）处于休眠状态时被激活了。

回文斯这个符号一度是围绕其父母符号运转的一颗卫星，正像一颗绕地球转的人造卫星一样，只是后者要大得多也重得多。然后进入一个中间阶段，此时虽然一个符号比另一个重要，但它们可以被看成是互相围绕的------有点像地球和月亮。最后，新符号变得非常自主了，此时它可以很容易地成为类符号，并开始有新的卫星绕着它转------一些相应于其它人的符号，他们还不为你所熟悉，但和回文斯有某些共同之处，因此你可用回文斯作为他们的临时规范型。等到你获得更多的信息后，这些新符号也随后都具有自主性了。

搞清符号间的纠葛是很难的

一个例在类中的生长以至最终脱离所经过的各个阶段，可以通过有关符号的联接方式的不同来区分。毫无疑问，有些时候很难分清在何处一个符号终止而另一个开始。一个符号怎样才算比另一个符号``更为活跃''？如果一个符号的激活可以独于另一个符号，那么称它们为自主的将是很明智的。

我们在前面用到了一个天文学的比喻。有趣的是，行星的运动是一个很复杂的问题------事实上，尽管进行了几个世纪的工作，三个天体通过万有引力相互作用（例如地球、月亮和太阳）的一般问题还远没有得到解决。但在下述情形之下，有可能得到一个良好的近似解：其中一个天体的质量远比另外两个大（在这里是太阳）。此时可以把这个天体看作静止的，另外两个天体围绕它转动，最后再加上两个卫星间的相互作用。但这种近似要求先把系统分解成太阳和一个``簇''：地球-月亮系统。这是一个近似，但能使我们对该系统有很深入的理解。这样，这个簇在多大程度上是现实的一部分，在多大程度上是人脑所臆造的，是人强加给宇宙的一种结构呢？这种在感知到的自主或半自主的簇之间所划的界限的``现实性''问题，当联系于大脑中的符号时，会带来无穷无尽的麻烦。

一个非常令人迷惑的问题是``名词复数''这样一个简单的问题。比如说，我们是怎样形成一个茶杯中有三只鸡，或一个电梯里有几个人的图像的？我们是不是从``鸡''的类符号开始，然后从它上面抹下三个``复本''？也就是说，我们是不是用类符号``鸡''作模板，构造了三个新的例符号？还是连带地激活符号``三''和``鸡''？如果往这些想象的场景中或多或少地添加些细节，那上述两种说法都会变得难于站住脚。例如，我们显然不是为所见到过的每个鼻子、胡子、盐粒等等都分别设立例符号。我们用类符号来处理这些为数众多的项目，因而当在街上遇到留胡子的人时，我们仅仅以某种方式激活类符号``胡子''，而不是铸造新的例符号，除非我们观察得很仔细。

另一方面，一旦我们开始区分不同的个体，就不能再依赖于用一个类符号（如``人''），以分时的方式来处理全部不同的人。很清楚，对于个别的人必须设置不同的例符号。若是想靠玩``杂耍''来完成这一任务------就是说，靠一个类符号在若干种不同的激活方式（每种分别对应于某个人）之间飞来飞去------那是很荒谬的。

在这两种极端情况之间，一定存在着许多种中间情况。很可能在大脑中有多种方式进行类与例的区分，从而产生具有不同程度的特殊性的符号及符号组织，而这些方式形成一个完整的层次结构。下列不同种类的个体和符号的连带激活可能会造成具有各种程度特殊性的心理映像：

⑴．单个类符号的各种不同激活方式或深度；

⑵．若干个类符号以某种并行方式同时激活；

⑶．激活单个例符号；

⑷．在激活单个例符号的同时激活若干个类符号；

⑸．以某种并行方式同时激活若干个例符号和若干个类符号。

这又把我们带回到一个老问题：``一个符号何时才能成为脑的一个可辨别的子系统？''例如，考虑上面第2种情况：若干个类符号以某种并行方式同时激活。在所考虑的概念是``钢琴奏鸣曲''时，上述情况很容易发生（至少相应于``钢琴''和``奏鸣曲''的两个符号正被激活）。但如果这对符号连带激活的次数足够多，则有理由假设它们之间的联系会变得很强，致使它们在以适当方式被共同激活时，将像一个单元一样活动。因此，在适当的条件下，两个或更多的符号能像一个符号一样活动，这意味着数出大脑中符号的数目要比预想的更困难。

有时会出现这样的情况：两个以前互不相关的符号以一种并行方式同时被激活。它们可能联系得很紧密，似乎是一个必然的共同体，而且两个老符号的强相互作用形成了一个独立的新符号。如果上述情况出现了，对新符号怎样说才好呢？是说``它一直存在但从未被激活''，还是说它已经被``制造''出来了？

如果说这听起来太抽象，让我们举一个具体例子：对话《螃蟹卡农》。在对话的创作过程中，两个已有符号------一个代表``螃蟹卡农音乐''，另一个代表``言语对话''------被同时激发，并在外力作用下以某种方式产生了相互作用。一旦这种情况出现，后面的事情就是不可避免的了：一个新符号------是个类符号------从二者的相互作用中产生，并从此能被独立激活。那么它是否本来就是我大脑中的一个休眠符号？如果是这样，那它一定也是任何人大脑中的休眠符号，只要他具有作为其成分的符号，哪怕它从未被从那些作为成分的符号中唤醒过。这就意味着要数出任何人大脑中的符号，就必须包括全部休眠符号------即全部已知符号的全部激活类型的全部可能的排列组合。这甚至应包括睡觉时大脑创造的那些稀奇古怪的软件产物------那些在它们的主人进入梦乡时被唤起的奇特的观念混合物\ldots\ldots 这些``潜在符号''的存在表明，把大脑想象成一个处于界说良好的激活状态的界说良好的符号集合，的确过于简单化了。而在符号层次上准确刻划大脑的状态则要远远困难得多。

符号------是软件还是硬件？

由于在每个人的大脑中都存放着数量众多甚至仍在不断增加的符号，你可能会问是否有一天大脑终于会饱和------那时新符号将不再有容身之处。推测起来，这种情况的出现条件是符号不彼此重叠。------如果一个给定的神经元从不身兼二任，符号就会像电梯里的人一样。``注意：这个大脑的最大容量是350275个符号！''

但是，这并不是大脑功能的符号模型的一个必然性质。事实上，符号间的重叠和缠结很可能已成惯例，因此每个神经元也许会成为上百个符号的功能部件，而远不是只属于一个符号。这有点让人摸不着头脑了，因为如果真是这样，那不是很容易变成每个神经元都是任何符号的一部分了吗？若是如此，无论什么样的符号定位都不可能存在------每个符号都将对应于整个大脑。像拉施利的鼠脑皮层切除试验那样的结果可以因此而得到解释------但这同时意味着放弃我们开始时的想法，即把大脑分成具有不同性质的物理子系统。我们前面把符号刻划成``概念的硬件实现''就可能是个过于简单化的想法了。事实上，如果每个符号都是由和其它符号相同的神经元组成的，那说``不同的符号''还能有什么意义呢？什么是一个给定符号的激活特征？就是说，怎样才能把符号A的激活与符号B的激活区别开？我们的整套理论不就要付之东流了吗？而且即使不存在符号间的完全重合，符号之间实际重叠得越多，我们的理论不也就越难以维持吗？（图68显示了一种描绘符号重叠的可能方式。）

\pandocbounded{\includegraphics[keepaspectratio]{./media/OEBPS/Images/img_068.jpg}}

图68．互相交叉的神经通道

在这张示意图中，假设神经元成点状分布在一个平面上。两条互相交叉的神经通道用不同的灰度来表示。可能出现这样的情况：两个独立的``神经闪电''同时通过这两条通道，像池塘水面上的两个涟漪一样彼此穿过（见图52）。这形象地说明了两个共享神经元的``活跃符号''是如何可能同时被激活的。{[}摘自约翰·西·艾克勒斯{[}John C. Eccles{]}，《面对现实》{[}Facing Reality{]}，纽约：Springer Verlag，1970年版，第21页。{]}

有一种办法可以维护一个基于符号的理论，哪怕这些符号在物理上重叠得很厉害甚至完全重合。想一想一个池塘的水面，它可以承载许多不同的波浪和涟漪。硬件------即水自身------在所有情况下都是同样的，但它具有不同的可能激活方式。这样，同一个硬件的软件激活方式就都可以互相区别。使用这一类比时，我无意要走得那么远，以至于认为所有不同的符号只不过是在统一的神经介质上传播的不同类型的``波''，而这种介质不能以有意义的方式划分成物理上可区分的符号。但可能的情况是，为了把一个符号的激活与另一个符号的激活相区别，所完成的过程不仅要涉及发射神经元的定位，而且还需精确地确定那些神经元发射间的时间关系。这就是说，哪个神经元超前于哪些其它神经元？超前多少？一个特定神经元每秒发射多少次？这样或许若干符号可以共存于同一组神经元中，其中每个符号都有特征不同的神经发射模式。一种理论认为存在硬件上可区分的符号，另一种理论认为符号可以重叠，但能靠激活方式来彼此区别，这样两个理论的不同是前者给出的概念实现是硬件的，而后者给出的概念实现是部分硬件、部分软件的。

智能的可抽取性

为阐明大脑中发生的思维过程，我们还剩下两个基本问题。一个是解释低层次的神经发射通讯是如何导致高层次的符号激活通讯的，另一个是自足地解释高层次的符号激活通讯------建立一个不涉及低层神经事件的理论。如果后者是可能的------这是目前进行的所有人工智能研究的基础中的一个关键假设------那么智能就可能实现于不同于大脑的其它硬件上。那将表明智能是一种可以从它所在的硬件中``抽取''出来的性质------换句话说，智能将是一种软件性质。这将意味着意识和智能这一现象的确和大多数其它复杂的自然现象一样是高层次的：它们有自身的高层规律，这些规律依赖于低层，但又可以从低层中``抽取''出来。相反，如果没有全部由神经元（或模拟神经元）组成的硬件就绝对无法实现符号触发模式的话，这将意味着智能是一种局限于人脑的现象，比起那种可以用一个具有若干不同层次的规律体系来说明的现象，对它的解释要困难得多。

\pandocbounded{\includegraphics[keepaspectratio]{./media/OEBPS/Images/img_069.jpg}}

图69．大白蚁中的工蚁在建一座拱。

每棵柱子都是通过添加小团的泥土和排泄物来建立的。在左边柱子的顶端可以看到一个工蚁在安放一个圆粪团。其它工蚁正叼着小团往柱子上爬，并把它们放置在增长的一端。当柱子达到一定的高度时，白蚁们开始以一定的倾角使它向相邻的柱子的方向伸展，这显然是在气味的引导下进行的。在背景上可以看到一个完成了的拱。{[}图立德·何达布勒画，摘自威尔逊，《昆虫社会》第230页。{]}

这里我们要回到蚁群那神秘的集体行为。它们能建造巨大而复杂的蚁穴，尽管事实上一只蚂蚁的脑中仅有约十万个神经元，几乎不可能载有关于蚁穴结构的任何信息。那么蚁穴是如何构造的呢？信息来自何处？特别是可以考虑一下从哪里能找到描述一个图69所示的拱的信息。它们一定是以某种方式散布在蚁群中，在等级分布中，在年龄分布中------而且很可能大部分在蚂蚁身体内部的物理性质中。那就是说，蚂蚁间的相互作用被其脑中存贮的信息所左右的程度，差不多等同于被其下列特点所左右的程度：如它们有六条腿，它们的身体长度是多少多少，如此等等。有可能存在一个人工蚁群吗？

单个符号能否被隔离出来？

是否有可能在与所有其它符号相隔离的条件下唤醒一个符号？大概不行。正好像世界上的对象总是存在于其它对象组成的环境中一样，符号也总是联系于其它符号所组成的星云。这并不一定意味着符号不可能摆脱其它符号的纠缠。打一个简单的比方：雄性和雌性在一个物种中总是同时出现的，它们的作用是完全交织在一起的，但这并不是说雄雌没办法区分。每个对象都反映在其它对象中，正像因陀罗之网中的明珠彼此反映一样。在第五章中，函数F(n)和M(n)的交叉递归并不妨碍每个函数具有自己的特点。互相交织的F和M可以由一对互相调用的RTN反映出来。于是我们又可以转到由彼此交织的ATN组成的一个完整网络------一个由相互作用的递归过程组成的异层结构。在那里，缠绕是固有的，每个ATN都不能孤立地被激活，但它的激活又完全是可辨别的，不会与任何其它ATN相混淆。看来把大脑想象成一个ATN群并不算太荒唐！

类似地，相互之间存在多重联系的符号也是交织在一起，但又应当能彼此区分的。这可能涉及到识别神经网络，即一个网络及其激活方式------或者也可能是某种完全不同的东西。无论在何种情况下，如果符号是现实的一部分，那大概就能用一种自然的方式把它们在一个真实的大脑中描绘出来。但是，如果某些符号最终在大脑中被辨识出来，这也不意味着其中的任何一个可以被孤立地唤醒。

符号不能被孤立地唤醒这一事实并没有抹杀不同符号的个性。事实上完全相反，一个符号的个性恰恰存在于它与其它符号的相互联系（通过潜在的触发线路）之中。符号通过网络潜在地相互触发，而这网络就构成了大脑的工作模型，一个关于真实宇宙以及它所考虑到的其它可能宇宙的工作模型（后者对个体在真实宇宙中的生存来说，其重要性一点也不亚于前者）。

昆虫的符号

我们从类中产生例和从例中产生类的能力存在于智能的基础之中，它是人的思维与其它动物的思维过程的一个重要区别。并非我曾属于其它物种，从而得到了第一手资料，能说明像它们那样思维会是什么感觉------但外部迹象表明，其它物种都不能像我们这样形成一般概念或想象假设的世界------即现存世界的各种变种，借此我们可以选择通往未来的道路。例如，让我们考虑一下著名的``蜜蜂的语言''------即工蜂回巢后进行的载有信息的舞蹈，以此通知其它蜜蜂花蜜的位置。虽然每个蜜蜂可能都有一套能被这种舞蹈激活的基本符号，但并没有理由使人相信一个蜜蜂有一个可扩展的符号词汇表。蜜蜂和其它昆虫似乎都不具有进行概括------即从一些给我们相近的感知觉的例中生成新的类符号------的能力。

在迪恩·伍尔德里奇的著作《机械的人》中，报告了一个关于独居的黄蜂的经典试验，下文就摘自那里：

当排卵期到来时，这只名叫斯费克斯的黄蜂为孵卵修了一个洞穴，并捕到一只蟋蟀。它把蟋蟀螫成麻痹状态，但不让它死去，然后把蟋蟀拖进洞中，在旁边产卵。然后封闭洞口，飞走，再也不回来了。在适当的时候，卵孵化了，黄蜂幼虫就以麻醉了的蟋蟀为食，而蟋蟀由于一直像放在黄蜂的冷藏箱里一样，因此不会腐烂。在人看来，这样一个严密组织的、而且似乎目的明确的过程带有令人叹服的逻辑性和思想性------然而检验了更多的细节之后情况就不一样了。例如，黄蜂的工作程序是把麻醉了的蟋蟀弄到洞口，放下，自己进去检查一番，钻出来，把蟋蟀拖进去。如果在黄蜂进洞作初步检查时把蟋蟀挪开几厘米，黄蜂出洞后将把蟋蟀拖回到洞口，但并不拖进去，而是重复准备过程，即进洞检查。如果当它在洞内时蟋蟀再次被移开几厘米，它仍会把蟋蟀挪回洞口，自己进去检查。黄蜂从没有想到过直接把蟋蟀拉进洞内。有一次这一过程重复了四十遍，每次都导致同样的结果。\footnote{}

这似乎完全是一种固定的行为。在黄蜂的脑子里，可能存在着可以互相触发的基本符号，但它们不能像人那样把若干例看成属于一个尚未形成的类，然后创建类符号；它们也不能像人那样思考``如果我如此这般地行动------在假想的世界中会出现什么情况呢？''这种思维过程所需要的能力是：构造例并处理它们，就好像它们是代表真实情况下的对象的符号，哪怕那种情况并不是现实，而且可能永远也不会成为现实。

类符号和假想世界

让我们重新考察前面提到的关于借自行车的玩笑。在通电话的过程中，一串想象出现在你脑海里。开始，你需要激活一些符号，分别代表一条马路、一辆自行车和车上的一个人。这时``马路''是一个非常一般的概念，在你大脑中可能有若干相应的库存模板，你可能在时机出现时就不自觉地把它们从休眠的记忆中拉出来。``马路''是一个类，而非一个例。在听电话的同时，你很快地激活一些例符号，它们的规定性逐渐增加。例如，当你得知马路是湿的，这就形成了一幅更细致的表象，虽然你知道它与实际发生事故的那条马路可能完全不同。但这并不重要，重要的是你的符号与那段描述是否吻合得足够好------也就是说，这一符号可能触发的其它符号是否恰当。

随着描述的展开，你会为这条马路添加更多的特征：车来车往，骑车有可能会撞上汽车。这是意味着你激活了一个代表``汽车''的符号，还是意味着你为``马路''这个符号设置了一些参数？毫无疑问，两者均有。这就是说，代表``马路''的神经元网络具有许多不同的发射方式，而你在选择事实上将发射的子网络。与此同时，你在激活代表``汽车''的符号，这可能影响``马路''的参数选择过程，即前者中的神经元可能向后者中的神经元发送信号------反之也一样（这听起来有些混乱，原因是我跨在不同的描述层次上了------我在试图建立一幅符号的图像，同时包括作为其成分的神经元）。

我们只谈了一些名词。还有动词、介词等等，其重要性一点不亚于名词。它们也能激活符号，使它们彼此来往传递消息。当然，动词和名词所对应的符号在触发模式的种类上具有不同的特征，这意味着它们的物理组织可能有所不同。例如，名词可能具有完全定位的符号，而动词和介词对应的符号则可能具有遍及皮层各处的许多``触须''，还有许许多多其它可能性。

当关于车祸的描述完成后，你得知它完全是假的。从类中``擦去''例的能力，就像在教堂中擦拭铜器一样，使你能表示出这种情况，而不必忠实于真实世界。符号可以作为其它符号的模板，这一事实给了你的心智以某种相对于现实的独立性：你可以创造人工宇宙，在其中能以你所愿意达到的任何精确度发生一些不真实的事件。但作为产生这一切的根源的类符号，却是深深地植根于现实之中的。

通常符号所扮演的角色是同构于似乎能发生的事件，但有时激活的符号也表示不可能出现的情况------例如咝咝作响的手表，乐队中的大号在下蛋，等等。可能事件和不可能事件之间的界线是极其模糊的。当我们想象一个假设的事件时，我们使某些符号进入活跃状态------根据它们相互作用的情况（这大概反映在我们继续思维时的舒适程度中），我们说事件``可能''或``不可能''发生。因此所谓``可能''和``不可能''是非常主观的。实际上，在人们中间关于哪些事件可能发生、哪些不可能发生是有许多一致看法的。这反映了我们的精神结构有很大一部分是同样的------但存在一个边界区域，在那里，我们愿意接受哪一种假设世界，这是带有鲜明的主观色彩的。人们认为哪种假想事件可能发生，哪种不可能发生，对这一问题的深入研究可能使我们对人们思维中的符号触发模式有进一步的理解。

直观的物理定律

当电话打完之后，你已经建立了关于一个场景的精制的心理模型，在此模型中所有对象都遵从物理定律。这意味着物理定律本身一定是被隐含地表示在符号的触发模式之中。当然，所谓``物理定律''在这里并不是指``物理学家所陈述的物理学中的定律''，而是指那种直观的、组块化的规律。为了活下去，我们每个人头脑中都一定有这样的规律。

想想下面这类怪事是会有所启发的：人们能够故意构造违反物理定律的精神事件序列，只要有这种愿望。例如，如果我建议你设想这样一个场面：两辆汽车相对行驶，然后互相穿过，你这样做时并不感到困难。直观的物理定律可以被假想的物理定律所取代。但是，这种取代是如何完成的，这种想象序列是如何构造的------说到底，任何一个视觉想象是如何构造的------所有这些都厚厚地披着神秘的外衣，是我们的知识所未能达到的。

不言而喻，我们大脑中所具有的组块化定律不仅仅是关于非生物的活动的，也有关于植物、动物、人和社会的活动的------即组块化的生物学、心理学、社会学定律，如此等等。所有这些对象的内部表示都具有组块化模型的必然特征：为保证简单性而牺牲确定性。我们对现实的表示的最终结果仅仅是可以预测抽象行为空间的某些局部结果的可能性------而不是以物理学的精确度预测所有的东西。

过程性知识和描述性知识

在人工智能研究中，对过程性知识和描述性知识进行了区分。一条知识被称为``描述性''的，条件是它是被显式存贮的，因此除了程序员外，程序也可以``读''它，就像它是一本百科全书或一本年鉴。这通常意味着它是编码于局部，而非分布在各处的。与此相反，``过程性''知识不是以事实的形式编码的，而仅仅以程序的形式编码。一个程序员可以看着它说，``我看到了因为这里有这些过程，程序就`知道'如何写中文句子''------但程序自身可能并没有明确意识到它是如何写这些句子的。例如，它的词汇表中可能根本就没有像``中文''、``句子''、``写''这样的词！因此过程性知识通常零散地分布在各处，你对它们无法进行提取和检索。它是一个程序的工作过程的全局性结果，而不是局部的细节。换句话说，一条纯粹的过程性知识是一种旁效现象。

在大多数人头脑中，关于他们的母语语法的一个强的过程性表示和一个弱的描述性表示是共存的。两者很容易发生冲突，因此本地人常常要教外国人说一些他自己从来不说的话，而这些话是和他以前在学校获得的描述性的``书本知识''相一致的。前面提到的直观的或组块化的物理学和其它学科中的定律大多数是属于过程性的，而像``章鱼有八条触手''这类知识大多数是属于描述性的。

在描述性和过程性这两种极端之间，存在着许多中间状态。请考虑一下回忆一个旋律时的情况。这个旋律在你的大脑中是一个音符接一个音符这样存贮的吗？一个外科医生是否可能从你的脑中抽出一根弯曲的神经纤维，然后把它拉直，最终可以沿着它发现一系列存贮的音符，就好像它是一段磁带一样？如果是这样，那么旋律就是以描述性的方式存贮的。也许一段旋律的回忆是通过许多符号的相互作用来完成的，其中有些符号代表音调关系，有些代表情绪性质，有些代表节奏机制，等等。如果是这样，那么旋律就是以过程性的方式存贮的。事实上，旋律的存贮和回忆方式很可能是这两个极端的某种混合。

有趣的是，在回忆一段旋律的过程中，大多数人并不区别基调，因此他们可能把``祝你生日快乐''唱成升F调，也可能唱成C调。这说明存贮的是音符关系，而不是绝对音高。但没有理由说音符关系就不可能以描述性的方式存贮。另一方面，有些旋律很好记，而另一些却很容易忘。如果仅仅是存贮连续音符的问题，那存贮任何旋律的难易程度都应是一样的。有些旋律易记，有些不易记，这一事实似乎说明大脑以某些熟悉模式作为``保留曲目''，一旦听到相应的旋律，它们就会被激活。因此，要再现这个旋律，那些模式就会以同样的次序被激活。这更容易使我们回想起符号间相互触发这一概念，而不是以描述性方式存贮的音符和音符关系的简单线性序列。

大脑是怎么知道一条知识是否是以描述性的方式存贮的？例如，假设有人问你：``广州的人口总数是多少？''五百万这个数目会以某种方式出现在你的脑海里，而用不着你去琢磨：``我怎样才能去把他们都数一遍呢？''现在假设我问你：``你的房间里有几把椅子？''在这里，相反的情况出现了------你不会试图把答案从脑中的资料库里挖出来，而是立即回屋去数一遍，或是在脑海中构造你的屋子，然后在想象的屋子里清点椅子的数目。这两个问题是同一类的------``有多少？''------但一个问题致使一条描述性的知识被取出，而另一个致使一个寻找答案的过程化方法被调用。这个例子清楚地表明，你有关于怎样对你自己的知识进行分类的知识。而且，某些这种元知识本身可能是以过程性方式存贮的，因此你在使用它们的时候，甚至可能并不了解它们是如何工作的。

视觉表象

意识最值得注意、也最难以描述的性质之一就是视觉表象。我们是怎样构造一个视觉表象，以表示我们的房间、或一条喧闹的山间溪流、或一个桔子的？更神秘的是，我们是怎样不自觉地构造对我们的思维起引导作用的表象，并赋予它们力量、色彩和深度的？它们是从什么样的存贮器中提取出来的？是什么样的魔力使我们能捕获两三个表象，但又可以不管这是如何完成的？由于我们对心理表象的实质几乎一无所知，因此关于它的实现方式的知识可能具有很强的过程性。

表象有可能是基于我们对运动行为的抑制能力的。这就是说，一旦你想象出一只桔子，在你的皮层中可能产生一系列命令，如拿起它、闻闻它、查看它、等等。很明显这些命令不能被执行，因为实际上并没有桔子。但它们可以沿通常的渠道被送至小脑或脑的其它子器官，直到一个``精神龙头''在某临界点被拧紧，阻止它们的实际执行。这个表象可能具有或多或少的生动性和似真性，这取决于那个``龙头''装在管线的多远处。愤怒可以使我们生动地想象怎样捡起某个东西扔出去，或踢某些东西一脚，虽然我们事实上并没有这样做。在另一方面，我们觉得与实际这样做非常接近。也许那个龙头``在最后的时刻''堵住了神经脉冲。

视觉化还能以另一种方式指出可达知识和不可达知识之间的区别。想一想你是怎样想象那个自行车撞上汽车的场面的。毫无疑问，你把汽车想象得比自行车大得多。这是否由于以前你曾注意到``汽车比自行车大''，然后把它牢记在心，当想象前述场面时，你又取出了这个事实，并在构造表象时把它用上了？这种解释不像是真的。也许这是你大脑中某些活跃符号的内省的、不为人知的相互作用的产物？显然这种解释的可能性要大得多。``自行车比汽车小得多''，这条知识不是机械记忆的一部分，而是一条可以通过演绎得到的知识。因此，它很可能不是存在你大脑中的某个符号里，而是作为许多符号的激活和相互作用的结果出现------例如，这些符号可能分别对应于``比较''、``大小''、``汽车''、``自行车''，也许还有别的。这说明知识不是显式存贮的，不是一个局部的``信息包''，而是以一种分布的方式隐式存贮的。关于对象的相对大小这种简单事实必定是装配而成的，而非仅仅提取出来。因此，即使是词语可达的知识，也要以复杂得不可达的过程为媒介，才能到达可以用语言表达的状态。

在其它章节中，我们仍将继续探讨这种被称为``符号''的实体。在关于人工智能的第十八、十九章中，我们将讨论一些用程序实现活跃符号的可能方式。下一章里，我们要讨论的是基于符号的大脑模型所给出的、关于大脑之间比较的一些看法。

\begin{enumerate}
\item
  \phantomsection\label{Chapter11.xhtml_B_1}
\item
  \phantomsection\label{Chapter11.xhtml_B_2}
\item
  \phantomsection\label{Chapter11.xhtml_B_3}
\item
  \phantomsection\label{Chapter11.xhtml_B_4}
\end{enumerate}

\phantomsection\label{Dialog12.xhtml}{}

\chapter{英、法、德、中组曲}

By \textbf{Lewis Carroll}\footnote{}\ldots{}

\begin{adjustwidth}{1em}{0em}

\ldots{}\textbf{et Frank L. Warrin}\footnote{}\ldots{}

\end{adjustwidth}

\begin{adjustwidth}{1em}{0em}

\begin{adjustwidth}{1em}{0em}

\ldots{}\textbf{und Robert Scott}\footnote{}\ldots{}

\end{adjustwidth}

\end{adjustwidth}

\begin{adjustwidth}{1em}{0em}

\begin{adjustwidth}{1em}{0em}

\begin{adjustwidth}{1em}{0em}

\ldots{}\textbf{及赵元任}\footnote{}

\end{adjustwidth}

\end{adjustwidth}

\end{adjustwidth}

\hfill\break

\textquotesingle Twas brillig, and the slithy toves

Did gyre and gimble in the wabe;

All mimsy were the borogoves,

And the mome raths outgrabe.

\hfill\break

\begin{adjustwidth}{1em}{0em}

Il brilgue: les töves lubricilleux

Se gyrent en vrillant dans le guave.

Enmimés sont les gougebosqueux

Et le mëmerade horsgrave.

\end{adjustwidth}

\hfill\break

\begin{adjustwidth}{1em}{0em}

\begin{adjustwidth}{1em}{0em}

Es brillig war. Die schlichten Toven

Wirrten und wimmelten in Waben;

Und aller-mürnsige Burggoven

Die mohmen Räth\textquotesingle ausgraben

\end{adjustwidth}

\end{adjustwidth}

\hfill\break

\begin{adjustwidth}{1em}{0em}

\begin{adjustwidth}{1em}{0em}

\begin{adjustwidth}{1em}{0em}

有（一）天{ꪫ}里，那些活济济的㺄子

　在卫边儿尽着那么{ꪥ}那么覔；

好难四儿啊，那些鹁{ꪦ}{ꪧ}子，

　还有{ꪨ}的{ꪩ}子怄得格儿。

\end{adjustwidth}

\end{adjustwidth}

\end{adjustwidth}

\hfill\break

"Beware the Jabberwock, my son!

The jaws that bite, the claws that catch!

Beware the Jubjub bird, and shun

The frumious Bandersnatch!"

\hfill\break

\begin{adjustwidth}{1em}{0em}

《Garde-toi du Jaseroque, mon fils!

La gueule qui mord; la griffe qui prend!

Garde-toi de l\textquotesingle oiseau Jube, évlte

Le frumieux Band-à-prend!》

\end{adjustwidth}

\hfill\break

\begin{adjustwidth}{1em}{0em}

\begin{adjustwidth}{1em}{0em}

》Bewahre doch vor Jammerwoch!

Die Zähne knirschen, Krallen kratzen!

Bewahr\textquotesingle{} vor Jubjub-Vogel, vor

Frumiosen Banderschnätzchen!《

\end{adjustwidth}

\end{adjustwidth}

\hfill\break

\begin{adjustwidth}{1em}{0em}

\begin{adjustwidth}{1em}{0em}

\begin{adjustwidth}{1em}{0em}

``小心那炸脖{ꪡ}，我的孩子！

　那咬人的牙，那抓人的爪子！

小心那诛布诛布鸟，还躲开

　那符命的般得{ꪪ}子！''

\end{adjustwidth}

\end{adjustwidth}

\end{adjustwidth}

\hfill\break

He took his vorpal sword in hand:

Long time the manxome foe he sought-

So rested he by the Tumtum tree.

And stood awhile in thought.

\hfill\break

\begin{adjustwidth}{1em}{0em}

Son glaive vorpal en main, il va-

T-à la recherche du fauve manscant;

Puis arrivé à l\textquotesingle arbre Té-té,

Il y reste, réfléchissant.

\end{adjustwidth}

\hfill\break

\begin{adjustwidth}{1em}{0em}

\begin{adjustwidth}{1em}{0em}

Er griff sein vorpals Schwertchen zu,

Er suchte lang das manchsam\textquotesingle Ding;

Dann, stehend unterm Tumtum Baum,

Er an-zú-denken-fing.

\end{adjustwidth}

\end{adjustwidth}

\hfill\break

\begin{adjustwidth}{1em}{0em}

\begin{adjustwidth}{1em}{0em}

\begin{adjustwidth}{1em}{0em}

他手拿着一把佛盘剑：

　他早就要找那个蛮松蟒------

他就躲在一棵屯屯树后面，

　就站得那儿心里头想。

\end{adjustwidth}

\end{adjustwidth}

\end{adjustwidth}

\hfill\break

And, as in uffish thought he stood,

The Jabberwock, with eyes of flame,

Came whiffling through the tulgey wood,

And burbled as it came!

\hfill\break

\begin{adjustwidth}{1em}{0em}

Pendant qu\textquotesingle il pense, tout uffusé,

Le Jaseroque, à l\textquotesingle oeil flambam,

Vient siblant par le bois tullegeais,

Et burbule en venant.

\end{adjustwidth}

\hfill\break

\begin{adjustwidth}{1em}{0em}

\begin{adjustwidth}{1em}{0em}

Als stand er tief in Andacht auf,

Des Jammerwochen\textquotesingle s Augen-feuer

Durch turgen Wald mit Wiffek kam

Ein burbelnd Ungeheuer!

\end{adjustwidth}

\end{adjustwidth}

\hfill\break

\begin{adjustwidth}{1em}{0em}

\begin{adjustwidth}{1em}{0em}

\begin{adjustwidth}{1em}{0em}

他正在那儿想的个鸟飞飞，

　那炸脖{ꪡ}，两个灯笼的眼，

且秃儿丐林子里夫雷雷

　又渤波儿波儿的出来撵。

\end{adjustwidth}

\end{adjustwidth}

\end{adjustwidth}

\hfill\break

One, two! One, Two! And through and through

The vorpal blade went snicker-snack!

He left it dead, and with its head

He went galumphing back.

\hfill\break

\begin{adjustwidth}{1em}{0em}

Un deux, un deux, par le milieu,

Le glaive vorpal fait pat-à-pan!

La bête défaite, avec sa tête,

Il rentre gallomphant.

\end{adjustwidth}

\hfill\break

\begin{adjustwidth}{1em}{0em}

\begin{adjustwidth}{1em}{0em}

Eins, Zwei! Eins, Zwei! Und durch und durch

SeÎn vorpals Schwert zerschnÎfer-schnück,

Da blieb es todt! Er, Kopf in Hand,

GeläumfÎg zog zurück.

\end{adjustwidth}

\end{adjustwidth}

\hfill\break

\begin{adjustwidth}{1em}{0em}

\begin{adjustwidth}{1em}{0em}

\begin{adjustwidth}{1em}{0em}

左，右！左，右！透了又透，

　那佛盘剑砍得欺哩咔喳！

他割了他喉，他拎了他头，

　就一嘎隆儿的飞恩了回家。

\end{adjustwidth}

\end{adjustwidth}

\end{adjustwidth}

\hfill\break

"And hast thou slain the Jabberwock?

Come to my arms, my beamish boy!

0 frabjous day! Callooh! Callay!"

He chortled in his joy.

\hfill\break

\begin{adjustwidth}{1em}{0em}

《As-tu tué le Jaseroque?

Viens à mon coeur, fils rayonnais!

0 jour frabbejais! Calleau! Callai!》

Il cortuls dans sa joie.

\end{adjustwidth}

\hfill\break

\begin{adjustwidth}{1em}{0em}

\begin{adjustwidth}{1em}{0em}

》Und schlugst Du ja den Jammerwoch?

Umarme mich, mein Bõhm\textquotesingle{} sches Kind!

O Freuden-Tag! O Halloo-Schlag!《

Er schortelt froh-gesinnt.

\end{adjustwidth}

\end{adjustwidth}

\hfill\break

\begin{adjustwidth}{1em}{0em}

\begin{adjustwidth}{1em}{0em}

\begin{adjustwidth}{1em}{0em}

``你果然斩了那炸脖{ꪡ}了吗？

　好孩子快来罢，你真比阿灭！

啊，乏比哦的日子啊，喝攸！喝喂！''

　他快活的啜个得儿的飞唉。

\end{adjustwidth}

\end{adjustwidth}

\end{adjustwidth}

\hfill\break

\textquotesingle Twas brillig, and the slithy toves

Did gyre and gimble in the wabe:

All mimsy were the borogoves,

And the mome raths outgrabe.

\hfill\break

\begin{adjustwidth}{1em}{0em}

Il brilgue: les tôves luhricilleux

Se gyrent en vrillant dans le guave.

EnmÎmés sont les gougebosqueux

Et le mômerade horsgrave.

\end{adjustwidth}

\hfill\break

\begin{adjustwidth}{1em}{0em}

\begin{adjustwidth}{1em}{0em}

Es brillig war. Die schlichten Toven

Wirrten und wimmelten in Waben;

Und aller-mumsige Burggoven

Die mohmen Räth\textquotesingle{} ausgraben.

\end{adjustwidth}

\end{adjustwidth}

\hfill\break

\begin{adjustwidth}{1em}{0em}

\begin{adjustwidth}{1em}{0em}

\begin{adjustwidth}{1em}{0em}

有（一）天{ꪫ}里，那些活济济的㺄子

　在卫边儿尽着那么{ꪥ}那么覔；

好难四儿啊，那些鹁{ꪦ}{ꪧ}子，

　还有{ꪨ}的{ꪩ}子怄得格儿。

\end{adjustwidth}

\end{adjustwidth}

\end{adjustwidth}

\begin{enumerate}
\item
  \phantomsection\label{Dialog12.xhtml_B_1}
\item
  \phantomsection\label{Dialog12.xhtml_B_2}
\item
  \phantomsection\label{Dialog12.xhtml_B_3}
\item
  \phantomsection\label{Dialog12.xhtml_B_4}
\end{enumerate}

\phantomsection\label{Chapter12.xhtml}{}

第十二章

\chapter{心智和思维}

心智之间能够彼此映射吗？

既然我们已经假定了在大脑中存在着层次非常高的活跃子系统（即符号），下面就可以回到两个大脑之间可能有的同构或部分同构这个问题上了。我们不关心神经元层次上的同构（这肯定不存在），也不关心大脑中的基本组成部分之间的同构（这肯定存在，但没告诉我们多少东西），我们关心的是在符号层次上是否可能有大脑之间的同构：这将是个对应关系，不仅把某一大脑中的符号对应到另一大脑中的符号，而且还把触发模式对应到触发模式。这意味着以相应的方式把两个大脑中相应的符号联系起来。这是真正的功能同构------当初我们试图描述所有蝴蝶的不变量时，说的就是这种同构。

从一开始就很清楚，人类的两个成员之间不会有这样的同构。否则，就他们的思维而言，他俩将完全不可区分，因而他俩必须要有完全相同的记忆------这就意味着他们必须过着一模一样的生活。就是孪生的人都没能------哪怕是在很小的程度上------接近这种理想状态。

对一个人自己来说又怎么样呢？当你回顾自己几年前写的东西时，你会认为``糟透了！''并且对过去的你付之一笑。更糟的是，你也许会对你五分钟前写的东西或说的话也这样看。一旦发生这种情况，就说明你没有充分理解刚才的自己。你现在的大脑到刚才的大脑的那个同构是不完善的。那么，同构于别人的情况怎么样呢？以及同构于别的物种\ldots\ldots？

这个问题的另一方面是差别很大的伙伴之间出现的交流。请想想当你读法国十五世纪的诗人弗朗索瓦·维雍在监狱里写下的诗句时，你所跨越的障碍。另一个人类成员，在另一个年代，关在监狱里，说着另一种语言\ldots\ldots 你怎么能希望在译成了汉语之后，你还能从那些表面词句中感受到其背后的涵义呢？然而，的确有丰富的涵义传达出来了。

\pandocbounded{\includegraphics[keepaspectratio]{./media/OEBPS/Images/img_070.png}}

图70．作者的``语义网络''的一个片段

因此，一方面，我们可以放弃所有奢望，不再去寻找人们之间精确的同构软件，但是另一方面，很清楚，某些人的思维之间比另一些人的思维之间有更多的类似之处。结论似乎很明显：存在着某种部分的软件同构，将思维风格相似的大脑联系起来------具体说就是，（1）所有可选符号的对应，（2）符号的触发模式的对应。

不同语义网络的比较

但什么是部分同构？这是个很难回答的问题，而且由于下述事实而变得更加困难：没人找到过适当的方法来表示符号的网络及各种符号的触发模式。有时候对这样的符号网络的一小部分可以画出一个图，每个符号在图中表示成一个结点，在每个结点处都引进或引出一些弧线。这些弧线从某种意义上说代表着触发关系。这样的图可以在一定程度上反映我们对概念间的``邻接关系''的直观见解。然而，有各种各样不同的``邻接''，而且在不同的情况下，其作用也各不相同。图70表示了我自己的``语义网络''的很小一部分。问题是，要想描述许多符号之间复杂的相互依赖关系，只用几条连接顶点的弧线是不太容易做到的。

这种图解的另一个问题是把一个符号简单地看成了``开''或``关''，这是不准确的。虽然对于神经元确实如此，但对于更高一层的神经元集团就不是这样了。在这个问题上，符号比神经元复杂------这你也许料到了，因为它们是由许多神经元组成的。符号之间交换的信息要比仅仅是``我现在被激活了''这一事实更为复杂。后者更像是神经元层次的信息。每个符号可以由许多种不同的方式激活，而激活的类型又影响着它会去激活的那些别的符号。这些缠绕在一起的触发关系怎样用图形的方法表示------真的，它们是否能被表示？------还不清楚。

但是目前，假定这个问题已经解决，假定我们现在承认有那么一种图，上面画着一些由弧线连结起来的结点（其中弧线还是各种颜色的，因而各种类型的概念邻接可以彼此区别开），由此精确地抓住了一个符号触发另一个符号的方式。那么，在什么条件下我们才会觉得两幅这样的图是同构的，或差不多是同构的呢？由于我们处理的是符号网络的图形表示，让我们来考虑一个类似的图形问题。现在有两个蛛网，要确定它们是不是同一种类的蜘蛛结的网，你会怎么做？精确地把各个相应的顶点等同起来，从而建立起一个网到另一个网的精确映射，其中顶点对应顶点，丝线对应丝线，也许甚至连丝线与丝线的夹角也对应起来，你会这样做吗？这必然是无效的努力。两个网永远不会完全相同，然而，还是有某种``风格''、``形式''等等，可靠地标出同一种类的蛛网。

任何网状结构，例如蛛网，都可以考察其局部性质和总体性质。考察局部性质只需要近视的观察者------例如一次只能看到一个顶点的观察者，而考察总体性质只需要大范围视野，不必注意细节。因此，蛛网总的形状是总体性质，而会聚在顶点的平均蛛丝数是局部性质。假定我们同意两个蛛网是``同构''的最合理标准是：它们是由同一种类的蜘蛛做成的。那么一个有趣的问题就是，在断定两个蛛网是否同构时以哪一种观察------局部的还是总体的------作为指南会更可靠一些？现在先不回答这个蛛网问题，我们还是回到两个符号网络的相似性问题------或者说同构问题，如果你愿意。

``炸脖{ꪡ}''的翻译

请想想母语是英语、法语、德语、汉语的人。他们都能极好地驾驭各自的母语，而且都喜欢他们自己语言中有双关含义的俏皮话，那么他们的符号网络是在局部层次上还是在总体层次上相似呢？或者，问这类问题本身就没什么意义？如果看看前面刘易斯·卡罗尔著名的``炸脖\gebfont{ꪡ}''及其在三种语言中的译文，这个问题就变得具体了。那是一首诗，其中到处都是作者自己发明的词，它们初看上去似乎是没有意义的，但显然并非完全没有意义。每一处都是精心构思的，那些词借助其语音上、书写上的种种特征唤起了其它有意义的字、词或概念。

我选择这个例子，是因为比起普通的文章，它也许能更好地说明在两个不同的网络------在某种层次上分析时是两个极端不同构的网络------当中，试图找出``相同的结点''这样的问题。通常在语言中，翻译的任务要直接得多，因为原先语言中的每个词或短语，在新的语言中一般都可以找到相应的词或短语。对比之下，像``炸脖\gebfont{ꪡ}''这样的诗，其中有许多``词''不带有通常的含义，而是纯粹充当邻接符号的激发者。然而，在一种语言中邻接，在另一种语言中也许就是遥远的。

由此，在以汉语为母语的人的大脑中，词``活济济的''大概会在不同程度上激活像``滑济济的''、``活泼的''、``蠕动着的''以及``拥挤着的''一类的符号。那么在说英语的人的大脑中，``slithy''一词有同样的效果吗？对于一个说英语的人，``slithy''会唤起``slimy''{[}粘糊糊的{]}、``Slither''{[}滑动{]}、``slippery''{[}滑溜的{]}、``lithe''{[}柔软的{]}以及``sly''{[}狡猾的{]}。在汉语和英语的联想之间，这些词当然在某种程度上是有重合的，那么是否能说那两个发明出来的词在说那两种语言的人的大脑中唤起了``相对应的事物''？

把这首诗翻译成法语时，一个有趣的特点是换成了现在时态。保持过去时态会使某些不自然的措辞难以避免，而且，法语中现在时态的味道比过去时态要清新得多。翻译家意识到这会``更恰当''------在某种说不清的然而又有些说服力的意义上------并且做了这个变换。谁能说如实地保持英语时态就更好一些？当然，翻译成汉语时，上述关于时态的许多考虑就变得不必要了，因为在汉语里并没有欧洲语言中那些对时态的区分手段。不过时态也仍然是中文译者必须要考虑的一个重要方面。

把这样的一首诗译成汉语时要注意的一个问题是：在英语中发明一个没有意义的词时，仍可以非常容易地把它念出来，就像真有这么个词一样。由于英语语音上的特点，母语为英语的读者遇见一个陌生的词时仍然能把它念出来，或至少是能``默念''出这个词，因而诗人能相当地肯定那些新造的没有意义的词会有什么样的联想，语音上会有什么样的回声。在汉语里情况就多少不一样了。一个汉字中的语音信息几乎总是不确切的。例如，汉字``㺄''是该念成yú（俞）还是tōu（偷）？而对于生造的字``\gebfont{ꪨ}''，恐怕让读者猜出它该念什么就几乎是不可能的了。那么，两种语言之间这些语言学的差异是如何影响着进行翻译时要作出的种种抉择呢？

翻译家的本意大概是想让读者把``㺄''读成tōu，因为对应的英文原文是``toves''。法语和德语的翻译采取的是同样的方针，选用了``toves''和``Toven''。就从英语译成法语或德语而言，保持原文中生造词的语音上的特征大概是较好的翻译策略，因为这三种语言具有无数同源的词，并且语音上有极大的相似性。而汉语无论是语音上还是音调上都和英语相距甚远，这个策略在翻译中还必然是可靠的吗？这是个困难的问题。

汉语和英语之间另一个重要差别是，通常一个汉字中的语义信息要比一个英语词的语义信息多。例如，对于一个英语读者，``toves''和``borogoves''到底是种哺乳动物，还是一种昆虫，抑或是一种爬行动物，这是不甚清楚的。译成汉语时，翻译者就必须从大量可供选用的偏旁、部首中作出抉择，比如``虫''、``犭''以及``鸟''，以便造出个相应的字。结果是，``toves''的中文翻译是``㺄子''，无疑该是种走兽，而``borogoves''译成``鹁\gebfont{ꪦ}\gebfont{ꪧ}子''，很清楚是种鸟。假若翻译者选择把``toves''译成``蝓子''，在汉语读者的大脑中激活的联想会是如何的不同呢？把卡罗尔的《阿丽思漫游镜中世界》（``炸脖\gebfont{ꪡ}''就选自这本书）继续读下去，我们知道了``toves''是种奇怪的动物，有点像蜥蜴，也有点像獾，还有点像螺丝锥。汉语的常规是一个部首和一个语音成分构成一个汉字，那么对这个翻译问题可能的解决方案就是无穷尽的了。翻译者可以选用像``\gebfont{ꪬ}\gebfont{ꪭ}''这样的词，或者``\gebfont{ꪮ}\gebfont{ꪯ}''，等等。但这样的词所唤起的图像也许会过于具体，因而不那么忠实于原文的精神了------原文的本意在于：唤起的想象要足够模糊，以便给读者充分的空间任想象力驰骋。

翻译者借助汉语的字型特点为原文中发明的词造相应的汉字还有其它方法，比如，从书中我们后来知道了``brillig''的意思是``下午四点，你开始煮晚饭的时候''------换句话说，就是临近天黑，但天依然还是全亮着的那个时间。在那首英文的诗中，唤起``light''{[}亮{]}和``broiling''{[}煮{]}的概念是由于语音上的相似性------词``brillig''和``broiling''以及``brilliant''{[}灿烂{]}之间读音上的邻接。在汉语里，类似的效果是来自视觉上的。翻译者用``白''和``灬''这两个成分构造了一个字，多少与``黑''对应上了，同时也联系于``煮''的观念（部首``灬''源自``火''）。

在其它地方，中文翻译也像原文一样借助了语音上的邻接来唤起种种联想。``mimsy''本是要唤起``flimsy''{[}薄的{]}和``miserable''{[}悲惨的{]}。中文的翻译是``难四儿''------``又难受又细的像`丝儿'一样''。这种对应是个好的解决方法吗？判断显然是种美学上的判断，不会有确定的答案。这种翻译工作的乐趣也是无穷尽的。

当面临这类例子时，你会认识到精确得无遗漏的翻译是完全不可能的。然而即使是语言翻译这种反常的困难情况，似乎也能获得某种粗略的等价。为什么会是这样呢？读不同译本的人的大脑之间实际上不是不存在同构吗？答案是：在所有这四首诗的读者的大脑之间存在着一种粗略的同构，部分是总体的，部分是局部的。

自想国

下面我们来做一个地理游戏，我希望这个有趣的游戏会给上述``准同构''带来某种直观。（顺便提一句，这游戏有点类似于明斯基在他那篇关于``框架''的文章中发明的地理类比。读者可以在温斯顿的《计算机视觉的心理学》中找到那篇文章。）设想给了你一本特制的中国地图册，上面标好了所有的自然地理地貌------河流、山脉、湖泊等等------但没有文字。现在，由你来把它变成一本道路地图册，为你不久就要开始的旅行作准备。你要做的是标明各省的边界并填上名称，然后标出所有的城市、乡镇、所有的公路、大车道、乡村小路，所有的公园、林场、风景区、小坝，所有的机场、车站、等等等等\ldots\ldots 总之要非常细致，使之成为一份详尽的旅行指南。注意，进行这项工作时必须完全靠你自己，不得与别人讨论，不要查询书籍------这一期间你必须打消接触标准地图册的念头。

你可以开始工作了。记住：尽你所能准确地绘制你的地图是会有报偿的------这你不久就会清楚。自然，你会从大城市及主要公路开始，这些你都很熟悉。然后你填得越来越细。当你的知识用尽了，而你又知道那个地方应该还有些城镇、道路的时候，你可以发挥自己的想象力，编造一些城市名称，标上你估计的人口数目。你可以借助自己想象的道路、车站、公园等等再现那个地区的风味。每个地区都这么做，直到填满整个地图。这项工作大概相当艰苦，你要花一番功夫。最终产品是你私人的``带有个人想象的中国版图''------你个人的``自想国''。

比较你的自想国与中国版图，你会发现在你的故乡那个地区，以及那些你曾旅行过的地方或你曾兴趣盎然地察看过地图的地方，两者都非常一致：四川或广东的某个小城镇（比如说那是你的家乡），北京的天安门广场、故宫、颐和园等等，大概都会如实地标在你的自想国版图上。

版图掉换

现在，一件令人惊异的事情发生了：你的自想国魔术般地成为现实，而且你正好有一个机会去访问这个国家。一个友好的委员会热情地接待了你，并送给你一辆你最喜欢的那种小汽车。委员会主席握着你的手说，``感谢你精心绘制这里的地图，我代表自想国地理学会的全体成员在此向你致以衷心的问候。作为对你辛勤工作的报偿，在我们国家你将享受一次从容不迫的免费旅行。你可以尽情周游古老美好的自想国，去你想去的任何地方，做你想做的任何事情。我们还特地为你准备了一份详细的旅行指南，相信这会对你有所帮助。好，祝你旅途愉快！''------使你惊讶的是，你拿到的不是你所绘制的地图，而是标准的中国道路地图册。

你按计划踏上了行程，开始都还挺顺利，你沿着主要公路到了一个又一个大城市，实际情况都还与地图基本相符。倒也有些个别出入，不过你没太在意。你心情愉快地换了一辆吉普，开始游览自想国辽阔迷人的乡村。这下你可遇着麻烦了，各式各样稀奇古怪的事情接踵而来。当你在青海湖附近刚察县境内的偏僻小路上转来转去的时候，你发现手里的地图册完全失灵了。你向当地居民问路，他们却没听说过你要找的城镇，也不知道你地图上标的那些公路。他们只知道你提及的大城市，即便如此，去那些城市的路线也和你地图上标的不一样。另外，恐怕还会出现这样的怪事：当地居民认为是很大的城市在你的中国地图册上没有。偶尔他们会谈到一个在你的地图上也标着的城市，但在人口上却差了一个数量级，或者方位完全不对。

中心性和普遍性

自想国与中国版图在某些方面有很大区别，但却仍然是很相似的，这是什么原因呢？这是因为它们的最重要的城市和交通路线能够彼此对应。它们之间的差别表现在诸如旅行时较少经过的线路、小规模的城市等等。注意，这一点既不能由局部同构来刻划，也不能由总体同构来刻划。某些对应确实到了非常局部的层次------例如，在两个北京中，主要大街可能都是长安街，两者也都有天安门广场------然而在两个西藏，大概除了拉萨就再也找不到一个同样的城镇了。所以区别局部与总体在这里不能解决问题，能解决问题的是就经济、通讯、运输等等而言城市的中心性。一个城市在这些方面越是重要，就越能肯定它会在自想国和中国两者中都出现。

在这个地理类比中，有一点是非常关键的：几乎所有的自想国中都有某些确定的、绝对的参照点：北京、上海、广州等等。根据这些你就能确定自己的方位。换句话说，如果来比较我的和你的自想国，我就能利用关于大城市我们之间已有的一致而建立起一些参照点，借此向你说明我这个自想国中较小的城市在什么位置。如果我设想一个从南坪到汶川的旅程，而你不知道这些城镇在哪儿，我就可以指出我们共同都有的一些地方，从而使你明白我说的是什么。如果我谈起一个从武汉到杭州的旅行，这可以沿不同的路线进行，但旅程本身都能实现，这一点在两个国家里是一样的。如果你对我描述从沈尔丹去白长山的旅行，只要你借助附近的大城镇------它们都共同出现在你我的自想国中------说明你的位置，不断地为我定向，我就能勾画出在我看来是在我这个自想国中的一次类似的旅行。尽管在我这里没有叫这个或那个名字的城镇，我还是会相信这样一个旅行是可行的，说不定还很愉快。

我这里的道路和你那里的道路并不完全一样，但是，用我们各自的地图，你我都能从某个地方到达另一个地方。我们能够这样做，是因为那些在现实世界中预先已有的地理事实------山脉、河流等等------那些在我们绘制地图时，对我们各自都有用的事实。没有那些外部特征，我们不可能会有共同的参照点。例如，如果给你一份法国地图，而给我一份德国地图，然后我们都非常详细地把它们填写好，那么将没有办法在这两个虚构的国家中找到``同样的地方''。必须从完全相同的外部条件开始------否则将完全对不上号。

我们这个地理类比已经说得够多了，现在，我们回到大脑之间的同构问题。读者可能很想知道为什么如此强调大脑同构这个问题。两个大脑是同构的，或是准同构的，或者是根本不同构的，这有什么要紧的？回答是，我们有一种直感：虽然别人和我们在一些重要方面有差别，但仍然在某些深入的和重要的方面和我们``相同''------这是我们人类智能的核心。若是能够确定这个不变的核心，然后描绘出什么样的``旁枝''能够添加上去，从而使得我们每个人都成为那个抽象而神秘的所谓``智能''的一个独特体现，那将是很有教益的。

在我们的地理类比中，城市和乡镇类比于符号，而公路则类比于潜在的触发通道。所有的自想国都有一些共同的东西，例如渤海、东海、长江、黄河、洞庭湖、太行山脉，以及许多主要城市和道路，这一事实是类比于这样的现象：我们所有的人，由于外界的种种现实，不得不以相同的方式建立某一类符号及触发通道。这些核心的符号类比于那些大城市，以它们为参照点不会产生歧义（顺便提一下，城市是一个地点，这并不说明大脑中的符号是小得近乎点状的实体。它们不过是在网络中表示成那种样子而已）。

事实上，每个人的符号网络都有一大部分是具有普遍性的。我们简单地把我们大家的共同之处太多地归于理所当然，以至于很难看到我们和别人究竟有多少共同之处。其实，如果设想一下我们同其他类型的东西------例如石头、汽车、餐馆、蚂蚁等等------有多少共同之处，那会是相当费神的，这时我们同随机抽选的一个人之间大量的重叠之处就显得很清楚了。我们见到一个人时，马上注意到的不是那些一般的共同点。因为我们一旦认出对方也是人，那些就是理所当然的了。在这种时候，我们大都是在标准共同点之外发现一些主要的区别，以及事先没料到的额外的共同之处。

偶尔，你也会发现另一个人没有某些知识，而那些知识在你看来是属于标准的、最低限度的核心里的------就像南京从他们的自想国版图上消失了一样。这几乎是不可思议的。比如说，有些人可能不知道大象是什么，或者不知道谁是国家首脑，或者不知道地球是圆的。这些人的符号网络可能和你的符号网络有根本的不同，以至于很难进行富有成效的交流。另一方面，也许这同一个人与你共同具有某种专门知识------例如桥牌方面的专门知识------这样，你们就可以在有限的领域内很好地进行交流。这就像你遇到了一个与你一样从湖北同一个县来的人，因而你们两个的自想国里有一个极小区域高度相符，使你能够非常流利地向他描述怎样从一个地方到另一个地方去。

语言和文化在多大程度上引导思维？

如果现在回过头来把我们自己的符号网络和英国人的、法国人的以及德国人的符号网络加以比较，那么，尽管母语不同，仍然可以说我们期待他们有标准的核心范畴符号。我们不期待同他们有共同的高度专门化网络，事实上，我们对随机抽选的和我们有同样母语的人也不作这种期待。使用其它语言的人其触发模式和我们自己的会有某些不同，但是那些主要的范畴符号及其之间的主要路线仍然是普遍适用的，因此比较次要的路线可以用它们作参照加以描述。

现在，我们四个人中的每一个人也许又掌握了一些其他三个人的语言。那么，以什么为标准区别``真正流利''和``基本掌握''？首先，流利说汉语的人使用的词汇多数都有一定的使用频率。一个外国人将从字典、小说或课本中学会一些词------这些词可能在某一时期是流行的或较可取的，但现在使用的频率大大减少了------例如，用``内人''而不用``妻子''，用``绍介''而不用``介绍''，等等。尽管这样说话一般仍能传达意思，但由于选择不寻常的词而带有异国调子。

假设一个外国人学会了在正常频率下使用所有的词，这会使他讲话真正流利吗？可能不会。在词汇层次之上还有联想层次，它作为一个整体隶属于文化------包括历史、地理、宗教、童话、文学、技术等等。例如，要能非常流利地说现代希伯来语，你需要精通希伯来语的《圣经》，因为这种语言吸收了大量《圣经》的短语和它们的涵义。这样的联想层次深深地渗入每一种语言。然而在流利的同时仍然有各式各样的变化余地------否则，真正流利说话的人就会是那样一种人：他们的思维是最没有生气的！

尽管应该认识到文化对思维有很深的影响，但却不应过分强调语言在引导思维过程中的作用。例如，可以被我们称作``手套''的两个东西会被说英语的人理解成两种不同类型的物体``Zgloves''{[}五指分开的手套{]}和``mittens''{[}联指手套{]}。母语是英语的人比我们更多地意识到这种差别------不过要注意，在中国，南方人会比北方人更多地意识到桔子、柑子以及橙子的差别------北方人可能把他们都叫做``桔子''，而引起这种感知差别的不是母语的不同，而是文化（或亚文化）的不同。

至少就考虑核心而言，母语不同的人其符号之间的相互关系有充分的理由非常相似，因为大家都生活在同一个世界上。当你深入到触发模式那些更为细致的方面时，你会发现共同之处就很少了。这很像这样一种情形：让一些从没有在广西住过的人们设计一些自想国，然后对其中广西的某些农村地区进行比较。不过这没什么关系，只要在主要城市和主要线路上有足够的一致，使整个地图上有共同的参照点就行了。

自想国中的旅行和旅行路线

在整个自想国类比中，对于什么是思维，我用了一个形象化的说法，而一直没有挑明------即，我暗暗地把``思维''与``旅行''相对应。旅行沿途经过的城镇代表被激活的符号。这不是完美的类比，但却是相当有说服力的。这里有个问题：一个想法在人的大脑中重现足够多次以后，它就会逐渐地组块化（凝结）而形成一个单独的概念。这和自想国里非常怪诞的一种事件相对应：经常重复的旅程会以某种奇异的方式变成新的乡镇或城市！如果还要继续使用自想国比喻，记住下面这一点是很重要的：城市不仅代表基本符号（例如，表示``草''、``房子''、``汽车''的那些符号），也代表着由于大脑的组块化能力而产生的符号------诸如那些表示``螃蟹卡农''、``回文斯''或者``自想国''这类深奥概念的符号。

现在，如果认可进行旅行是恰当地对应于思维过程，那么就有下述难题：实际上任何一条路线，从第一个城市到第二个城市，再到第三个，如此下去，都是可行的，只要你记住这些路线都会经过一些中间城市。这对应于符号的任意序列的激活过程：一个接一个，同时允许加入一些额外的符号------那些出现在半途中的符号。也就是说，实质上能按任何期望的次序激活任何一个符号序列。这样一来，大脑似乎就是个不分青红皂白的系统了，它能吸收或产生任何一种思想。但我们都知道并非如此。事实上有一些思想------我们称之为知识或信念------与随意幻想或幽默开心的荒唐念头有着完全不同的作用。我们怎样才能揭示出梦、瞬间即逝的想法、信念以及各种知识之间的异同呢？

可能的、潜在的、反常的通道

有些通道（这里``通道''既可以看作是自想国里的，也可以看作是大脑中的）是从一个地方到另一个地方的常规路线。还有一些通道必须跟着向导才能走。这种通道称为``潜在的通道''，只有出现特殊的外部环境时我们才走它。那些一次又一次被人们采纳的通道组成了知识------这里我指的不仅是关于事实的知识（描述性知识），也指关于怎样做的知识（过程性知识）。这些稳定可靠的通道构成了知识的内容。各项知识逐渐与信念融合，后者同样由可靠的通道代表，不过大概更容易被替换------比如说，一座桥塌了，或者有大雾的时候。剩下的是幻想、谎言、虚假、谬论以及其它花样，它们对应着一些怪异的路线，比如从北京途径张家口、太原、郑州到济南去。这的确是可能的路线，但大概不是一般的路线------不是人们旅行时通常所采用的路线。

这个模型的一个奇怪而有趣的引申就是，从根本上说，上面列出的各式各样``脱离常轨''的思维完全是由信念及各项知识所组成的。也就是说，任何离奇的、曲折的、间接的通道都可分解为一些不离奇、不曲折、直接的路段，而这些短小的、直接了当的符号连线就代表人们可以信赖的简单思维------信念及各项知识。要是认真想一想，这其实并不令人惊奇。道理很简单：我们能够虚构出来的事情，都是以某种方式根植于我们所体验到的现实中的，不论两者有多大的差异。梦也许就是这种在精神里的自想国的随意漫游。局部地看，它有意义------但是整体上看时\ldots\ldots{}

小说翻译的不同风格

读一首类似``炸脖\gebfont{ꪡ}''的诗，就像在自想国中做一次不切实际的旅行，沿着非常古怪的路线迅速地从一个地方跳到另一个地方。那些译文传达了诗的这一侧面，而不是精确地表示那些触发了的符号序列（尽管对此他们做了最大的努力）。在翻译一般的文章时，这样的跨越并不常见。但这类翻译问题的确是有的。假定你在将一本俄文小说翻译成中文，你碰见一个句子，照字面意思翻译出来就是``她喝了一杯伏特加酒''。现在也许你认为你的读者对于什么是伏特加酒没有概念。你可能会试着用他们的文化中``相应''的词来替代------这样，你可以把它译成：``她喝了一杯茉莉花茶''。如果你认为这个例子是个笨拙的夸张，那么你看看陀斯妥耶夫斯基的小说《罪与罚》的第一句话，先看俄文原文，再看几种不同的中文译文。碰巧这部作品有三种不同的中文译本，其中有不少很有意思的事情。

这第一个句子中出现了街道名字``С-й переулок''{[}C胡同{]}和``К-н чост''{[}K桥{]}。这是什么意思？如果一位细心的陀斯妥耶夫斯基作品的读者，他了解列宁格勒（过去叫``圣彼得堡''------或者应该说是``彼得格勒''？），仔细地核对书中其它地理名称（那些地名也只是给出首字母），就会发现那条胡同必定是``СтолярнЫй переулок''（斯陀里亚尔尼胡同）。而那座桥的名字不太容易搞清，我们只能知道是个以``K''开头的人名。陀斯妥耶夫斯基大概是想逼真地讲他的故事，但并不想真实到人们能从文字中得知那些地方的地址------在那里，他使得种种罪行及其它事件发生。在任何情况下，都有个翻译上的问题，或者更确切地说，我们在几个不同层次上都有各种翻译上的问题。

首先，我们是否应该保持首字母，使在这本书第一个句子里就出现了的半神秘气氛得以再现？要是这样，我们就应该用``C胡同''和``K桥''。三位译者中只有一个采取了这个方针。这个译本的读者会感觉到原文的``俄国味''，因为那些地方的名字仍然用的是西里尔字母，从而代表着俄文词。

但若是翻译者想保持原作者隐匿真实地名的意图，同时又觉得熟悉西里尔字母的人太少，又会如何？译者可以选用罗马字母，因为中文读者对罗马字母是相当熟悉的，拼音是罗马字母，数学及各门科学的教科书中也大量地用罗马字母表示变量及未知数。这种方针将导致选用``S胡同''和``K桥''------依然是有种很强烈的``外国味''。第二个译者就是这么做的。这也是我高中时读的那部《罪与罚》英译本所采用的方针。我永远不会忘记，当我开始读这篇小说，遇到了那些仅仅由字母作为名称的地方时，我所体验到的迷惑之感。这本书的开头就使我有了某种莫名的不适。我肯定我漏掉了一些本质的东西，然而我不知道那是什么\ldots\ldots 我断定所有的俄国小说都是离奇的。

第三位译者取了个折中，保持了``S胡同''，但给出了桥的名字：``康桥''。除去这几种办法，我们也可以坦诚地对待读者（可以假定，读者大概对于到底这些地方是真的还是虚构的没有哪怕是一丁点想法！）：利用我们的现代学识，写下``斯陀里亚尔尼胡同''。但假若我们再往前走一步，把那地名的意思翻译过来------而非音译------会怎么样呢？``столяр''的意思是``木匠''，``нЫй''是个形容词词尾，于是这就可以译成``木匠胡同''。这样，我们就可以设想我们是在北京，不是在彼得格勒，处在老舍虚构的情境中，而不是陀斯妥耶夫斯基的。这是我们想要的吗？也许我们应该干脆换一本老舍的作品来读，并为自己辩护说这是一部``中文中相应的作品''。从足够高的层次上来看，它是陀斯妥耶夫斯基小说的一种``译本''------事实上，是可能有的最好的一个！谁需要陀斯妥耶夫斯基？

我们已经从丝毫不差地忠实作者风格的尝试，一直谈到了保持韵味的高层次翻译。现在如果对第一个句子已经做到了这一点，你能设想在书的其余部分会怎样进行下去吗？当一个德国老板娘操着德国味儿的俄语开始大声喊叫时怎么办？你怎么把带着德国口音说出的断断续续的俄语译成汉语呢？

然后你或许还要考虑俗语及口语表达方式的翻译问题。你是该去找一个``类似''的短语，还是该满足于字对字的翻译？如果你是去找一个类似的短语，那么你就要冒着犯``茉莉花茶''类错误的风险。但如果你字对字地翻译每一个习惯用语，那么你这汉语听起来就会有外国味道。也许这是你所期望的，因为俄国文化对于说汉语的人来说是外国文化。但是一个读这种译本的说汉语的人将由于这种不寻常的词序而不断地体会到一种陌生的------人为的------感觉，不是作者想要表达的，也不是俄文原本的读者所体验到的。

或者换一个问题。设想有这么一本英文书，其中有无数英语中的双关语和文字游戏。有人想把这本书译成中文。书中有幽默的对话，涉及了藏头诗、词首字组合等种种形式与内容上的游戏。译者是否只须字对字把原文译出来，添加无数的脚注，用令人感到折磨的说明文形式来解释原文中那些令人轻松愉快、赏心悦目的机巧？无疑这并不忠实于原文的精神。也许译者应该努力在汉语中重建那些双关机制，在这一过程中改换了许多内容上的细节部分，但保持了原文中的核心内容。进一步来讨论的话，假定这本书中有一段涉及对一本俄文小说在英文中的各种译本进行比较。译者是否应该仅是原原本本把这一段译成汉语，然后加许多脚注来解释那些英语词的意思？或者换一种方式------即重写这个段落，对这本小说在汉语中的几种不同译本进行比较------是不是更好？毕竟，这种选择是旨在用一种易懂并有说服力的方式，提出一些有趣的翻译问题以供讨论，而非让读者皱着眉头费劲地去弄懂一件不是用他们的母语描述的事情。但若是译者走得那么远了，有什么必要还是讨论原来那本俄文小说？海明威的小说或者加缪的小说是不是也行呢？也许也行。什么时候你才离原文的精神实质太远了呢？什么是可以变通的，什么是实质性的？

像这样的问题会使人们联想到有关计算机翻译的种种说法，例如，第一批计算机翻译的提倡者之一------华伦·维佛尔在二十世纪四十年代后期曾提出：``{当我看一篇用俄文写的文章时，我就说，`这实际上是用英文写的，只是用一些奇怪符号进行了编码。我现在要进行解码工作。'}''\footnote{}维佛尔的话不能仅从字面上理解。相反，它必须被看成是以一种富于启发的方式在说明：有一种可以客观地描述的意义------或至少是某种相当接近客观的东西------隐藏在符号之中。所以，没有理由假定即使有编制得足够好的程序，计算机也无法捕捉到意义。

程序之间的高层次比较

维佛尔说的是不同的自然语言之间的翻译。现在我们来考虑两种计算机语言之间的翻译问题。比如说，假定有两人分别编写了运行在不同计算机上的两个程序，而我们想知道这两个程序是否执行同样的任务。我们怎样才能知道呢？我必须比较这两个程序。但应该在什么层次上比较？也许一个程序员用机器语言写，另一个程序员用编译语言写。这样的两个程序是可比较的吗？当然可以。可是怎么比呢？方法之一就是编译那个编译语言程序，产生一个该计算机上的机器语言程序。

现在我们有两个机器语言程序。但还有一个问题：这里是两台计算机，因此是两种不同的机器语言------而且可能是极不相同的。一台机器可能是十六位字长，而另一台则可能是三十六位字长。一台机器可能有内部堆栈处理指令（推入与弹出），而另一台没有。两台机器硬件上的差别也许会使这两种机器语言看上去不可比较------而我们依然觉得它们是在执行同样的任务，而且希望一眼就看出这一点。显然，我们现在看这两个程序时离得太近了。

我们需要做的是后退一步，离开机器语言，采用较高层次的、组块化的观点。我们希望从这一有利角度能看出程序中的组块，这些组块使程序看上去是经过了合理的全局------而非局部------设计。也就是说，程序是由这些组块很好的配合在一起而组织成的，使我们能领会程序员的目的。让我们假定这两个程序最初都是用高级语言写的，那么某种组块就已经有了。可是我们会陷入其它麻烦。高级语言的数目在迅速增多：Fortran、Algol、Lisp、APL以及众多其它的语言。你怎么比较用APL写的程序与用Algol写的程序？当然不能是一行一行地比。你将再次在心里对这些程序进行组块，寻找能彼比对应的概念功能单位。所以，你不是在比较硬件，也不是在比较软件------你是在比较``以太件''------那些在软件背后的纯粹概念。在你能实行任何一种有意义的比较之前，你必须从低层次中拔出某种抽象的``概念骨架''，不论你是在比较用两种不同的计算机语言写成的程序，还是比较两种动物，或者是比较两种不同的自然语言写成的句子。

这就把我们带回到一个以前提出来的问题上，那是关于计算机与大脑的：我们怎样才能理解计算机或大脑的一个低层次描述？在这么复杂的系统里，有没有一种从随便什么合理意义上说都是客观的方法，可以从低层描述中抽出高层描述？在计算机的情形中，完全展示存储器的内容------这称为存储器转储------是容易做到的。用计算机进行计算的早期年代里，程序出错时通常都打印出全部存储器内容。然后程序员就不得不回家去，花几个小时钻研存储器的内容，努力去理解存储器的每个小片断都表示了什么。实质上，这个程序员要做的工作与一个编译程序所做的工作正好相反：他要把机器语言翻译成高层次语言，一种概念级别上的语言。结果，这个程序员理解了这个程序的目标，并能用高层次的语汇来描述------例如，``这个程序把英文小说翻译成中文''，或者``这个程序根据输入的任何主题写出八声部赋格曲''。

大脑之间的高层次比较

现在我们该从大脑的角度来研究上述问题了。这时问题是这样提出的：``人的大脑也能在高层次上被`读出'吗？有没有对大脑内容的客观描述？''在《蚂蚁赋格》中，食蚁兽宣称，他能通过察看奔忙的蚁群------蚂蚁的组成部分------而得知蚂蚁在想什么。会不会有什么超人------也许，是个``食神经元兽''------存在，可以想象他俯视我们的神经元，把所看到的进行组块，得出对我们思维的分析？

当然，回答必须是肯定的，因为我们在任何时候都能用关于组块的词语（即：非神经元的）去描述我们的意识活动。就是说我们有一种机制，使得我们能够对我们的大脑状态进行粗略的组块，并对之进行功能描述。更确切地说，我们并不对所有的大脑状态都进行组块------我们只是对活跃部分进行组块。然而，如果别人就某一问题向我们提问，而那个问题编码在我们大脑中当前不活跃的区域里，那么我们几乎能在瞬间里进入那个休眠区域，并对那个问题给出组块化的描述------就是说，给出关于那个问题的信念。注意，我们的描述中没有任何关于大脑中那一部分的神经元层次上的信息：我们的描述是高度组块化的，以至于我们甚至一点都不清楚我们描述的是大脑中的哪一部分。这恰好和那个程序员的例子成为鲜明对照：他的组块化描述来自对存储器全部内容的每个部分进行有意识的分析。

如果一个人能对自己大脑中的任何部分做出组块化描述，那么一个局外人，掌握了进入那个大脑的一些非破坏性手段，为什么不能做同样的事情呢？那局外人应该不仅能对那个大脑的有限部分进行组块，而且实际上还能给出那个大脑的完整的组块化描述------换句话说就是，完整地列出一个其大脑是可达的人的全部信念。显然这样的描述会有天文数字的规模，不过在这里没关系。我们感兴趣的是，原则上是否存在关于大脑的一个界说良好的、高层次的描述？或者反过来，是否神经元层次的描述------或者某种同样是生理上的、但在直观上还不清楚的什么描述------是原则上存在的描述中最好的一个？毫无疑问，如果我们要探究我们是否能了解自己，那对这个问题的回答将是至关重要的。

潜在信念，潜在符号

有可能进行组块化描述，这是我的论点。但有了这个描述后，并不是一切就突然清楚明白了。问题是，为了从大脑状态中抽出组块化描述，我们需要一种语言来描述我们发现的东西。看起来，描述大脑的最恰当方式是枚举它能具有的各个种类的思想，以及它不可能具有的种种思想------或者也许是，枚举它相信的事情及它不相信的事情。如果这是我们进行组块化描述时的奋斗目标，那就很容易看出我们将碰上什么样的麻烦了。

假设你想枚举出在一个自想国里能进行的所有旅行。这会有无穷多种可能。可是你怎么确定哪种可能是似乎合理的？而且，``似乎合理''是什么意思？在试图找出什么是大脑中符号到符号的``可能通道''时，我们要碰到的恰恰就是这种困难。我们可以设想四脚朝天的一条狗从空中飞过，嘴里还叼着一根雪茄------或者公路上两个巨大的油煎鸡蛋撞上了------或任何其他的荒唐情景，要多少有多少。我们大脑中可以采纳的不自然的通道，其数目是没有限制的，就像在自想国里能计划无数条不切实际的旅行路线一样。但对于一个给定的自想国，到底是什么构成了一条``合情合理''的旅行路线？而对于一个给定的大脑状态，到底是什么构成了一个``合理''的想法？大脑状态自身是不禁止任何路径的，因为对于任何一条通道，都会有一些环境使之被采纳。大脑的物理状况所提供的信息------如果是正确读出的------不是说明哪条通道可被采纳，而是沿这条通道走时会遇到多少阻力。

在一个自想国里，可以沿着两条或更多的可选择合理路线进行旅行。比如说，从北京到广州的旅行既可以沿东面路线走，也可以沿西面路线走。两者都是合理的。但在不同的情况下人们有各自的倾向。你在这个时候看了地图，并不能得知在遥远的将来的某个时候哪条路线更为可取------那取决于进行旅行时的外界环境。同样，``读出''一个大脑状态将会揭示出，对于给定的一个符号集合，往往有不止一条合理的可选通道能用来连接符号。不过，这些符号之间的旅行未必是要马上发生的，可以只是所读出的大脑状态中的亿万个``潜在的''旅行当中的某一个。由此得出一个重要结论：大脑状态本身不具有说明哪条路线将被采纳的信息。外界环境在决定路线的选择时扮演着极为重要的角色。

这意味着什么？这意味着在不同的环境下，同一个大脑中可以产生出彼此冲突的想法。任何有价值的高层次大脑状态读出形式都必须包括对所有这类彼此抵触的想法的刻画。实际上这是显而易见的------我们全都是一大堆矛盾的混合体，为了把他们统一起来，我们总是在一个特定时刻只表现出其中的一个方面。选择什么是无法预料的，因为事先并不知道什么条件会对选择施加影响。大脑状态所能提供的是------如果恰当地读它------对路线选择的一个有条件的描述。

我们来考察一个例子，看看《前奏曲》中描写的螃蟹的境况。他对奏出乐曲有各种各样的反应。有时候他几乎无动于衷，因为他太了解那支曲子了。另一些时候，他会相当激动，但这样的反应需要外界恰当的触发------比如说，有一位热心的听众在场，对于这位听众来说作品是新的。高层次地读出螃蟹的大脑状态大概就会揭示出那些潜在的激动（以及会导致激动的条件），也会揭示出潜在的麻木（以及会导致麻木的条件）。然而，大脑状态本身不会说出下一次听这支曲子时会出现什么状况。它只能说：``如果有了这样那样的条件，那么就会导致激动；否则\ldots\ldots''。

所以，一个对大脑状态的组块化描述将给出一个目录，列出依环境而有条件地引发的信念。由于不可能枚举所有可能的环境，我们只好满足于列出那些我们认为是``合理''的信念。更进一步，我们还不得不满足于对环境本身的组块化描述，因为显然不可能------也不应该------降到原子层次上去刻划环境。因此，我们不可能做出精确的、决定论的预言，指出在给定的组块化环境下，从大脑状态中会抽出什么信念。概括地说就是：一个大脑状态的组块化描述将由一个带或然性的登记表构成，其中列着一些信念（或符号），它们在各式各样``多半合理的''环境中------这些环境本身也是在组块层次上描述的------最可能被唤起（或触发）。脱离背景环境而试图对一个人的信念进行组块是愚蠢的，这恰恰就像不涉及配偶而试图描述一个人``潜在后代''的分布一样。

同样类型的问题也出现在枚举一个给定人脑中所有符号的时候。一个大脑中不仅潜在地有无穷多的通道，还有无穷多的符号。就像前面已经指出来的，新概念总能从已有概念中形成，而且可以争辩说表示这种新概念的符号不过就是潜存在每个个体中等待着被唤醒的休眠符号。它们也许在这个人的一生中从未被唤醒，但却可以宣称那些符号一直是存在于那里的，只等着恰当的环境来触发它们的合成。不过，如果这种可能性非常小，似乎把``休眠''用于这种情况就不太符合实际了。为了使这一点更为清楚，可以试图去想象当你醒着的时候蛰居在你头颅里的所有``休眠的梦''。是否有可能存在一个判定过程，在给定你的大脑状态时能区分开``不可梦想的主题''与``潜在地可梦想的主题''？

自我意识在哪里？

回顾我们讨论过的问题，你也许会自言自语，``这些关于大脑和心智的思考都挺好，但是涉及到意识的那些知觉怎么样呢？这些符号尽可以随心所欲地相互触发，但除非某个人觉察整个这件事，否则就没有意识''。

直观上这在某个层次上有道理，但在逻辑上却不太说得通。因为这样一来就会迫使我们去寻找一种解释这种机制（觉察所有活跃符号的机制）的说法，如果这种说法还没有包括在我们至今所给出过的描述里。当然，一位``唯灵论者''不必看得这么远------他只需断言所有这些神经活动的觉察者就是灵魂，而灵魂是不能用物理词汇描述的，那就够了。不过，我们将试图给出一个``非唯灵论者''的解释，说明意识出现在哪里。

我们这个有别于``唯灵论者''的解释------这也是个令人困惑的解释------是停在符号层次上说：``这就是它------这就是所谓的意识。意识是系统的一种性质，每当系统中有服从触发模式------多少有点像前几节描述过的触发模式------的符号时，这种性质就会出现''。这么简单，似乎太不够了。这怎么能说明关于``我''的知觉，即自我意识呢？

子系统

没有理由期望``我''，或者``自我''，不该由一个符号所表示。事实上，表示自我的符号可能是大脑的所有符号中最复杂的一个。出于这个原因，我愿意把它放在分层结构中的一个新层次上，并称之为子系统，而不称符号。确切地说，所谓``子系统''，我指的是符号的集群，其中的每个符号都能在子系统本身的控制下被激活。我想传达出来的子系统图像是这样的：它起的作用几乎像个独立的``子脑''，具有自己的可选符号集，而那些符号则可以在它内部相互触发。当然，子系统和``外界''------也就是大脑的其余部分------也有大量通讯。``子系统''不过是一个生长过度的符号的别称，这种符号已变得非常复杂，其中有许多子符号彼此相互作用。所以，在符号与子系统之间没有严格的层次区别。

由于子系统与大脑的其余部分有广泛的联系（不久就会描述几种这样的联系），所以很难在子系统和外界之间划出明确的界线。但即使边界模糊，子系统仍然是相当实在的东西。有趣的是，子系统一旦被激活，然后撒手不管，它就会自己一直工作下去。这样，一个个体的大脑中两个或更多的子系统可能会同时运转。我就注意到过偶尔发生在我自己大脑中的这类事：有时我发觉我的脑海里响起两支不同的曲子，争夺着``我的''注意力。每一支曲子都是在我大脑中的一个单独隔间里以某种方式制造或``演奏''出来的。两个负责从我的大脑中提取出曲子的子系统大概都各自激活了一系列符号，一个接着一个，完全不顾另一个子系统也在做同样的事情。然后两者都试图同我大脑的第三个子系统------我的自我符号------进行通讯，就在这个时候我大脑中的``我''察觉到发生了什么。换句话说，``我''开始获得了一个关于``那两个子系统的活动"的组块化描述。

子系统和共用编码

典型的子系统大概是那些代表我们亲密朋友的子系统。那些人在我们的大脑中被表示得非常复杂，以至于代表他们的那些符号扩张到了子系统的级别，变得能自主地活动并且还可以利用我们大脑中的一些资源来维持自己。在这里我的意思是，一个将朋友符号化的子系统能像我自己一样在我的大脑里激活很多符号。比如，我可以启动我的一个代表某好朋友的符号，然后就几乎真的感觉到自己处于他的地位上，冒出他会有的种种想法，激活那些更多地是反映他的思维模式的符号序列。可以说我这个朋友的模型------在我大脑中体现为一个子系统------就构成了我对他大脑的一个组块化描述。

那么，对于那些我认为在他大脑中会有的符号，这个子系统是否包括了一些与之分别对应的符号呢？这其实是多余的，很可能这个子系统会大量利用我大脑中已有的符号。例如，当我大脑中表示``山''的那个符号被激活时，它就可以被这个子系统借去。这时，子系统使用那个符号的方式不必等同于我整个大脑使用它的方式。具体说，如果我在和我的朋友谈论亚洲中部的天山山脉（我们俩都没去过），并且我知道几年前他在阿尔卑斯山有过一段极棒的远足经历，那么我心里对他的那些谈论的解释就会部分地带有我印象中他从前那次阿尔卑斯山经历的色彩，因为我会去努力设想他会如何想象那个地方。

使用我们在本章建立起来的词汇，我们可以说，符号``山''在我这里的激活是受控于我的那个代表他的子系统。其效果是在我的记忆中打开一个通常情况下我不用的窗口------就是说，我的``缺席选择开关''从我通常的记忆范围扳向了我对他的记忆的记忆。当然，我对于他的记忆的表示只是他的实际记忆的一个近似，构成他记忆的那些大脑符号彼此激活的形态十分复杂，我是无法全部了解的。

他的记忆在我这里的表示也是处于复杂的大脑符号彼此激活的形态，当然这里是我的大脑中的符号------即那些``原始的''概念，例如草、树、雪、天、云等等。我必须假定在他那里对这些概念的表示与在我这里是``等同的''。我还必须假定对更基本的观念我们俩也有类似的表示，如重量、呼吸、疲劳、颜色等感觉。登上峰顶瞭望四周景色时的愉悦感虽不那么基本，却也几乎是每个人都会有的。

因此，我的大脑中产生这种愉悦感的复杂机制可以直接搬到``朋友子系统''中去，而且不会太走样。

我们可以接着做下去，尝试去描绘我怎样完整不漏地理解朋友讲给我的一个故事，一个充满了复杂的人际关系和心灵体验的故事。可是我们的术语很快就会变得不够用了。当涉及对这个或那个事情的他所表示的我所表示的他的表示时，就会现出复杂的递归。假如故事中还有我们共同的朋友，那么我会在我自己关于他们的形象和我关于他那里对他们形象的表示之间寻找调和。纯粹递归在处理这种类型的符号混合物时显然是种不适用的形式手段。况且我们只不过才触及到问题的表面！我们目前实在缺乏词汇去描述符号之间可能有的复杂的相互作用。所以在我们陷入窘境之前还是止步为好。

不过我们应该注意，计算机系统在某些方面已经开始达到同样类型的复杂性，而这些概念中的一部分也因此有了名称。例如，我的``山''符号就类似于计算机术语中的共用（可重入）编码------那些可同时被运行在一台计算机上的两个或更多的分时程序使用的编码。当一个符号同时是几个不同子系统的部分时，激活它就会有不同的结果。这一事实可以作这样的解释：那个符号的编码在被不同的解释程序处理。所以，``山''符号中的触发模式不是绝对的，而是和在哪个系统中激活符号有关。

有些人可能会怀疑这种``子脑''的真实性，下面是艾舍尔在谈论他怎样创作他那些充溢在平面里的循环画时说的一段话，也许对弄清我这里指的是什么现象会有所帮助：

当我作画时，有时会觉得似乎自己是个巫师，被我召唤出来的灵物所控制。好像是它们在决定选择什么形状来显现自己。它们在诞生过程中，一点也不考虑我的评判，它们将怎样发展我也施加不了多少影响。它们通常都是些很难对付并且很固执的家伙。\footnote{}

这是个理想的例子，说明大脑中某类子系统一旦被激活，就几乎是自主的了。对于艾舍尔来说，他的子系统几乎能统治他的审美判断力。当然，不能把这一点看得太绝对，因为那些强有力的子系统是作为一个结果，来自他多年的训练并严格服从那种塑造了他的审美意识的力量。总之，如果把艾舍尔大脑中的那些子系统与艾舍尔本人或他的审美判断力分隔开，那就错了。在他的美感中，``他''是这个艺术家的完整自我，而那些子系统是必不可少的组成部分。

自我符号与意识

自我子系统活动时有一个非常重要的旁效，即它在下述意义上能扮演``灵魂''的角色：在不断地与大脑中其他子系统及符号进行通讯时，它密切注意着哪个符号处于激活状态，以及是以什么方式在活动。这就是说它必定有表示心智活动的符号------换一种说法，就是那些表示符号的符号，以及表示符号活动的符号。

当然，这并没有把意识或感知提高到什么``有魔力的''非物理的层次上。这里感知是我们所描述过的那些复杂的硬件软件的一个直接后效。这种方式描述感知------即把感知作为大脑本身的一个子系统对大脑活动的监控------尽管源于世俗，仍然和我们都知道并称为``意识''的那种近乎不可描述的感觉相像。读者肯定能看出这里的复杂程度已经足以产生许多意外的效果。例如，很可能一台具有这种结构的计算机会谈论自己，那些话极像人们通常谈论自己时说的，包括固执地强调它有自由意志，它不能被解释成``部分的总和''等等。（有关这方面的问题，可参阅明斯基的文章《物质、心智和模型》）。

我一直假定有一个表示自我的子系统。怎么能保证这样的子系统实际存在于我们的大脑里呢？如果没有自我符号的发展演化，一整套复杂的符号网络------比如我们上面描述过的那种------也能够发展演化吗？如果没有一个表示宿主有机体的符号，这些符号及其活动怎么能实现那些``同构''于周围世界中的实际事件的心理事件呢？进入系统的所有刺激在片刻中都集中在一个很小的空间范围之中。要是没有一个表示有机体的符号，处于这个有机体中的大脑符号结构就会有一个很显眼的洞，需知，在这个大脑符号结构所反映的事件中，那个有机体扮演着一个比其它对象都更重要的角色。事实上，仔细想一想的话，理解那个围绕着一个生物体的世界的唯一途径，似乎就是在与周围物体的关系中理解这个生物体所起的作用。这就必须要有一个自我符号存在。而且，从符号到子系统这一步骤也不过是反映了自我符号的重要性，并非质的变化。

我们与卢卡斯的初次会面

牛津哲学家卢卡斯（和我们前面谈到的``卢卡斯数''没有关系）于1961年写过一篇引人注目的文章，题目是``心智、机器和哥德尔''。他的观点和我的观点几乎完全对立，然而他用来论证他的观点的材料中有许多和我的材料相同。下面这段节录与我们刚才的讨论有很大关系：

当一个人第一次进行最简单的哲学探讨时，他就会陷入下面一系列问题：如果某人知道了一件事，他是否知道他知道这件事？一个人自我反思的时候他在想什么，当时是什么机制在进行思维？经过长时间的困惑与挫折之后，他学会了不去纠缠这些问题：他暗暗体会到有意识物与无意识物的概念是不同的。当我们说一个有意识物知道什么事情的时候，我们说的不仅是他知道那件事，还包括他知道他知道那件事，以及他知道他知道他知道那件事，如此下去，直到我们不再纠缠这个问题------在这里我们识别出了无穷，但并不是什么不好意义下的无穷回归，因为逐渐消失掉的并不是答案，而是那个没多少意义的问题。感到这些问题没有意义，是因为在其概念自身中包含着对这些问题的回答能无限进行下去的思想。

虽然有意识物有能力进行下去，我们不希望仅仅作为一系列他们能从事的工作而展示这一点，也不把心智看作是一个自我和超自我和超超自我的无穷序列。相反，我们坚持这样的观点：有意识物是个整体，虽然我们也谈论心智的各个部分，但那只是一种比喻的说法，不能按字面意义理解。

意识悖论的出现是因为有意识物除了能意识到别的东西，还能意识到自身，并且的确不能看作是可分解的。这意味着有意识物可以用某种手段处理哥德尔式的问题，而机器却不能用同样手段。这是因为有意识物可以既考虑自己也考虑其工作，同时依然是在那里进行工作，而不是在干别的。一台机器可以被制造成那种所谓能``考虑''它在干什么的机器，但它无法在工作时把它的所作所为``考虑进去''，除非它变成一台新的机器，也就是说，这台机器加上了一个``新的部分''。但在我们关于有意识的心智的观念中，能反思自己和评判自己的所做所为是它固有的特点，不需要外加的部分：它已经是完备的了，没有固有的致命弱点。

这样，我们的论题开始变成更多地是属于概念分析问题，而非数学发现。这一点由图灵提出的另一个论据所证实。至今，我们只造出了相当简单的并且是可预知其行为的人工制品。随着我们增加机器的复杂度，也许会有令人惊讶的事情在等着我们。他用核反应堆的裂变作比较。在某个``临界点''之下，几乎没什么事情发生，但是一旦超过临界点，就会火花四溅。大脑和机器或也许是这样。多数大脑和所有的机器目前都是在``临界点''之下的------它们对受到的剌激冷漠而迟钝，而且毫无主见，只能作出平庸的反应------但现在有少数大脑，以及将来可能有的一些机器，是超临界点的，能独自焕发出光彩。图灵的意思是，这只是个复杂程度问题，超过某种复杂水平就会出现质的变化。因此``超临界''的机器会与人们迄今设想的简单机器非常不同。

也许是这样。复杂程度的确经常引起质的不同。尽管听起来难以置信，结果可能就是：在某种复杂度级别之上，一台机器甚至从本质上都不再是可预知的，它开始自行其是，或者换一种富于启发性的说法，就是，当它不再是完全可预知和完全可驾驭时，它也许就开始有它自己的心智了，它能做出一些我们承认是需要智能的事情，而非仅仅是出错或随便瞎猜，同时那还不是我们事先编好程序的。但若是这样，就那些行为所处的背景而言，它已不再是台机器了。在机械论者的争辩中，关系重大的不是心智怎样形成，或心智可能会怎样形成，而是心智怎样活动。对于机械论者的论点来说，关键是心智的机械模型将按照``机械原则''运转，也就是说，我们可以从各个部分的运转去理解整体的运转，而每个部分的运转或者取决于它的初始状态和机器的结构，或者是在一定数目的确定的运转之间作一个随机选择。如果机械论者造出了一台机器，复杂得这一切都不能很好地保持下去，那么对于我们现在的讨论而言，它就不再是一台机器了，无论它是怎样建造的。相反，我们应该说是创造了一个心智，和现在人类生育是同一种意义。于是产生新的心智的方式就会有两种：传统的方式------即妇女生小孩------和新的方式，即建造非常非常复杂的系统：比如说电子管和继电器组成的系统。当谈到第二种方式时，我们应该注意强调，虽然创造出来的东西看上去像台机器，但却不是真正的机器，因为它不单是部分的总和。一个人仅仅知道它的构造及其各部分的初始状态，是无法说出它将会做些什么的，甚至无法说出它能力的极限在哪里，因为即使给它一个哥德尔型问题，它也会正确回答。事实上我们应该简单地说，没有被哥德尔问题难倒的系统就不是图灵机，即：在``对哥德尔问题的反应''这一点上，它们不是机器。\footnote{}

在读这段文章时，我不断地被那一连串的话题、引喻、暗指、混淆以及结论弄得不知所措。我们从卡罗尔式的悖论跳到哥德尔再跳到图灵再到人工智能再到整体论和简化论，全都在短短的两页纸上。卢卡斯这套玩意儿唯一的作用就是刺激我们的思想。以下各章里，我们还将回到在这篇奇怪的文章中许多撩人地一掠而过的话题。

\begin{enumerate}
\item
  \phantomsection\label{Chapter12.xhtml_B_1}
\item
  \phantomsection\label{Chapter12.xhtml_B_2}
\item
  \phantomsection\label{Chapter12.xhtml_B_3}
\end{enumerate}

\phantomsection\label{Dialog13.xhtml}{}


\end{document}
